

\medterm{S Phase} The stage of the cell cycle when a cell copies all of its DNA before dividing. This duplication allows each new cell to receive a complete set of genetic instructions.

\medterm{S-Adenosylmethionine} A substance made naturally in the body that helps chemical reactions involving cell membranes, genetic material, and brain chemicals. It is also sold as a supplement, but its benefits and interactions vary by condition and require medical advice.

\medterm{S3 Gallop Pathology} An extra heart sound heard just after the normal second heart sound as the lower chambers fill rapidly. It can be normal in children and pregnancy but may suggest excess fluid or heart failure in an older adult.

\medterm{S4 Gallop Pathology} An extra heart sound heard just before the normal first heart sound when an upper chamber pushes blood into a stiff lower chamber. It may occur with long-standing high blood pressure, a narrowed aortic valve, or thickened heart muscle.

\medterm{SA Node} A small group of specialized cells in the upper right heart chamber that normally starts each heartbeat. It produces regular electrical signals that spread through the heart and coordinate its rhythm.

\medterm{SAA} Usually an abbreviation for serum amyloid A, a blood protein that rises during inflammation. Persistently high levels can contribute to abnormal protein deposits in organs during some long-lasting inflammatory diseases.


\medterm{Sabin Vaccine} The oral polio vaccine containing live attenuated poliovirus strains. It produces strong intestinal immunity but, rarely, vaccine-derived poliovirus can occur, so vaccine policy depends on local public-health guidance.

\medterm{Saccade} A quick, coordinated eye movement that shifts the eyes from one target to another. Abnormal saccades can occur with disorders affecting the eye muscles, balance system, or brain and nerves.


\medterm{Saccular Aneurysm} A rounded pouch bulging from one side of an artery because that part of the vessel wall is weak. It may enlarge, form clots, or burst; risk depends on its location, size, shape, and growth.

\medterm{Sachet Poisoning} Illness caused by swallowing, inhaling, or getting the contents of a sachet on the skin or in the eyes. Effects depend on the substance and amount; some sachets contain medicines or dangerous chemicals.

\medterm{Sacral} Sacral means relating to the sacrum, the triangular bone at the base of the spine that connects with the pelvis.

\medterm{Sacral Dimple} A small indentation in the skin over the lower spine that is present at birth. Most are harmless, but a deep, large, high, hairy, discolored, or draining dimple may need imaging for an underlying spinal abnormality.
\synonyms
Pilonidal Dimple

\medterm{Sacral Fistula} An abnormal channel near the lower spine that connects the skin with deeper tissue or another body space. It may cause persistent drainage, pain, repeated infection, or a skin opening and can result from birth differences, injury, or pilonidal disease.

\medterm{Sacral Nerve Stimulation} Treatment that sends mild electrical pulses to nerves near the lower spine to improve bladder or bowel control. It may help selected people with persistent urinary or fecal leakage, retention, or constipation after other treatments have not worked.

\medterm{Sacral Nerve Stimulator} An implanted device that sends mild electrical signals to nerves near the lower spine to improve bladder or bowel control in selected people. It may help incontinence or some emptying problems when other treatments have not worked.

\medterm{Sacral Nerves} The spinal nerves arising from the sacral spinal cord and supplying the pelvis, pelvic floor, bladder, bowel, and lower limbs.
\medterm{Sacral Region} The lower back and posterior pelvic region overlying the sacrum, a fused triangular bone that connects the lumbar spine with the pelvis.

\medterm{Sacral Vertebra} One of the fused vertebrae forming the sacrum, which supports the spine, joins the pelvic bones, and transmits trunk weight to the lower limbs.

\medterm{Sacral Vertebrae} Five spinal bones that fuse to form the sacrum at the back of the pelvis. They transfer weight from the spine to the hips and contain openings through which sacral nerves pass.

\medterm{Sacral Vertebrae (S1-S5)} The five sacral vertebrae that fuse to form the sacrum, connect the spine to the pelvis, and provide foramina for sacral nerve roots.
\medterm{Sacrococcygeal} Sacrococcygeal means relating to the sacrum and coccyx at the lower end of the vertebral column.

\medterm{Sacrococcygeal Teratoma} A tumor present from birth at the base of the spine and made of several tissue types. It may grow outside or inside the pelvis; a large fetal tumor can strain the heart and increase amniotic fluid.

\medterm{Sacrocolpopexy} An operation that supports a prolapsed top of the vagina or uterus by attaching it to the strong tissue over the sacrum, usually with surgical mesh. It may be performed through open or small incisions; bleeding, infection, mesh problems, bowel or bladder injury, and recurrent prolapse are possible.

\medterm{Sacroiliac Joint} The joint connecting the triangular sacrum at the base of the spine to each hip bone. It transfers weight between the spine and legs and may become painful or inflamed after injury, pregnancy, arthritis, or infection.

\medterm{Sacroiliac Joint Dysfunction (Sacroiliac Joint Pain)} Alternate index label for Sacroiliac Joint Pain.

\medterm{Sacroiliac Joint Dysfunction (SI Joint Pain)} Pain or abnormal movement from the joint between the lower spine and pelvis. It may cause pain in the lower back, buttock, or thigh and can follow injury, pregnancy, arthritis, or changes in walking.

\medterm{Sacroiliac Joint Pain} Sacroiliac-joint pain arises near the junction of the lower spine and pelvis and may follow injury, pregnancy, arthritis, leg-length or gait changes, or inflammation. Pain can radiate to the buttock or thigh; examination distinguishes it from hip or spine disease.

\synonyms
Sacroiliac Joint Dysfunction (Sacroiliac Joint Pain)


\medterm{Sacroiliitis} Inflammation of one or both joints connecting the lower spine to the pelvis. It can cause pain in the lower back, buttock, or back of the thigh and may result from inflammatory arthritis, infection, injury, or wear.

\medterm{Sacrospinous Ligament Fixation} An operation through the vagina that attaches the top of the vagina or cervix to a strong pelvic ligament to treat prolapse. It avoids an abdominal incision but may cause buttock pain, bleeding, infection, nerve injury, or recurrent prolapse.

\synonyms
Sacrospinous Colpopexy; Sacrospinous Suspension

\medterm{Sacrum} A large triangular bone at the back of the pelvis formed by five fused spinal bones. It connects the spine to the hip bones, transfers body weight to the legs, and contains openings for nerves supplying the pelvis and lower limbs.

\medterm{SAD} An abbreviation for seasonal affective disorder, a recurring pattern of depression or mood change linked to a particular season, often winter and reduced daylight. Treatment may include therapy, medicines, or supervised light therapy.
\synonyms
Seasonal Affective Disorder (SAD)

\medterm{Sad Mood} Sad mood is a feeling of unhappiness or low mood that may be a normal response or a symptom of depression, grief, illness, or another condition.

\medterm{Saddle Pulmonary Embolism} A large blood clot lodged where the main artery from the heart divides toward both lungs. It can severely block blood flow, causing sudden breathlessness, chest pain, fainting, low blood pressure, or cardiac arrest and requires emergency treatment.

\synonyms
Saddle Embolus; Bifurcation Pulmonary Embolism

\medterm{Sadism} Sexual arousal or gratification from inflicting pain, humiliation, or suffering; clinical significance depends on distress, impairment, or nonconsent.


\medterm{Safe Sex} Sexual practices that reduce the risk of sexually transmitted infections and unintended pregnancy, including condoms or other barriers, testing, vaccination, and informed consent.


\medterm{Safety Of Drugs} The process of balancing a medicine’s expected benefits against side effects, interactions, and reasons it may be unsafe. Safe use includes checking the person, dose, allergies, other medicines, monitoring needs, and instructions.
\medterm{Safety Plan} A collaboratively written, practical plan for responding to escalating suicidal thoughts
or emotional crisis. It identifies warning signs, internal coping strategies, supportive
contacts, professional resources, and methods to restrict access to lethal means.

\medterm{Safety Study} A safety study evaluates whether a medicine, device, procedure, or treatment causes harmful effects and identifies risks before or during clinical use.



\medterm{Safrole} Safrole is a natural aromatic compound found in some plants and oils; high or prolonged exposure is toxic and has carcinogenic concern.

\medterm{Sagittal} Oriented in a vertical direction that divides the body or a structure into right and left parts. A midsagittal plane divides the parts equally; a parasagittal plane does not.

\medterm{Sagittal Plane} An imaginary vertical plane dividing the body or an organ into right and left portions. A midsagittal plane divides them equally; a parasagittal plane divides them unequally.

\medterm{Salaam Attacks} Brief repeated bending spasms or sudden head drops in infants. They may be seizures and require urgent specialist assessment because early treatment can protect development and reduce the risk of continuing epilepsy.


\medterm{Salbutamol (Albuterol)} A quick-relief inhaled medicine that relaxes tightened airway muscles in asthma and chronic obstructive lung disease. It eases wheezing and breathlessness but may cause shaking, a fast heartbeat, or palpitations and does not replace treatment that controls airway inflammation.

\synonyms
Albuterol

\medterm{Salicylate} A group of chemicals that includes aspirin and can reduce pain, fever, or inflammation. Excessive amounts may cause poisoning with vomiting, ringing ears, fast breathing, confusion, dehydration, or dangerous blood chemistry changes.

\medterm{Salicylate Toxicity} Poisoning caused by too much aspirin or another salicylate. It may cause vomiting, ringing in the ears, fever, sweating, fast or deep breathing, dehydration, confusion, seizures, or coma and requires urgent medical advice.

\medterm{Salicylic Acid} A medicine applied to the skin to soften and remove excess surface skin. It is used for acne, psoriasis, warts, corns, and similar thickened or blocked skin; irritation is common and large exposures can be toxic.

\medterm{Saline} A sterile solution of sodium chloride in water. Different concentrations are used for intravenous fluids, wound or eye irrigation, inhalation, laboratory work, or dilution of medicines.
\synonyms
Normal saline

\medterm{Saline Infusion Sonohysterography} An ultrasound test in which sterile salt water is placed inside the uterus to outline its cavity. It can show polyps, fibroids, scar tissue, or an unusually shaped cavity and may cause temporary cramping or light bleeding.

\medterm{Saline Nasal Washes} Rinsing the nasal passages with sterile, distilled, or appropriately treated saline to loosen mucus and reduce nasal irritants; unsafe water can cause serious infection.

\medterm{Saliva} Fluid made by the salivary glands that moistens food, begins digestion, protects the mouth and teeth, and helps swallowing and taste. Too little saliva can cause dry mouth, mouth soreness, and tooth decay.

\medterm{Salivary Duct Stones} Calcified deposits obstructing a salivary duct, most often the submandibular duct, causing recurrent painful swelling of a gland during meals and sometimes infection.

\medterm{Salivary Gland} A gland that makes saliva and releases it into the mouth. The three major paired glands are the parotid, submandibular, and sublingual glands; many smaller glands lie in the mouth lining.

\medterm{Salivary Gland Biopsy} Removal of salivary-gland tissue or a needle sample for microscopic examination of suspected autoimmune disease, inflammation, infection, or neoplasm.

\medterm{Salivary Gland Cancer} Cancer arising in a salivary gland, most often the parotid near the ear. It may cause a persistent lump, pain, numbness, dry mouth, or facial weakness; imaging and biopsy identify the type and guide treatment.

\medterm{Salivary Gland Infection} Infection of a salivary gland, usually bacterial and often involving the parotid or
submandibular gland. Painful swelling, fever, dry mouth, and purulent drainage from the
duct may occur.

\medterm{Salivary Gland Infections} Infections of a salivary gland or its duct, commonly bacterial or viral, causing painful swelling, tenderness, fever, dry mouth, or purulent drainage.

\medterm{Salivary Gland Neoplasm} An abnormal growth in a salivary gland that may be noncancerous or cancerous. It can cause a lump, pain, facial weakness, dry mouth, or swallowing problems and may require imaging and biopsy.

\medterm{Salivary Gland Scintigraphy} An imaging test that measures salivary-gland function after a small radioactive tracer is given. It may help assess dryness from Sjögren syndrome, blocked ducts, inflammation, radiation injury, or other gland problems.

\medterm{Salivary Gland Tumors} Tumors arising in the salivary glands. Many are benign, especially in the parotid gland;
pleomorphic adenoma is a common benign type. Malignant types include mucoepidermoid,
adenoid cystic, and acinic cell carcinoma.

\synonyms
Tumor, Salivary Gland

\medterm{Salivary Glands} Glands that make saliva, including the parotid, submandibular, and sublingual glands. Saliva moistens food, begins digestion, protects teeth, helps swallowing, and supports taste and mouth health.
\medterm{Salivation} Production and release of saliva by the salivary glands. Saliva moistens food, begins digestion, protects teeth, and helps swallowing; excessive or reduced salivation can have many causes.

\medterm{Salk Vaccine} The inactivated poliovirus vaccine, given by injection to protect against poliomyelitis. It contains killed virus and cannot cause polio; it differs from the older oral live-attenuated vaccine.
\medterm{Salmon Patches} Flat pink or red capillary birthmarks, commonly on the eyelids, forehead, or nape; many fade during infancy.
\medterm{Salmonella} A genus of Gram-negative, facultatively anaerobic, rod-shaped bacteria in the Enterobacteriaceae family.
It includes non-typhoidal strains that cause gastroenteritis and typhoidal strains that can
cause enteric fever.


\medterm{Salmonella Choleraesuis} A Salmonella serotype adapted to pigs that can cause severe, invasive infection in humans, including bloodstream infection and pneumonia.


\medterm{Salmonella Enterica} A major Salmonella species causing human disease. It can produce diarrhea and fever, enteric fever, or bloodstream infection after contaminated food, water, animals, or infected people spread the bacteria.
\medterm{Salmonella Enterocolitis} Salmonella enterocolitis is intestinal infection causing diarrhea, fever, abdominal cramps, and sometimes bacteremia after contaminated food, water, or animal exposure.

\medterm{Salmonella Food Poisoning (Salmonellosis)} Alternate index label for Salmonellosis.


\medterm{Salmonella Infection} An infection caused by Salmonella bacteria, usually acquired through contaminated food, water, animals, or hands. It commonly causes diarrhea, fever, cramps, and vomiting but can spread into blood or other organs.

\medterm{Salmonella Infections} Infections caused by Salmonella bacteria, ranging from gastroenteritis to enteric fever or invasive bloodstream disease.

\medterm{Salmonella Typhi} The bacterium that causes typhoid fever, usually acquired through contaminated food or water and capable of spreading through the bloodstream.

\medterm{Salmonella Typhimurium} A common Salmonella serotype that usually causes foodborne gastroenteritis but can cause invasive disease in vulnerable people.

\medterm{Salmonella Typhisuis} A Salmonella serotype mainly associated with swine; human infection is uncommon but may cause enteric or invasive disease.

\medterm{Salmonella Vaccines} Vaccines used mainly to prevent typhoid fever, including typhoid conjugate, Vi polysaccharide, and oral live-attenuated vaccines; routine vaccination does not prevent all Salmonella gastroenteritis.

\medterm{Salmonellosis} Infection with Salmonella bacteria (non-typhoidal), typically acquired from contaminated
food (eggs, poultry, meat). Presents with acute gastroenteritis: diarrhea, fever, and
abdominal cramps within 6-72 hours of ingestion. Self-limited in immunocompetent
patients. Bacteremia can occur, especially in immunocompromised.

\synonyms
Salmonella Food Poisoning (Salmonellosis)

\medterm{Salpingectomy} Surgical removal of a fallopian tube. It may be used to treat a tubal ectopic pregnancy,
hydrosalpinx, or other tubal disease. The approach and whether the tube is removed depend
on the clinical condition, fertility plans, and the state of the opposite tube.

\medterm{Salpingitis} Inflammation or infection of one or both fallopian tubes, usually as part of pelvic inflammatory disease spreading upward from the vagina or cervix. It may cause lower-abdominal pain, fever, unusual discharge, infertility, or pregnancy outside the uterus.

\medterm{Salpingitis Isthmica Nodosa} Small thickened areas and side pockets in the narrow part of a fallopian tube, often after inflammation. They can narrow the tube, making pregnancy harder to achieve and increasing the risk of an ectopic pregnancy.

\synonyms
SIN

\medterm{Salpingo-Oophorectomy} Surgery to remove one or both fallopian tubes and ovaries. It may treat cancer, an ovarian or tubal problem, or reduce cancer risk in people with certain inherited gene changes. Removing both ovaries causes permanent loss of ovarian hormone production.

\medterm{Salpingo-OöPhorectomy} Surgical removal of a fallopian tube and ovary, also called salpingo-oophorectomy.
\medterm{Salpingography} Contrast imaging of the fallopian tubes, usually performed as part of hysterosalpingography to assess whether the tubes are open and to examine the uterine cavity.

\medterm{Salpingostomy (Salpingotomy)} An operation that opens a fallopian tube to remove an ectopic pregnancy while leaving the tube in place. It may preserve fertility in selected cases, but persistent pregnancy tissue, bleeding, scarring, or another ectopic pregnancy can occur.

\synonyms
Salpingotomy

\medterm{Salt} An ionic chemical compound; dietary sodium chloride is essential in small amounts but excess intake can increase blood pressure.

\medterm{Saltatory Conduction} The rapid movement of an electrical signal along a nerve fiber covered by myelin. The signal travels between small uncovered gaps in the covering, allowing fast communication with other nerves and muscles.

\medterm{Salter-Harris Fracture} A break involving the growth plate of a child or adolescent’s bone. The fracture pattern helps predict the risk of disturbed growth; pain, swelling, and inability to use the limb require prompt assessment and follow-up.

\medterm{Salvage Therapy} Treatment used when a disease, especially cancer, has not improved with earlier treatment or has returned. The choice depends on the disease, previous treatments, overall health, and treatment goals.

\medterm{Sample Fixation} Preserving a tissue or cell sample soon after collection so it does not decay or change shape. The preserved sample can then be processed, stained, and examined under a microscope for diagnosis.


\medterm{Sandfly Fever} An acute viral illness transmitted by phlebotomine sandflies, causing fever, headache, muscle pain, and sometimes rash.
\medterm{Sandhoff Disease} A rare inherited disorder in which fatty substances build up in nerve cells. It usually begins in infancy or childhood with loss of skills, weakness, seizures, movement problems, and sometimes an enlarged liver or spleen.

\medterm{Sandpaper Tooth} A tooth with a rough, grainy enamel surface that can hold plaque and feel uneven to the tongue. It may result from weak enamel, early mineral loss, fluorosis, or wear and should be checked by a dental professional.

\medterm{Sanfilippo Syndrome} A rare inherited disorder in which the body cannot break down part of a complex sugar. The buildup mainly affects the brain and causes progressive developmental and behavioral changes.
\synonyms
Mucopolysaccharidosis Type III (MPS III)


\medterm{Sanguineous} Containing or resembling fresh blood. The word often describes bright-red drainage from a wound; a large or rapidly increasing amount after surgery may indicate active bleeding and needs urgent assessment.
\medterm{Sanitation} Public-health measures that safely manage human waste, water, food, hygiene, and environmental contamination to prevent infection and disease transmission.

\medterm{Sano Shunt} A temporary tube connecting the lower-right heart chamber to the lung artery during the first operation for some babies with severe left-sided heart defects. It sends blood to the lungs while the heart is reconstructed; blockage or abnormal flow can be dangerous.

\synonyms
RV-PA Conduit; Norwood-Sano Modification

\medterm{Saphenous} Relating to the superficial veins of the lower limb, especially the great or small saphenous vein.
\medterm{Saphenous Vein} A superficial vein of the lower limb that drains the foot and leg. The great and small saphenous
veins are clinically important in venous disease and may be used as graft material in selected
vascular or cardiac bypass procedures.

\medterm{Small Saphenous Vein} A superficial vein running up the back of the calf and usually joining the popliteal vein behind the knee. It helps return blood from the foot and calf to the heart.
\medterm{SAPHO Syndrome} A rare inflammatory condition combining bone or joint pain with acne or pus-filled skin lesions. It often affects the breastbone, collarbone, spine, or jaw and can resemble infection or cancer, so diagnosis and treatment require specialist assessment.

\medterm{Sapovirus} A virus that can cause vomiting and watery diarrhea, especially in children. It spreads through contaminated food, water, surfaces, or close contact.

\medterm{Sapovirus Infection} An acute enteric infection caused by sapoviruses, positive-sense single-stranded RNA viruses in the family Caliciviridae that infect humans across all age groups. Clinical manifestations include acute watery diarrhea, nausea, vomiting, abdominal cramping, and low-grade fever, often responsible for outbreaks of foodborne or waterborne gastroenteritis in daycares, schools, and hospitals that resolve with supportive rehydration within 48 to 72 hours.

\synonyms
Sapovirus Gastroenteritis; Human Sapovirus Infection

\medterm{Sapporo Virus} An older name for sapovirus, a virus that can cause vomiting and watery diarrhea.

\medterm{Sapropterin} A synthetic form of tetrahydrobiopterin used for some people with phenylketonuria and certain tetrahydrobiopterin deficiencies; response and diet requirements vary.

\medterm{Saquinavir} An HIV protease inhibitor used with other antiretroviral medicines; drug interactions, gastrointestinal effects, and cardiovascular concerns limit its use.

\medterm{Saquinavir Mesylate} The mesylate salt form of saquinavir, an HIV medicine that blocks viral protein processing. It is used with other antiretroviral drugs, but interactions, liver problems, and dangerous heart-rhythm changes require careful monitoring.




\medterm{Sarcocystis Hominis} A protozoan parasite acquired by eating undercooked beef; people may develop intestinal symptoms, while cattle serve as an intermediate host.

\medterm{Sarcocystis Suihominis} A protozoan parasite acquired by eating undercooked pork that can cause intestinal sarcocystosis; pigs serve as an intermediate host.

\medterm{Sarcocystosis} Parasitic infection by Sarcocystis species, acquired from undercooked meat or contaminated food and sometimes causing muscle pain, fever, diarrhea, or no symptoms.

\medterm{Sarcoglycan Complex} A group of proteins that helps strengthen the outer covering of muscle cells during movement. Inherited problems with these proteins can cause progressive muscle weakness and sometimes heart-muscle disease.

\medterm{Sarcoglycanopathy} A group of inherited muscle diseases caused by changes in sarcoglycan genes. They usually cause progressive weakness of the hips and shoulders and may affect the heart.
\synonyms
Sarcoglycan-Deficient Limb-Girdle Muscular Dystrophy; LGMD R3-R6

\medterm{Sarcoidosis} An inflammatory disease in which clusters of immune cells form in one or more organs, most often the lungs and nearby lymph nodes. It can also affect the skin, eyes, heart, nervous system, or other organs.

\synonyms
Arthritis Sarcoid (Sarcoidosis)

\medterm{Sarcoidosis Skin Manifestations} Skin involvement in sarcoidosis can cause firm bumps, raised or discolored patches, scars that change, or lupus pernio on the nose or face. The appearance varies and may require a skin examination or biopsy.

\medterm{Sarcolemma} The thin outer covering of a muscle cell. It carries electrical signals that start contraction and helps maintain the cell’s structure and internal chemical balance.

\medterm{Sarcoma} A malignant tumor arising from connective or supporting tissues, including bone, cartilage, fat,
muscle, and blood vessels. Examples include osteosarcoma, liposarcoma, and rhabdomyosarcoma.

\medterm{Sarcoma, Bone Cancer} Alternate index label; see \textit{Bone cancer}.

\medterm{Sarcoma, Soft Tissue} Alternate index label; see \textit{Soft tissue sarcoma}.

\medterm{Sarcomere} The smallest working unit that makes skeletal and heart muscle contract. It contains protein threads that slide past one another, shortening the muscle and producing force.

\medterm{Sarcopenia} Loss of muscle strength and muscle mass that reduces physical function, often with aging but also with inactivity, poor nutrition, or chronic disease. It can cause weakness, slower walking, falls, and difficulty with daily activities.

\medterm{Sarcopenic Obesity} Excess body fat together with low muscle mass, strength, or physical ability. It can increase the risk of falls, disability, metabolic disease, and complications after surgery.

\medterm{Sarcoplasm} The material inside a muscle cell surrounding its specialized structures. It contains energy-producing parts, stored fuel, oxygen-holding myoglobin, and the protein machinery that produces contraction.

\medterm{Sarcoplasmic Reticulum} A network inside muscle cells that stores calcium. It releases calcium to start contraction and takes calcium back up so the muscle can relax.

\medterm{Sarcoptes} The mite genus containing Sarcoptes scabiei, the cause of scabies.
\medterm{Sarcoptes Scabiei} The microscopic mite that causes scabies. It burrows into the outer layer of skin and lays eggs, producing intense itching and a rash or thin lines. It spreads mainly through prolonged close skin contact.

\medterm{Sarcoptic Mange (Scabies)} Alternate index label for Scabies.




\medterm{Sargramostim} A recombinant granulocyte-macrophage colony-stimulating factor used to promote white-cell recovery after selected cancer treatments or transplantation; bone pain and fever may occur.

\medterm{Sarnat Staging} A three-level examination used to describe the severity of brain dysfunction in a newborn after reduced oxygen or blood flow. Mild disease causes increased alertness, moderate disease causes poor responsiveness and seizures, and severe disease causes coma and weak or absent reflexes; the stage helps guide urgent treatment and prognosis.

\synonyms
Sarnat and Sarnat Grading; Sarnat Score; HIE Staging

\medterm{SARS} Severe acute respiratory syndrome, a serious lung infection caused by the SARS coronavirus. It can cause fever, cough, breathing difficulty, and pneumonia and spreads mainly through close contact with respiratory droplets; suspected cases require public-health and medical guidance.
\medterm{SARS Virus} The coronavirus that causes severe acute respiratory syndrome, a serious respiratory infection that can spread through respiratory droplets and close contact.

\medterm{SARS-CoV-2} The coronavirus that causes COVID-19. It spreads mainly through particles released when an infected person breathes, talks, coughs, or sneezes and can cause anything from no symptoms to severe lung or other organ disease.

\medterm{Sassafras Oil Overdose} Poisoning from swallowing or otherwise being exposed to too much sassafras oil. It may cause nausea, vomiting, seizures, breathing problems, liver injury, or other serious effects and requires urgent medical advice.

\medterm{Satellite DNA} Repeated stretches of DNA found mainly in tightly packed parts of chromosomes. They help organize chromosome structure and can vary between people, usually without providing instructions for making proteins.

\medterm{Satellite Virus} A small virus-like infectious agent that can multiply only when a separate helper virus is present. It uses the helper’s machinery and may change how severe or persistent the combined infection becomes.

\medterm{Satiation} The feeling of fullness that terminates a meal, produced by gastrointestinal stretch, nutrients, hormones, and brain signals regulating food intake.

\medterm{Satiety} The feeling of fullness that reduces the desire to eat after a meal. It is influenced by stomach stretching, food composition, hormones, and brain signals; persistent early fullness can indicate illness.

\medterm{Satiety - Early} Feeling full sooner than expected while eating, sometimes after only a small amount of food. It may result from stomach or digestive disease and should be assessed if persistent or associated with weight loss.




\medterm{Saturated Fat} A dietary fat whose fatty-acid chains contain no carbon-carbon double bonds. It is found in foods such as fatty meat, butter, full-fat dairy, coconut oil, and palm oil; replacing some saturated fat with unsaturated fat can help lower LDL cholesterol and cardiovascular risk.


\medterm{Saturation} A state in which available binding, transport, or processing capacity is nearly full. Once capacity is reached, extra medicine or substance may remain in the blood or produce a disproportionately larger effect.

\medterm{Sausage Digit} Diffuse swelling of an entire finger or toe, often caused by inflammation in joints, tendons, and surrounding tissues.

\synonyms
Dactylitis


\medterm{Saxitoxin} A powerful poison made by some marine algae that can build up in shellfish. Eating contaminated seafood may cause tingling, weakness, trouble breathing, or paralysis.

\medterm{SBBO} Alternate index label; see \textit{Small intestinal bacterial overgrowth}.

\medterm{Scab} A crust of dried blood, serum, pus, or tissue fluid over a healing skin injury.
\medterm{Scabies} An infestation caused by tiny mites that burrow into the skin, causing intense itching that is often worse at night and a rash or thin burrows. It spreads mainly through prolonged close skin-to-skin contact.

\synonyms
Sarcoptic Mange (Scabies)


\medterm{SCAD} Alternate index label; see \textit{Spontaneous coronary artery dissection}.

\medterm{Scaffold} A three-dimensional framework used in tissue engineering to hold cells in place while they grow and organize into new tissue. It may be made from natural or synthetic material and may gradually dissolve.

\medterm{Scaffold Seeding} Placing living cells onto or inside a three-dimensional framework so they can attach, multiply, and begin forming tissue in laboratory or regenerative-medicine research.

\medterm{Scalded Mouth Syndrome} Alternate index label; see \textit{Burning mouth syndrome}.

\medterm{Scalded Skin Syndrome} Alternate index label; see \textit{Scalded-Skin Syndrome}.

\medterm{Scalded-Skin Syndrome} A serious infection, usually in young children, in which toxins from Staphylococcus bacteria make the outer skin layer separate. The skin becomes painful, red, and easily blistered or peeled; prompt treatment is needed to prevent dehydration and infection.
\medterm{Scalds} Burns caused by hot liquids or steam. They may damage only the surface skin or extend into deeper tissue; severity depends on depth, size, location, and age. Large, deep, facial, genital, hand, or circumferential burns need urgent specialist assessment, and unusual patterns in children require safeguarding evaluation.
\medterm{Scale} Visible flakes or plates of stratum corneum shed from the skin surface, commonly caused by abnormal keratinization, inflammation, infection, or dryness.

\medterm{Scales} Scales are flakes or plates of shed skin on the surface, occurring with dryness, eczema, psoriasis, fungal infection, or other dermatologic conditions.

\medterm{Scaling} Scaling is flaking or shedding of the outer skin layer and may occur with dryness, eczema, psoriasis, fungal infection, or other skin disease.

\medterm{Scaling And Root Planing} A deep dental cleaning that removes hardened plaque and bacteria above and below the gumline and smooths tooth roots. It reduces gum inflammation and may slow periodontal disease, but severe damage may need additional treatment.

\medterm{Scalp} The skin and underlying tissues covering the skull. It contains hair follicles, glands, blood vessels, nerves, and connective tissue and can be affected by injury, infection, or skin disease.
\medterm{Scalp And Body (Ringworm)} Alternate index label for Ringworm.

\medterm{Scalp Ringworm} A contagious fungal infection of the scalp, also called tinea capitis. It can cause scaling, itching, broken hairs, bald patches, or swollen tender areas and usually needs prescription antifungal medicine taken by mouth to prevent persistent hair loss or spread.

\medterm{Scalp Psoriasis} Psoriasis affecting the scalp, causing sharply demarcated red plaques covered by silvery scale that may extend beyond
the hairline. It can itch, crack, or bleed and is not contagious. Treatment may include medicated shampoos, topical
corticosteroids, vitamin-D analogues, or systemic therapy for extensive disease.

\synonyms
Psoriasis Of The Scalp (Scalp Psoriasis)

\medterm{Scalpel} A sharp surgical knife used to make precise incisions or cuts in tissue. It may be a reusable handle with a detachable blade or a disposable single-piece instrument. Blade designs and numbers differ in size and shape; for example, a No. 10 blade is commonly used for larger incisions, a No. 11 for puncture or stab incisions, and a No. 15 for shorter precise cuts.

\synonyms
Surgical Knife

\medterm{Scalpel Handle} The reusable handle that holds a detachable scalpel blade. Different handle and blade sizes are chosen for the type and precision of the incision.

\medterm{Scan} An imaging examination that creates pictures of body structures or functions using methods such as ultrasound, X-rays, CT, MRI, or radioactive tracers. The method and interpretation depend on the clinical question.

\medterm{Scan Gallium} A nuclear-medicine scan using radiolabeled gallium to identify areas of inflammation, infection, or some tumors; uptake is interpreted with anatomy and clinical context.

\medterm{Scanning} Systematic examination of a body region or specimen using imaging, microscopy, or another detection method to locate, characterize, or monitor abnormalities.

\medterm{Scanning Electron Microscopy} A method that scans a specimen’s surface with an electron beam to produce a highly magnified, three-dimensional-looking image of its shape and surface details.

\medterm{Scanning Speech} Speech that is unusually slow, broken into separate syllables, and spoken with uneven rhythm or emphasis. It can occur when disease affects the cerebellum or other parts of the nervous system.
\medterm{Scanning Techniques} Methods used to create images or measurements of the body or a sample. Examples include ultrasound, X-ray, CT, MRI, and nuclear-medicine scans; the method is chosen according to the clinical question.
\medterm{Scaphoid} A small wrist bone on the thumb side, commonly injured when someone falls onto an outstretched hand. Pain in the hollow just behind the thumb may indicate a fracture, which sometimes heals poorly because part of the bone has limited blood flow.

\medterm{Scaphoid Abdomen} A scaphoid abdomen is an unusually hollowed abdominal contour, classically seen with severe diaphragmatic hernia in a newborn but also possible with marked wasting.

\medterm{Scaphoid Bone} The largest proximal-row carpal bone on the thumb side of the wrist; its waist is vulnerable to fracture and avascular necrosis because of its blood supply.

\medterm{Scaphoid Fractures} Breaks in the small wrist bone near the base of the thumb, often after falling onto an outstretched hand. Pain may be felt in the hollow beside the thumb; some fractures heal poorly because part of the bone has a limited blood supply.

\medterm{Scaphycephaly} A long, narrow head shape caused when the major seam running from front to back of the skull closes too early. It is a form of craniosynostosis and may require specialist assessment to protect brain and skull development.

\medterm{Scapula} The shoulder blade, a flat triangular bone at the back of the chest. It connects the upper arm to the collarbone and provides attachment for muscles that move and stabilize the shoulder.

\medterm{Scapulothoracic Joint} The functional gliding connection between the shoulder blade and the back of the chest wall. It is not a true joint but works with nearby joints to position and move the arm.

\medterm{Scar} Fibrous tissue formed during healing after injury, inflammation, or surgery.
\medterm{Scar Excessive (Keloid)} Alternate index label for Keloid.

\medterm{Scar Management} Measures used to reduce symptoms and abnormal scar formation. Options include silicone
products, pressure therapy, massage, corticosteroid injections, or other specialist treatments, depending on the scar
and its cause.

\medterm{Scar Revision} A procedure that improves the appearance, contour, flexibility, or symptoms of a scar through excision, rearrangement, grafting, laser, injection, or other selected techniques.


\medterm{Scarlatina} Scarlet fever, an illness caused by toxin-producing group A Streptococcus and characterized by sore throat, fever, and a fine rough-textured rash.

\synonyms
Scarlet Fever (Scarlatina)

\medterm{Scarlatina (Scarlet Fever (Scarlatina))} Alternate index label for Scarlet Fever (Scarlatina).

\medterm{Scarlet Fever} An illness caused by toxin-producing group A Streptococcus. It usually causes sore throat, fever, and a widespread fine, rough-textured rash; a red tongue and pale skin around the mouth may also occur.

\medterm{Scarlet Fever (Scarlatina)} Alternate index label for Scarlatina.


\medterm{Scarring Alopecia} Permanent hair loss caused when inflammation or injury destroys hair follicles and replaces them with scar tissue. Early scalp redness, scaling, itching, pain, or spreading hair loss needs dermatologic assessment because lost follicles cannot regrow hair.

\medterm{Scars} A scar is fibrous tissue formed during healing after injury, surgery, inflammation, or burns. Scars may be flat, depressed, raised, or contractured and can itch, hurt, restrict movement, or affect appearance; treatment depends on maturity and type.

\medterm{Scatter Correction} A medical-imaging method that estimates and reduces errors caused when detected radiation changes direction after passing through tissue. It can improve the accuracy of a nuclear-medicine image or measured activity.

\medterm{Scatter Events} Radiation particles that change direction after interacting with tissue or another material. In medical imaging they can place signals in the wrong location or reduce contrast, so scanners use design and correction methods to limit their effect.
\medterm{Scavenger Receptor} A protein on or inside a cell that recognizes and helps remove altered fats, germs, cell debris, or dying cells. It supports early immune defense but can also help fatty material build up in artery walls.

\medterm{Scavenging System} Equipment that collects and removes waste anesthetic gases from an operating room. It protects staff from repeated exposure and may use a vacuum system or passive venting to carry gases safely away.

\medterm{Scedosporium} A genus of environmental molds that can cause difficult-to-treat lung, sinus, skin, bone, or disseminated infections, especially in immunocompromised people.

\medterm{Scedosporium Apiospermum} A mold that can cause localized or disseminated infection, including mycetoma, pneumonia, sinusitis, and brain abscesses; treatment often requires specialist antifungal care.

\medterm{Scedosporium Prolifican} An environmental mold, now commonly classified as Lomentospora prolificans, that causes severe disseminated infection and is often resistant to antifungal drugs.

\medterm{Schaalia Radingae} A normally slow-growing anaerobic bacterium of the human mouth and skin that can rarely cause abscesses or other invasive infections.

\medterm{Schaalia Turicensis} A filamentous anaerobic bacterium that may be part of normal flora but can cause actinomycosis-like abscesses and other opportunistic infections.

\medterm{Schatzki Ring} A thin ring of tissue narrowing the lower esophagus, often near a hiatus hernia. It may cause occasional difficulty swallowing solid food or a piece of food becoming stuck; treatment can widen the narrowed area.

\synonyms
Ring Schatzki (Schatzki Ring)
Strictures Esophagus (Schatzki Ring)

\medterm{Schedule I} A United States controlled-substance category for drugs with high misuse potential and no accepted medical use
under federal law.

\medterm{Schedule II} A United States legal category for medicines with accepted medical uses but high misuse, dependence, or overdose potential. Prescribing, dispensing, refills, and storage are subject to strict rules.
\medterm{Schedule III} A United States legal category for medicines with accepted medical uses and less misuse potential than Schedule II substances, although dependence and misuse remain possible. Rules vary by medicine and jurisdiction.
\medterm{Schedule IV} A United States controlled-substance category for drugs with accepted medical uses and lower misuse potential
than Schedule III drugs.

\medterm{Schedule V} A United States controlled-substance category for drugs with the lowest level of misuse potential among the
federal schedules.

\medterm{Scheduled Notification} A planned alert or message delivered at a specified time to remind a person about an appointment, medicine, monitoring task, or health-related action.

\medterm{Scheuermann Disease} A juvenile osteochondrosis of the spine in which vertebral end plates develop irregularities and anterior wedging during skeletal growth.
It commonly produces a rigid thoracic kyphosis, although the thoracolumbar or lumbar spine may also be involved.
Symptoms can include back pain, fatigue, limited spinal motion, and postural deformity that becomes more apparent during adolescence.
Diagnosis is based on clinical examination and characteristic radiographs; management depends on deformity severity, skeletal maturity, pain, and progression.

\medterm{Scheuermann's Disease} A growth disorder beginning in adolescence that causes several spinal bones to become wedge-shaped. This can produce a lasting rounded upper back, stiffness, pain, and sometimes breathing or activity limitations.
\medterm{Schiff Base} A chemical compound formed when a primary amine reacts with an aldehyde or ketone; Schiff bases occur in biochemistry and drug chemistry.

\medterm{Schiller Test} An examination of the cervix in which iodine solution is applied to highlight areas of abnormal tissue. Areas that do not stain may need further examination or biopsy; the test alone cannot diagnose cancer.

\medterm{Schilling Test} An older test that measured how well the body absorbed vitamin B12, sometimes used to investigate pernicious anemia or intestinal disease. It is rarely performed now because blood tests and antibody testing have largely replaced it.

\medterm{Schirmer Test} A test in which standardized filter paper strips measure tear production over a set time, commonly used to evaluate aqueous-deficient dry eye.

\medterm{Schistocyte Count} Measurement of fragmented red blood cells on a blood smear, supporting evaluation of microangiopathic hemolysis, mechanical valve damage, disseminated intravascular coagulation, or severe burns.

\medterm{Schistocytes} Broken, irregular, or helmet-shaped red blood cells seen on a blood-smear test. They form when red cells are damaged in abnormal small vessels, by artificial heart valves, or in other illnesses and may signal serious red-cell destruction.

\medterm{Schistomicide} A medicine that kills parasitic blood flukes responsible for schistosomiasis. Treatment choice depends on the parasite species and the organs affected, and medical advice is needed because infection may cause lasting damage.
\medterm{Schistosoma} A group of parasitic worms that live in blood vessels. Their young forms can enter through skin exposed to contaminated fresh water and later cause intestinal, urinary, liver, or other disease.

\medterm{Schistosoma Haematobium} A parasitic blood fluke that can cause urinary schistosomiasis, including hematuria and chronic
injury of the urinary tract.

\medterm{Schistosoma Intercalatum} A parasitic blood fluke that causes intestinal schistosomiasis, mainly in parts of Central and West Africa, after freshwater exposure.

\medterm{Schistosomiasis (Bilharzia)} An infection caused by parasitic worms called \textit{Schistosoma}. The worms live in freshwater snails, and their larvae can penetrate the skin when a person swims, wades, or works in contaminated freshwater. The infection may cause an itchy rash, fever, abdominal pain, diarrhea, blood in the stool or urine, and painful urination. Long-term infection can damage the intestine, liver, urinary tract, or other organs.

\synonyms
Bilharzia

\medterm{Schizencephaly} A rare condition present from birth in which a fluid-filled cleft extends through part of the brain. It results from abnormal early brain development and may cause seizures, weakness, movement problems, or delayed development, depending on the cleft’s size and location.

\medterm{Schizoaffective Disorder} A mental-health disorder in which a person has symptoms of schizophrenia, such as hallucinations or disorganized thinking, together with major depression or mania. Psychotic symptoms also occur for a period without a major mood episode.

\medterm{Schizogony} A way some single-celled parasites reproduce without mating. One parasite cell grows and divides internally to produce many new cells, which are then released to infect other cells.
\medterm{Schizoid Personality Disorder} A long-standing pattern of limited interest in close relationships and restricted emotional expression. It differs from schizophrenia because hallucinations and fixed delusions are not central features.
\synonyms
Personality Disorder, Schizoid

\medterm{Schizophrenia} A serious, long-term mental disorder that can affect a person’s thoughts, perceptions, emotions, behavior, and ability to function. Symptoms are often grouped as psychotic, negative, and cognitive. Psychotic symptoms include hallucinations, such as hearing or seeing things that others do not, delusions, which are firmly held beliefs not supported by shared evidence, and disorganized thought or speech. Negative symptoms may include reduced emotional expression, motivation, pleasure, or social interest. Cognitive symptoms may affect attention, memory, planning, and decision-making. Symptoms and their severity vary between people and may come and go or persist over time.

\medterm{Schizophrenia, Adolescent} Alternate index label; see \textit{Childhood schizophrenia}.

\medterm{Schizophrenia, Childhood} Alternate index label; see \textit{Childhood schizophrenia}.

\medterm{Schizophrenia, Juvenile} Alternate index label; see \textit{Childhood schizophrenia}.

\medterm{Schizotypal Personality Disorder} A long-term pattern of discomfort with close relationships, unusual beliefs or perceptions, unusual speech or behavior, and difficulty interpreting social situations. It is different from schizophrenia because persistent hallucinations and major loss of reality testing are not typical.

\synonyms
Personality Disorder, Schizotypal

\medterm{Schlemm’S Canal} A circular drainage channel at the edge of the cornea that carries fluid from the front of the eye into blood vessels. Poor drainage can raise pressure inside the eye and contribute to glaucoma.

\synonyms
Canal of Schlemm

\medterm{Schnitzler Syndrome} A rare inflammatory disease causing a long-lasting hives-like rash, repeated fever, bone or joint pain, and raised inflammation tests. It is usually associated with an abnormal IgM blood protein and requires specialist evaluation.

\medterm{School Age Test Or Procedure Preparation} Preparing a school-age child for a test or procedure includes age-appropriate explanations, comfort measures, caregiver support, and following any fasting or medicine instructions.

\medterm{School Sores (Impetigo)} Alternate index label for Impetigo.

\medterm{School-Age Children Development} School-age development includes advances in growth, coordination, language, learning, independence, social relationships, and emotional regulation, with broad normal variation.

\medterm{School-Based Rehabilitation} Therapy and other support services provided in a school setting to help a child participate in education and daily school activities.

\medterm{Schwann Cell} A supporting cell of the peripheral nervous system. It wraps around many nerve fibers to form their insulating myelin covering and helps maintain or repair peripheral nerves.
\medterm{Schwann Cells} Glial cells of the peripheral nervous system that form myelin around many axons and support the repair and maintenance of peripheral nerves.

\medterm{Schwannoma} A usually noncancerous, slow-growing tumor arising from cells that support nerves. It may cause a lump, pain, numbness, or weakness; a tumor involving the hearing and balance nerve can cause hearing loss, ringing, or dizziness.

\synonyms
Neurilemmoma

\medterm{Schwannomatosis} A condition in which multiple schwannomas, usually slow-growing tumors of the covering of nerves, develop. The main problem is often long-lasting pain, although numbness, weakness, or other nerve symptoms can occur depending on tumor location.

\synonyms
Neurinomatosis

\medterm{SCI} Alternate index label; see \textit{Spinal cord injury}.

\medterm{Sciatic Nerve} The body's largest nerve, running from the lower spine through the buttock and down the back of the leg. It supplies movement and sensation to much of the leg and foot.

\medterm{Sciatic Neuritis (Sciatica)} Alternate index label for Sciatica.

\medterm{Sciatica} Pain, tingling, numbness, or weakness that travels from the lower back or buttock down the leg, usually because a nerve root is irritated by a disk problem or narrowing of the spine. Progressive weakness, numbness around the inner thighs or genitals, or new loss of bladder or bowel control may indicate severe nerve compression and requires urgent assessment.

\synonyms
Sciatic Neuritis (Sciatica)

\medterm{Scimitar Syndrome} A birth defect in which some veins from the right lung drain to the wrong large vein instead of directly to the heart. The right lung may be smaller, and the condition can cause repeated chest infections, breathing problems, high pressure in the lung arteries, or heart failure.

\medterm{Scintigraphy} Nuclear-medicine imaging in which a small amount of radioactive tracer is given and a camera records where it collects. The pattern can show the function, blood flow, or activity of organs such as bone, thyroid, kidneys, or heart.

\medterm{Scintillation} The brief release of light by a material after it absorbs radiation. Detectors measure these flashes to estimate radiation and create images in some medical scans.

\medterm{Scintillation Crystal} A material in some medical detectors that changes incoming radiation into tiny flashes of light. The flashes are measured to create images or estimate the amount of radioactive tracer present.

\medterm{Scintillations} Flashes, sparkles, or zigzag lines seen in the field of vision. They often occur during a migraine aura but may also result from an eye or brain problem; sudden new flashes with floaters or a curtain-like shadow need urgent eye assessment.

\medterm{Scirrhus} An older term for a hard cancer containing a large amount of scar-like tissue. It was used especially for some breast or stomach cancers and describes texture rather than a separate cancer type.
\medterm{Scissors} Instruments with two blades used to cut tissue, sutures, dressings, or other material. Surgical types vary in length, blade shape, and sharpness; examples include Mayo and Metzenbaum scissors.

\medterm{Scissors (Surgical)} Surgical instruments designed to cut tissue, sutures, dressings, or other materials. Different shapes and blade types are selected for delicate, deep, curved, or heavy-duty work.
\synonyms
Surgical scissors

\medterm{Scissors (Vascular)} A surgical instrument with two blades used to cut tissue, sutures, dressings, or other materials
during vascular procedures.

\synonyms
Vascular scissors

\medterm{Sclera} The sclera is the tough white outer coat of the eye that protects the globe and provides attachment for eye muscles.

\medterm{Scleral Buckling} A retinal-detachment repair procedure in which a silicone band or buckle is placed around the outside of the eye to support the retinal break and reduce traction.

\medterm{Scleredema Diabeticorum} Scleredema diabeticorum is thickening and hardening of the skin on the upper back and neck, usually associated with long-standing diabetes and obesity.

\medterm{Scleritis} Painful inflammation of the white outer coat of the eye. It is often associated with autoimmune disease and can damage the eye or threaten vision, so prompt eye assessment is needed.

\synonyms
Inflammation Sclera (Scleritis)

\medterm{Scleroderma} A group of autoimmune diseases that cause inflammation and too much collagen to build up, making skin and sometimes deeper tissues hard, thick, and tight. Localized scleroderma, also called morphea or linear scleroderma, mainly affects the skin and tissues just beneath it. Systemic scleroderma, also called systemic sclerosis, can affect the skin, small blood vessels, and internal organs such as the esophagus, lungs, kidneys, and heart. The cause is not fully known, and symptoms vary with the type and organs involved.

\synonyms
Arthritis Scleroderma (Scleroderma)
Sclerosis, Systemic


\medterm{Scleroma} A long-lasting infection, usually of the nose and upper airway, caused by Klebsiella rhinoscleromatis. It can produce scar tissue, blocked breathing, nasal discharge, voice changes, or facial deformity.

\medterm{Scleromyxedema} A rare long-term skin disorder in which many small, firm, waxy bumps and thickened skin develop, especially on the face, neck, and hands. It may also affect muscles, joints, swallowing, lungs, or other organs.

\medterm{Sclerosing Cholangitis} Inflammation and scarring that narrow or block the tubes carrying bile from the liver. It may cause itching, jaundice, abdominal pain, repeated infections, and liver damage; the cause guides treatment and follow-up.

\medterm{Sclerosing Mesenteritis} A rare inflammatory condition causing thickening and scarring in the tissue that supports the intestines. It may cause abdominal pain, nausea, weight loss, or bowel blockage, although some people have no symptoms.

\medterm{Sclerosing Peritonitis} Thickening and scarring of the lining inside the abdomen, sometimes after long-term peritoneal dialysis. The scar tissue can restrict the intestines and cause pain, vomiting, poor nutrition, or bowel blockage.
\medterm{Sclerosing Sialadenitis} Long-term inflammation that makes a salivary gland swollen and hard because scar-like tissue forms inside it. It may be related to an immune disorder and can cause dry mouth, pain, or repeated swelling.

\medterm{Sclerosis} Abnormal hardening or scarring of tissue, often from excess connective tissue, chronic inflammation, degeneration, or neurologic plaques; the affected organ determines the specific disease.

\medterm{Sclerosis, Multiple} Alternate index label; see \textit{Multiple sclerosis}.

\medterm{Sclerosis, Systemic} Alternate index label; see \textit{Scleroderma}.

\medterm{Sclerosteosis} A rare inherited condition in which bones, especially the skull, become abnormally thick. It may cause a broad or tall build, joined fingers, headaches, hearing loss, facial weakness, or pressure on nerves inside the skull.

\medterm{Sclerotherapy} A treatment in which an irritating or closing agent is injected into a blood vessel or body cavity to make it scar and seal, commonly used for varicose veins and selected vascular lesions.


\medterm{Scoliosis} An abnormal sideways curve of the spine that may look like the letter “C” or “S.” The bones of the spine may also rotate, changing the shape of the back or rib cage. In many children and teenagers, the cause is unknown; this is called idiopathic scoliosis. Other causes include differences present at birth, nerve or muscle conditions, spinal injury, or age-related changes. Mild scoliosis may cause no symptoms, while larger curves may cause uneven shoulders or hips, a prominent shoulder blade, back pain, reduced movement, or, rarely, breathing problems.

\synonyms
Spinal Curvature



\medterm{Scopolamine} Scopolamine is a muscarinic antagonist used mainly to prevent motion sickness and postoperative nausea; dry mouth, blurred vision, urinary retention, confusion, and other anticholinergic effects may occur.

\medterm{Scopolamine Hydrobromide} The hydrobromide salt of scopolamine, an anticholinergic medicine used mainly to prevent motion sickness and postoperative nausea.


\medterm{Scopulariopsis} Scopulariopsis is a mold genus that rarely causes opportunistic infection, especially in people with immune suppression or implanted devices.

\medterm{Scorbutic} Relating to scurvy caused by severe vitamin C deficiency. It can cause easy bruising, bleeding gums, poor wound healing, anemia, bone or joint pain, and weakness; treatment with vitamin C and an adequate diet usually leads to recovery.
\medterm{Scorpion Fish Sting} A painful injury from the venomous spines of a scorpion fish, causing severe local pain and swelling and sometimes whole-body symptoms. The wound needs urgent medical assessment because spines or venom may remain in the injury.

\medterm{Scorpion Sting} A venomous injury from a scorpion, usually causing immediate pain, burning, redness, or tingling. Some species can cause muscle twitching, abnormal eye movements, breathing difficulty, or dangerous autonomic symptoms.
\synonyms
Stings, Scorpion

\medterm{Scorpions} Venomous arachnids whose stings can cause local pain and, depending on the species, sweating, muscle spasms, breathing problems, or other serious symptoms.

\medterm{Scotoma} A blind or blurred spot within an otherwise usable field of vision. It may result from retinal disease, optic-nerve damage, migraine, stroke, or another eye or brain disorder.

\medterm{Scrape} A scrape is a superficial abrasion in which friction removes part of the skin surface; cleaning, protection, and monitoring for infection support healing.


\medterm{Screen Time And Children} Time children spend using televisions, computers, phones, tablets, or game devices; excessive or poorly timed use can affect sleep, activity, learning, mood, and family routines.

\medterm{Screening} Testing people without symptoms to identify a disease or an increased risk of disease. The appropriate test,
age range, interval, and follow-up depend on the condition, individual risk, and current
guidance; screening benefits must be balanced against harms.

\medterm{Screening (Geriatric)} Checking older adults who have no symptoms for diseases or risks that may benefit from early detection. Decisions should consider expected benefit, possible harm, life expectancy, health goals, and preferences; examples include falls, depression, osteoporosis, and selected cancers.

\synonyms
Geriatric screening; screening in older adults

\medterm{Screening And Diagnosis For HIV} Testing for human immunodeficiency virus using an antigen-antibody assay or nucleic-acid test, followed by confirmatory testing and assessment of immune status when positive.

\medterm{Screening Procedure} A test or examination offered to people without symptoms to detect disease or risk early enough for effective prevention or treatment; eligibility and intervals depend on evidence and risk.

\medterm{Screening Program} An organized public-health service that systematically invites a defined population to
undergo testing for an asymptomatic condition, with quality-assured follow-up,
diagnosis, treatment, and evaluation of benefits and harms.

\medterm{Screening Study} A planned investigation of an asymptomatic population to identify early disease, precursor lesions, or risk factors using a validated test and follow-up pathway.

\medterm{Screening Test} A test applied to people without symptoms or a known diagnosis to identify those at
increased risk of a condition. A screening result usually requires a diagnostic or
confirmatory evaluation before treatment is initiated.

\medterm{Screws} Threaded metal devices used to hold broken bones together or attach plates and other implants securely to bone.

\medterm{Script} A script is a written set of instructions or a prescription; in healthcare, a prescription specifies the medicine, dose, route, timing, and duration.

\medterm{Scrivener's Palsy} An older term for task-specific hand dystonia or cramp occurring during writing.
\medterm{Scrofula} Tuberculous infection of cervical lymph nodes, producing chronic neck swellings that may suppurate or form draining sinuses; nontuberculous mycobacteria can cause a similar presentation.

\medterm{Scrofulous} Scrofulous is an old term for tuberculous infection of lymph nodes, especially in the neck, which can cause chronic swelling and draining sinuses.

\medterm{Scrotal Calcinosis} Firm, usually painless calcium deposits in the skin of the scrotum that appear as one or more whitish or yellowish lumps. They may become tender, ulcerate, or release a chalky material; removal is considered if they cause symptoms.

\medterm{Scrotal Infection} Infection of the scrotal skin or underlying structures, which may cause redness, swelling, warmth, pain,
or drainage. Rapidly worsening pain, fever, discoloration, or severe swelling is an emergency
because serious soft-tissue infection must be excluded.

\medterm{Scrotal Masses} Lumps or swellings in the scrotum, arising from the skin, fluid around a testicle, veins, infection, hernia, or a testicular tumor. A new, firm, painless, or rapidly growing lump needs prompt examination and ultrasound.

\medterm{Scrotal Pain} Pain in the scrotum caused by infection, trauma, hernia, fluid collections, varicocele, testicular torsion, or other conditions affecting the testes and surrounding structures.

\medterm{Scrotal Swelling} Enlargement of the scrotum caused by fluid, inflammation, infection, hernia, varicocele, hydrocele, tumor, trauma, or testicular torsion; sudden pain is an emergency.

\medterm{Scrotal Ultrasound} Ultrasound imaging of the scrotum and testes used to evaluate pain, swelling, masses, torsion, infection, trauma, and blood flow.

\medterm{Scrotum} The scrotum is the external skin-and-muscle sac containing the testes and epididymides; its muscles help regulate testicular temperature, and swelling or pain may indicate infection, torsion, hernia, or tumor.


\medterm{Scrub In} To scrub in is to perform surgical hand antisepsis and put on sterile clothing and gloves before participating in a sterile procedure.

\medterm{Scrub Up} To scrub up means to clean and antiseptically prepare the hands and forearms before a sterile procedure.

\medterm{Scrubs} Scrubs are protective work garments commonly worn by healthcare staff; they reduce clothing contamination but do not replace hand hygiene or other precautions.

\medterm{Scrub Typhus} An acute febrile zoonotic infection caused by the intracellular bacterium Orientia tsutsugamushi and transmitted to humans through the bite of larval trombiculid mites (chiggers). Characteristic manifestations include high fever, severe headache, generalized lymphadenopathy, maculopapular rash, and a diagnostic necrotic skin lesion with a black crust (eschar) at the bite site; severe cases can progress to pneumonitis, acute kidney injury, myocarditis, or acute encephalitis syndrome.

\synonyms
Tsutsugamushi Disease; Mite-Borne Typhus

\medterm{Scurf} Fine flakes or scales shed from the skin, especially the scalp. It may occur with dry skin, dandruff, eczema, psoriasis, infection, or irritation and may be accompanied by itching or redness.
\medterm{Scurvy} A disease caused by severe vitamin C deficiency. It can cause bleeding or swollen gums, easy bruising,
small skin bleeds, poor wound healing, anemia, bone or joint pain, and fatigue.

\medterm{Sea-Blue Histiocyte} An immune cell that appears blue-green under a microscope because it contains stored fats or cell-waste pigment. It may be seen in some storage, blood, or metabolic disorders, but the finding alone does not identify a disease.

\medterm{Sea-Sickness} Motion sickness caused by conflicting signals from the inner ear, eyes, and body during travel on a boat or ship. It may cause nausea, vomiting, dizziness, sweating, headache, and pallor.
\medterm{Seasonal Affective Disorder} Depression that repeatedly begins during a particular season, most often autumn or winter, and improves at another time of year. Symptoms may include low mood, loss of interest, tiredness, oversleeping, increased appetite, and social withdrawal.

\medterm{Seasonal Affective Disorder Syndrome} Alternate index label for Seasonal Affective Disorder.
\medterm{Seasonal Allergies (Hay Fever)} Alternate index label for Hay Fever.

\medterm{Seasonal Allergy} Alternate index label; see \textit{Hay fever}.

\medterm{Sebaceous Adenoma} A sebaceous adenoma is a usually benign tumor of sebaceous glands that can appear as a yellowish skin nodule and may be associated with Muir-Torre syndrome.

\medterm{Sebaceous Carcinoma} An uncommon cancer arising from sebaceous, or oil-producing, glands. It often appears as a firm or yellowish eyelid or skin nodule and may require biopsy and specialist treatment.

\medterm{Sebaceous Cyst} A common but imprecise term for an epidermoid cyst, a keratin-filled cyst beneath the skin.

\medterm{Sebaceous Epithelioma} A benign sebaceous neoplasm, often presenting as a yellowish papule or nodule. Multiple or unusual lesions may warrant assessment for a DNA mismatch-repair syndrome such as Muir-Torre syndrome.

\medterm{Sebaceous Gland} A small gland in the skin that secretes sebum (oil) into hair follicles to lubricate and
waterproof the skin and hair.

\medterm{Sebaceous Gland Carcinoma} Sebaceous gland carcinoma is an uncommon skin cancer arising from oil glands, often appearing as a firm yellowish eyelid or skin lump.

\medterm{Sebaceous Glands} Skin glands that produce sebum, an oily substance usually released into hair follicles. They are most abundant on the face and scalp and contribute to acne when blocked or inflamed.
\medterm{Sebaceous Hyperplasia} A benign enlargement of sebaceous glands, presenting as small yellowish papules with central umbilication, usually on the face of adults.

\medterm{Seborrhea} Seborrhea is excessive or altered oil production that contributes to greasy scale and inflammation, commonly in seborrheic dermatitis of the scalp, face, or chest. Gentle cleansing, medicated shampoos, and topical anti-inflammatory or antifungal treatment may help.

\medterm{Seborrheic Dermatitis} A common long-term inflammation of oily skin, especially the scalp, eyebrows, sides of the nose, chest, and skin folds. It causes red or pink patches with greasy white or yellow scale; dandruff is a mild scalp form.

\synonyms
Dermatitis, Seborrheic

\medterm{Seborrheic Keratoses} Common, benign, epidermal skin tumors presenting as waxy, stuck-on, pigmented plaques
with a rough, verrucous surface, typically appearing on the trunk and face of
middle-aged and older adults. They are the most common benign skin neoplasm and are
caused by keratinocyte proliferation.

\medterm{Seborrheic Keratosis} A benign epidermal tumor appearing as stuck-on, waxy, hyperpigmented plaques with a
verrucous surface. They are extremely common in older adults and increase in number with
age. Seborrheic keratoses are not premalignant.

\synonyms
Keratosis, Seborrheic

\medterm{Seborrhoea} A chronic inflammatory skin condition affecting sebum-rich areas: scalp, face
(nasolabial folds, eyebrows), chest, and intertriginous areas. See Seborrheic
Dermatitis.

\medterm{Seborrhoeic} Relating to increased or altered activity of sebaceous glands or to areas rich in these glands, such as the scalp, face, chest, and upper back.

\medterm{Seborrhoeic Keratosis} A common harmless skin growth that may be brown, black, or light-colored with a waxy, rough, or stuck-on appearance. It is not contagious or cancerous, but changing lesions should be checked.

\medterm{Sebum} An oily substance made by skin glands that lubricates and protects the skin and hair. Excess production can contribute to oily skin, acne, and dandruff, while changes may accompany some skin disorders.
\medterm{Seckel Syndrome} A rare inherited condition causing severe growth restriction before and after birth, a very small head, characteristic facial features, and often developmental or intellectual disability. Bone, blood, and other birth differences may also occur.

\medterm{Secobarbital Overdose} Poisoning from taking too much secobarbital, a sedative medicine. It can cause extreme sleepiness, confusion, slowed breathing, coma, or death and requires emergency care.

\medterm{Second Malignancy} A new, separate cancer that develops after an earlier cancer. The risk may be influenced by inherited factors, previous radiation or chemotherapy, age, and other exposures; it is different from recurrence of the first cancer.

\medterm{Second Messenger Systems} Internal signaling systems that carry a message from a hormone, nerve chemical, or growth factor at the cell surface to processes inside the cell. They can change contraction, secretion, metabolism, gene activity, or cell growth.

\medterm{Second Messengers} Small molecules inside cells that pass along messages received from hormones, nerve chemicals, or other outside signals. They help control contraction, secretion, metabolism, gene activity, and cell growth.

\medterm{Second Opinion} An independent medical assessment by another qualified clinician to help confirm or clarify a diagnosis, prognosis, or treatment plan.

\medterm{Second Stage Of Labor} The part of childbirth that begins when the cervix is fully open and ends when the baby is born. It includes descent of the baby and pushing; its length varies with previous births, pain relief, the baby’s position, and the health of parent and baby.

\medterm{Second Trimester} The middle part of pregnancy, from 14 through 27 completed weeks. The fetus grows rapidly, movement is often felt, and ultrasound and blood screening may assess development and certain conditions; the timing and tests vary by pregnancy and local care.

\medterm{Second-Degree AV Block} A heart-rhythm problem in which some electrical signals from the upper chambers fail to reach the lower chambers. The heartbeat may be slow or irregular; certain patterns can progress to complete heart block and need urgent assessment.

\medterm{Second-Degree Burn} A burn that damages the outer skin and part of the layer beneath it. The area is usually painful, red, moist, and blistered; deeper burns may look pale and heal slowly. Large, deep, or sensitive-area burns need medical assessment.

\medterm{Second-Hand Smoke} Smoke inhaled from burning tobacco products or exhaled by a person who is smoking. It contains toxic
substances and increases health risks in exposed people.

\medterm{Second-Look Surgery} Second-look surgery is a planned or follow-up operation performed to inspect healing, remove remaining disease, assess treatment response, or manage complications.

\medterm{Secondary Adrenal Insufficiency} Reduced production of adrenal hormones because the pituitary or brain does not provide enough stimulation, often after long-term steroid treatment is stopped suddenly. It can cause fatigue, weakness, nausea, low blood pressure, and low blood sugar; skin darkening and salt loss are less typical than in primary adrenal failure.

\medterm{Secondary Amenorrhea} Absence of menstruation for at least three months in someone with previously regular cycles or six months with previously irregular cycles, caused by pregnancy, endocrine, ovarian, uterine, or systemic factors.

\medterm{Secondary Attack Rate} The proportion of susceptible people who become infected after contact with an initial case during a defined period. It helps estimate how readily an infection spreads in households, schools, healthcare settings, or other close-contact groups.

\medterm{Secondary Brain Cancer} Alternate index label; see \textit{Brain metastases}.

\medterm{Secondary Hyperparathyroidism} Overactivity of the parathyroid glands in response to low calcium, abnormal phosphate, vitamin-D deficiency, or kidney disease. It can weaken bones, cause bone pain, and contribute to calcium deposits in blood vessels.

\medterm{Secondary Hypertension} High blood pressure caused by another condition, medicine, or substance. Possible causes include kidney disease, hormone disorders, sleep apnea, narrowed arteries, pregnancy, steroids, stimulants, or certain pain medicines.

\medterm{Secondary Hypothyroidism} Low thyroid-hormone production because the pituitary or nearby brain area does not provide enough thyroid-stimulating signal. It may cause tiredness, feeling cold, constipation, dry skin, slow heart rate, or menstrual changes.

\medterm{Secondary Immune Response} A faster and stronger defense when the body meets a germ or foreign substance again. Memory immune cells from the first exposure respond quickly and usually produce antibodies and other defenses more effectively.

\medterm{Secondary Immunodeficiency} Reduced immune function caused by an acquired condition or exposure rather than a primary genetic disorder. Causes
include HIV infection, cancer, malnutrition, medicines, and chronic disease; the
manifestations depend on which immune functions are impaired.

\medterm{Secondary Infections} Infections that develop during or after another disease, injury, or treatment because of tissue damage, altered defenses, or immunosuppression.


\medterm{Secondary Osteoporosis} Bone loss caused by another illness, medication, or prolonged immobility.

\medterm{Secondary Parkinsonism} Secondary parkinsonism is slowed movement, rigidity, tremor, or postural instability caused by another condition, medicine, toxin, vascular injury, or structural brain disease.

\medterm{Secondary Prevention} Finding and treating a disease early to stop it worsening or causing complications. Examples include screening for some cancers, checking blood pressure, testing for early diabetes, and treating an abnormality before it produces serious symptoms.

\medterm{Secondary Sexual Characteristics} Physical changes that develop during puberty and are not directly part of the reproductive organs, such as breast development, facial or body hair, voice changes, and changes in body shape.

\medterm{Secondary Progressive Multiple Sclerosis} A phase of multiple sclerosis in which disability gradually worsens over time after an earlier relapsing pattern. Relapses may still occur, but progressive walking, strength, vision, bladder, thinking, or other problems can develop even between attacks.

\synonyms
SPMS; Secondary-progressive MS

\medterm{Secondary Syphilis} The stage of syphilis that occurs weeks or months after the first sore when the bacteria spread through the body. It may cause a widespread rash, including the palms and soles, swollen glands, fever, tiredness, mouth patches, or moist lesions and requires treatment.

\medterm{Secondary Systemic Amyloidosis} Secondary systemic amyloidosis results when chronic inflammation or infection drives production of serum amyloid A, which deposits in organs such as the kidneys, liver, and spleen.

\medterm{Secondary Tumor} A tumor formed when cancer cells spread from the original site to another part of the body.

\medterm{Secretin} A hormone released by the first part of the small intestine when acidic stomach contents enter it. Secretin prompts the pancreas and bile ducts to release alkaline fluid, helping neutralize acid and support digestion.

\medterm{Secretin Stimulation Test} A test in which secretin is given and changes in blood or digestive fluid are measured. It may assess pancreatic function or investigate a rare tumor that releases excess gastrin and stomach acid.

\medterm{Secretion} The release of a substance from cells or glands into the blood, a body passage, or outside the body. Examples include saliva, hormones, digestive juices, mucus, and substances removed into urine or bile.

\medterm{Secretory Diarrhea} Watery diarrhea caused when the intestine releases too much water and salt or absorbs too little. It may continue even without eating and can result from infection, medicines, bile-acid problems, hormone-producing tumors, or inflammation.

\medterm{Section} A defined part or subdivision of an organ, document, study, or procedure; in anatomy and imaging, the term identifies a specific slice or region examined.

\medterm{Section, Cesarean} Alternate index label; see \textit{Cesarean Section}.

\medterm{Sed Rate} A blood test measuring how quickly red blood cells settle, used as a nonspecific marker of inflammation.

\synonyms
Erythrocyte sedimentation rate; ESR

\medterm{Sedation} The use of medicines to make a person relaxed, drowsy, or less aware during a procedure while maintaining appropriate monitoring and breathing.

\medterm{Sedation Dentistry} The use of medicine during dental care to reduce fear, discomfort, or awareness. It may involve inhaled gas, tablets, or medicines through a vein; deeper sedation can affect breathing and requires trained staff and continuous monitoring.

\medterm{Sedation Levels} The depth of medicine-induced drowsiness, ranging from minimal relaxation with normal responses to deep sedation with limited response and possible breathing or airway support. Sedation can deepen unexpectedly, so appropriate monitoring is required.

\medterm{Sedative-Hypnotic} A medicine or substance that slows brain activity to reduce anxiety or cause sleep. Examples include benzodiazepines, barbiturates, and some sleep medicines. Too much can cause confusion, coma, or dangerously slow breathing, especially with alcohol or opioids.

\medterm{Sedative-Hypnotic Toxidrome} A pattern of poisoning caused by substances that slow brain activity. It may cause extreme sleepiness, slurred speech, poor coordination, slow breathing, coma, or loss of airway protection, especially when combined with alcohol or opioids.

\medterm{Sedatives} Medicines or substances that slow brain activity and produce calmness, sleepiness, or reduced anxiety. Excessive amounts can impair coordination, lower blood pressure, or dangerously slow breathing, especially with alcohol or opioids.
\medterm{Sedimentation Rate} The erythrocyte sedimentation rate, a nonspecific blood test measuring how quickly red cells settle in a tube and sometimes supporting assessment of inflammation when interpreted with other findings.
\medterm{Segmental Pressures} Blood-pressure readings taken at several points on the arms and legs to compare circulation. A large difference between levels can help locate a narrowed or blocked artery in peripheral artery disease.

\medterm{Segmentation} Repeated squeezing of short sections of the intestine that mixes food with digestive fluids and brings it against the intestinal lining for absorption. Unlike peristalsis, it mainly mixes contents rather than pushing them forward.

\medterm{Segmentectomy} Surgery removing one or more defined sections of a lung rather than an entire lobe. It may treat a small, localized lung cancer or another focal problem while preserving more lung tissue; bleeding, air leakage, infection, and reduced breathing capacity are possible.

\synonyms
Sublobar Resection; Anatomical Segmentectomy; Pulmonary Segmentectomy

\medterm{SEID} Alternate index label; see \textit{Myalgic encephalomyelitis/chronic fatigue syndrome (ME/CFS)}.

\medterm{Seizure} A temporary episode caused by abnormal electrical activity in the brain. It may cause shaking, stiffening, unusual sensations, staring, confusion, or loss of awareness and can be triggered by fever, illness, injury, medicines, or epilepsy.

\medterm{Seizure Disorder} Alternate index label; see \textit{Epilepsy}.

\medterm{Seizure Disorders} Conditions characterized by recurrent unprovoked seizures or an enduring predisposition to them, classified by seizure type and cause and treated according to epilepsy syndrome and risk.

\medterm{Seizure Febrile (Febrile Seizures)} Alternate index label for Febrile Seizures.

\medterm{Seizure Fever Induced (Febrile Seizures)} Alternate index label for Febrile Seizures.

\medterm{Seizure, Absence} Alternate index label; see \textit{Absence seizure}.

\medterm{Seizure, Febrile} Alternate index label; see \textit{Febrile seizure}.

\medterm{Seizure, Grand Mal} Alternate index label; see \textit{Tonic-clonic (grand mal) seizure}.

\medterm{Seizure, Petit Mal} Alternate index label; see \textit{Absence seizure}.

\medterm{Seizure, Temporal Lobe} Alternate index label; see \textit{Temporal lobe seizure}.

\medterm{Seizures} Seizures are transient episodes of abnormal excessive brain electrical activity causing changes in awareness, movement, sensation, behavior, or autonomic function.


\medterm{Selective Deficiency Of IgA} A usually inherited immune problem in which the antibody IgA is very low or absent while other antibodies are near normal. Some people have repeated sinus, chest, or bowel infections, allergies, or autoimmune disease; others have no symptoms.

\medterm{Selective Estrogen Receptor Modulator} A medicine that acts like estrogen in some tissues and blocks estrogen effects in others. Different medicines in this group are used for conditions such as osteoporosis, breast-cancer risk reduction, or breast cancer.

\medterm{Selective Estrogen Receptor Modulators} Medicines that act like estrogen in some tissues and block estrogen effects in others, depending on the organ.

\synonyms
SERMs

\medterm{Selective IgA Deficiency} An immune condition with very low or absent immunoglobulin A, an antibody protecting mucus-lined surfaces. Some people have repeated respiratory or intestinal infections, allergies, autoimmune disease, or no symptoms.

\medterm{Selective Antibody Deficiency With Normal Immunoglobulins} An immune problem in which the total amount of antibodies is normal but the body does not make enough protective antibodies against some germs after vaccination. Repeated ear, sinus, or chest infections may occur.

\medterm{Selective Mutism} A childhood anxiety disorder in which a person speaks comfortably in some settings but repeatedly cannot speak in others, such as school or unfamiliar social situations. It lasts at least one month and can interfere with learning, relationships, or daily life.


\medterm{Selective Serotonin Reuptake Inhibitors} Antidepressants that increase serotonin signaling by slowing its reabsorption by nerve cells. They treat depression and several anxiety disorders, but effects take time and may include nausea, sexual problems, sleep changes, or withdrawal symptoms.

\synonyms
SSRIs

\medterm{Selective Serotonin-Reuptake Inhibitors (SSRIs)} Antidepressants that increase the availability of serotonin, a chemical used for communication between nerve cells. They treat depression and anxiety disorders; nausea, sleep changes, sexual problems, and withdrawal symptoms can occur.
\medterm{Selegiline} A medicine that blocks monoamine oxidase B, used for Parkinson disease and sometimes depression; interactions require caution.
\medterm{Selenium} A mineral needed for antioxidant enzymes, thyroid function, reproduction, and immune activity. Deficiency is uncommon in many regions, while excessive supplements can cause hair or nail changes, stomach upset, nerve problems, or toxicity.

\medterm{Selenium In Diet} Selenium in diet provides a trace mineral needed for antioxidant enzymes, thyroid hormone metabolism, and reproduction. Excess supplements can cause toxicity, so intake should remain within recommended limits.

\medterm{Selenium Sulphide} A skin-applied medicine that reduces certain fungi and excess scaling. It is used for dandruff and seborrheic dermatitis and may cause irritation, dryness, discoloration, or an unpleasant smell if misused.
\medterm{Selenium-75} A radioactive form of selenium used mainly as a calibration or testing source in imaging and measurement equipment. It releases radiation as it decays and has a half-life of about 120 days.

\medterm{Self Catheterization - Female} Intermittent insertion and removal of a catheter through the female urethra to empty the bladder when normal urination is inadequate, reducing retention and protecting kidney function.

\medterm{Self Catheterization - Male} Intermittent insertion and removal of a catheter through the male urethra to drain the bladder when neurologic, obstructive, or other conditions prevent adequate emptying.

\medterm{Self Gratification Sexual (Masturbation)} Alternate index label for Masturbation.

\medterm{Self Harming} Alternate index label; see \textit{Self-injury/cutting}.

\medterm{Self-Diagnosing Body Dysmorphic Disorder} Alternate index label for BDD.

\medterm{Self-Efficacy} A person's confidence in their ability to manage a health condition, follow a treatment plan, and handle symptoms. It can influence participation in care and health outcomes.

\medterm{Self-Harm} Deliberate injury or poisoning of oneself. It may be related to overwhelming emotions, mental illness, trauma, or suicidal thoughts, and can include cutting, burning, hitting, or taking too much medicine. Every episode deserves compassionate safety assessment.

\medterm{Self-Help Group} A peer group in which members support one another while changing a harmful behavior or coping with a problem.

\synonyms
Support group

\medterm{Self-Injury/Cutting} Deliberate injury to one's own body, often used to cope with overwhelming emotions and requiring compassionate assessment of safety and mental health.

\synonyms
Self Harming

\medterm{Self-Management} Actions a person takes to monitor and manage a health condition, including treatment use, symptom monitoring, and lifestyle measures.

\medterm{Self-Poisoning} Intentional ingestion or exposure to a potentially harmful substance, often requiring urgent medical and mental-health assessment.
\medterm{Self-Testing For COVID-19} Testing a sample collected at home, usually from the nose, to check for COVID-19 infection. A negative result does not always rule out infection, especially early in illness.

\medterm{Sella Turcica} A saddle-shaped hollow in a bone at the base of the skull that holds the pituitary gland. The pituitary controls many hormones, so changes in this area can affect growth, reproduction, metabolism, and the body's response to stress.
\medterm{Semaglutide} Semaglutide is a glucagon-like peptide-1 receptor agonist used for selected patients with type 2 diabetes or chronic weight management. It slows gastric emptying and reduces appetite and glucose levels; nausea, gallbladder disease, pancreatitis risk, and contraindications require review.

\medterm{Semen} The fluid released during ejaculation, containing sperm from the testes mixed with fluids from the seminal vesicles, prostate, and other glands. It carries and supports sperm; changes in amount, color, blood, or pain may need assessment.

\medterm{Semen Analysis} A laboratory examination measuring semen volume and sperm number, movement, shape, and vitality. It is a common first test when assessing fertility; results vary, so an abnormal result is often repeated and interpreted with the couple’s history.

\medterm{Semen Blood (Blood In Semen)} Alternate index label for Blood In Semen.

\medterm{Semen DNA Analysis} A test that examines genetic material inside sperm, usually for selected infertility or repeated-pregnancy-loss evaluations. It may detect increased DNA damage, but its usefulness and treatment implications depend on the clinical situation.

\medterm{Semi-Fowler Position} A semiupright position in which the head and upper body are elevated about 30 to 45 degrees.

\medterm{Semicircular Canals} Three curved, fluid-filled tubes in each inner ear that sense head rotation and help maintain balance. Problems affecting them can cause spinning dizziness, nausea, or difficulty keeping the eyes steady when the head moves.

\medterm{Semilunar Cartilages} The two crescent-shaped pads of tough cartilage inside the knee, commonly called the menisci. They cushion the joint, spread weight, and improve stability; a tear can cause pain, swelling, catching, or locking.
\medterm{Semilunar Valves} The aortic and pulmonary valves, which open during ventricular contraction and close during ventricular relaxation to prevent backflow into the ventricles.

\medterm{Seminal Vesicle} One of a paired set of male accessory reproductive glands behind the urinary bladder.
Its fructose-rich secretion contributes substantially to semen and joins the vas
deferens through the ejaculatory duct.

\medterm{Seminal Vesicles} Paired glands behind the bladder that produce much of the fluid in semen. Their secretions contain nutrients and other substances that support sperm movement and help form the ejaculate.

\medterm{Seminiferous Tubules} Tiny, coiled tubes inside the testes where sperm are made. Supporting cells nourish developing sperm and help create the environment needed for sperm production.
\medterm{Seminoma} A cancerous tumor arising from sperm-forming cells in the testis. It often grows slowly and is highly treatable with surgery, radiation, chemotherapy, or careful monitoring depending on its stage.
\medterm{Senescence} Changes and gradual loss of function associated with aging. In cells, senescence is a lasting state in which a cell stops dividing but remains active; an accumulation of such cells may contribute to aging and disease.

\synonyms
Biological aging

\medterm{Senile Lentigines} Flat, sharply defined brown or tan spots caused by chronic ultraviolet exposure, commonly appearing on sun-exposed skin with age.

\medterm{Senility} Senility is an outdated, imprecise term for cognitive decline in older age; modern assessment identifies specific causes such as dementia, delirium, depression, or medical illness.
\medterm{Senna} A plant-derived stimulant laxative containing sennosides, used short term for constipation and sometimes to empty the bowel before a procedure. It increases intestinal activity and may cause abdominal cramps or diarrhea; prolonged or excessive use can cause dependence or fluid and electrolyte disturbances.
\medterm{Sensate Focus Techniques} Structured, gradual touch exercises used in sex therapy to reduce performance pressure, improve body awareness, and rebuild comfortable intimacy. Partners agree on boundaries and initially avoid goals such as intercourse or orgasm.
\synonyms
Sensate focus

\medterm{Sensation} Awareness of changes such as touch, pain, temperature, sound, movement, or body position. Special nerve endings detect these changes and send messages to the brain, where they are experienced and interpreted.
\medterm{Sense} A sense is a capacity to detect and interpret stimuli such as touch, temperature, pain, sound, light, smell, taste, position, or internal body signals.

\medterm{Sensitisation} A process in which repeated exposure makes the body’s immune, nerve, or behavioral response stronger or easier to trigger. Sensitization can contribute to allergy, chronic pain, and some learned responses.
\medterm{Sensitivity} In testing, the proportion of people who truly have a condition and receive a positive result. In other settings, it means how strongly a person, cell, or instrument responds to a stimulus.
\medterm{Sensitivity Analysis} Sensitivity analysis tests how much a study's conclusion changes when assumptions, measurements, or input values are varied.

\medterm{Sensitivity, Gluten (Nonceliac Gluten Sensitivity (Intolerance))} Alternate index label for Nonceliac Gluten Sensitivity (Intolerance).

\medterm{Sensitized} Having antibodies against human-cell markers called HLA. This can happen after pregnancy, transfusion, or transplantation and may make it harder to find a suitable organ donor or increase the risk of transplant rejection.

\medterm{Sensorimotor Polyneuropathy} Damage to many nerves causing both altered sensation and weakness, usually starting in the feet or hands. Numbness, burning pain, poor balance, reduced reflexes, and muscle wasting may occur; diabetes is one possible cause.

\medterm{Sensorineural Deafness} Sensorineural deafness is hearing loss from dysfunction of the inner ear, auditory nerve, or central pathways rather than blockage of sound conduction.

\medterm{Sensorineural Hearing Loss} Hearing loss caused by damage to the inner ear or the nerve carrying sound to the brain. It may cause difficulty understanding speech, especially in noise, and can result from aging, loud noise, infection, inherited conditions, or medicines.

\medterm{Sensorium} A person’s level of consciousness, awareness, attention, and understanding of surroundings. Changes may include confusion, drowsiness, agitation, disorientation, or reduced responsiveness and can signal illness, medicines, intoxication, or brain injury.

\medterm{Sensory} Relating to sensation or to nerves and pathways that carry information such as touch, pain, temperature, position, hearing, or vision to the brain.
\medterm{Sensory Cortex} Areas of the brain that receive and interpret information about touch, pain, temperature, pressure, and body position. Damage can cause numbness, altered sensation, or difficulty recognizing where a body part is.
\medterm{Sensory Deprivation} A major reduction in normal input from sight, sound, touch, movement, or other senses. Prolonged deprivation can affect mood, thinking, sleep, orientation, and perception, especially in isolated or vulnerable people.
\medterm{Sensory Level} A clear line on the body below which feeling is reduced or changed. It can suggest a problem in the spinal cord, especially when accompanied by weakness, altered reflexes, or bladder and bowel problems, and needs prompt assessment.

\medterm{Sensory Nerves} Nerves that carry information such as touch, pain, temperature, and position from the body to the spinal cord and brain. Damage may cause numbness, tingling, burning pain, or reduced awareness of injury.

\medterm{Sensory Neuron} A nerve cell that carries information from receptors toward the brain and spinal cord. It helps detect touch, pain, temperature, pressure, body position, sound, light, smell, and taste.

\medterm{Sensory Processing Disorder} Difficulty noticing, interpreting, or responding to sound, touch, movement, light, or body position. It may affect behavior, learning, coordination, or daily activities and can occur with other developmental conditions; assessment and support are individualized.

\medterm{Sensory Receptors} Specialized cells or nerve endings that detect stimuli such as touch, temperature, pain, light, sound, or chemical changes.

\medterm{Sentinel Event} An unexpected healthcare event that causes death or serious physical or psychological harm, or creates a major risk of such harm. It requires prompt investigation and system changes to reduce the chance of recurrence.
\medterm{Sentinel Injury} An unexpected injury in an infant who cannot yet move independently, such as bruising, a mouth injury, or a fracture. It may be an early sign of serious harm and requires careful examination, explanation, safeguarding assessment, and documentation.

\medterm{Sentinel Lymph Node} The first lymph node or group of nodes most likely to receive fluid draining from a tumor. Removing and examining them can show whether cancer has begun to spread and help determine its stage.

\medterm{Sentinel Lymph Node Biopsy For Melanoma} Sentinel lymph-node biopsy identifies and removes the first draining lymph node or nodes from a melanoma site to detect microscopic spread. The result helps stage disease and guide further treatment; it does not remove all regional nodes routinely.

\medterm{Sentinel Lymph-Node Biopsy} Surgical removal and pathologic examination of the first draining lymph node or nodes from a primary tumor. It is used for selected melanomas, breast cancers, and other tumors to assess regional metastatic risk and guide staging.

\medterm{Sentinel Node Biopsy} Alternate index label for sentinel lymph-node biopsy: removal and examination of the first draining node or nodes to which a tumor is likely to spread, used to assess regional metastasis.

\medterm{Sentinel Surveillance} Ongoing collection of health information from selected clinics, laboratories, or communities to track disease trends. It can provide early warning of outbreaks but does not count every case in the population.

\medterm{Separation Anxiety} Separation anxiety is developmentally excessive fear or distress when apart from an attachment figure, causing avoidance, physical complaints, sleep problems, or impairment. Treatment may include gradual exposure, family-based support, cognitive behavioral therapy, and selected medication.

\medterm{Separation Anxiety Disorder} Excessive and persistent fear about being apart from a parent, partner, or other attachment figure. It may cause repeated reassurance-seeking, refusal to leave home, nightmares, physical symptoms, or difficulty attending school or work.

\medterm{Separation Anxiety In Children} Developmentally excessive distress when separated from an attachment figure, with worry, refusal, physical symptoms, nightmares, or school avoidance that impairs daily life.

\medterm{Separation Of The Iris} A tear that pulls the colored part of the eye away from its normal attachment, usually after a blunt injury. It may cause double vision, glare, an irregular pupil, bleeding, or increased eye pressure and needs eye-specialist assessment.


\medterm{Sepsis} A life-threatening condition in which the body's response to an infection damages its own organs. It may cause confusion, very fast breathing or heartbeat, low blood pressure, reduced urine, or severe weakness and needs emergency treatment.

\synonyms
Septicemia (Sepsis)

\medterm{Sepsis Syndrome} A range of severe responses to infection, from infection-related organ dysfunction to septic shock and failure of several organs. The condition can worsen quickly and requires urgent assessment and treatment.

\medterm{Septal Defect} An abnormal opening in the septum between the heart's chambers that permits blood to pass between the right and left sides of the heart.
\medterm{Septal Defect, Atrioventricular} Alternate index label; see \textit{Atrioventricular canal defect}.

\medterm{Septal Defect, Ventricular} Alternate index label; see \textit{Ventricular septal defect}.

\medterm{Septal Hematoma} A collection of blood inside the wall dividing the two sides of the nose, usually after an injury. It causes blockage and painful swelling and needs urgent drainage because untreated blood can destroy cartilage and deform the nose.

\medterm{Septate Uterus} A congenital difference in which a band of tissue partly or completely divides the inside of the uterus. Many people have no symptoms, but it can increase miscarriage or some pregnancy risks; imaging confirms it, and selected cases are treated by removing the dividing tissue.

\synonyms
Uterine Septum

\medterm{Septic} Relating to a serious infection or to harmful germs in a normally sterile body site. For example, septic arthritis is an infected joint, while septic shock is life-threatening organ and circulation failure caused by an overwhelming infection.

\medterm{Septic Abortion} A miscarriage or induced pregnancy loss complicated by infection of the uterus or surrounding tissues. Fever, chills, worsening abdominal pain, foul discharge, fast heartbeat, or faintness may occur; it is an emergency requiring antibiotics, supportive care, and removal of infected retained tissue when necessary.

\synonyms
Septic Miscarriage; Infected Abortion

\medterm{Septic Arthritis} Infection of a joint, usually bacterial, causing acute pain, swelling, restricted movement, and systemic illness.

\synonyms
Arthritis Infectious (Septic Arthritis)

\medterm{Septic Embolism} A piece of infected material travels through the bloodstream and blocks a smaller blood vessel, spreading infection to another organ. It often starts with an infected heart valve and can cause abscesses or tissue damage.

\medterm{Septic Pulmonary Infarction} Death of part of the lung after infected clots or other infected material block small lung arteries. It may cause fever, chest pain worse with breathing, cough, coughing blood, and lung abscesses, and requires urgent treatment of the infection.

\medterm{Septic Shock} A life-threatening complication of sepsis in which infection causes dangerously low blood pressure and poor blood flow to organs, even after initial fluids. It requires immediate hospital treatment with antibiotics, fluids, and medicines to support circulation.
\medterm{Septicaemia} An older term for bloodstream infection, often used to describe bacteremia with systemic illness.
\medterm{Septicemia} Septicemia is an older term for serious bloodstream infection; modern assessment distinguishes bacteremia from sepsis, which is life-threatening organ dysfunction caused by infection.

\medterm{Septicemia (Sepsis)} Alternate index label for Sepsis.

\medterm{Septo-Optic Dysplasia} A condition present from birth involving underdevelopment of one or both optic nerves, parts of the brain’s midline, and sometimes the pituitary gland. It may cause poor vision, involuntary eye movements, seizures, developmental differences, or hormone deficiencies.

\medterm{Septoplasty} Septoplasty is surgery to straighten or reshape the nasal septum to improve airflow or address structural obstruction.


\medterm{Septum} A wall or partition separating two spaces or tissues, such as the nasal septum or the septum between heart chambers.
\medterm{Septum Deviated (Deviated Septum)} Alternate index label for Deviated Septum.

\medterm{Septum, Heart} The wall separating the right and left chambers of the heart, consisting of atrial and ventricular portions and preventing normal mixing of oxygenated and deoxygenated blood.

\medterm{Sequela} A lasting consequence or condition that remains after an earlier disease, injury, infection, or treatment has ended, such as paralysis after stroke or scarring after infection.

\medterm{Sequelae} Lasting effects or complications that remain after an illness or injury. Examples include weakness after a stroke, breathing problems after severe infection, or pain after an injury; sequelae may be physical, mental, emotional, or social.
\medterm{Sequential Therapy} Treatment in which different medicines or interventions are given one after another rather than at the same time. The sequence may improve safety, effectiveness, or tolerance.

\medterm{Sequester} To sequester means to isolate or set apart; in medicine, sequestered tissue or material is separated from its normal circulation or surroundings.

\medterm{Sequestrum} A piece of dead bone that has separated from healthy bone, usually because of a long-lasting bone infection. It can prevent healing, cause drainage or pain, and may need surgery after infection is treated.
\medterm{Ser} Ser is the standard abbreviation for serine, an amino acid incorporated into proteins and involved in metabolism and cell signaling.


\medterm{Serine} Serine is a nonessential amino acid used in protein synthesis, one-carbon metabolism, membrane formation, and signaling.

\medterm{Serious Diseases And Health Problems (Symptoms Of Serious Diseases And Health Problems)} Alternate index label for Symptoms of Serious Diseases and Health Problems.

\medterm{Serious Adverse Event} An unwanted medical event that results in death, is life-threatening, requires or prolongs hospitalization, causes major disability, or causes a birth defect. It may occur during treatment without being caused by the treatment.

\medterm{SERMs} Selective estrogen receptor modulators that have estrogen-like effects in some tissues and blocking effects in others.

\synonyms
Selective estrogen receptor modulators

\medterm{Seroconversion} Development of detectable antibodies in blood after infection, vaccination, or another immune-triggering event, showing a changed immune status.
\medterm{Serodiagnostic} Relating to diagnosis using antibodies, antigens, or other markers measured in blood serum. These tests can show infection, immunity, inflammation, or other conditions but must be interpreted with timing and clinical findings.

\medterm{Serologic Test} A blood test that detects antibodies, antigens, or other immune markers to help diagnose infection, assess immunity, or monitor response to treatment.

\medterm{Serology} The laboratory evaluation of serum for antibodies, antigens, or other immune markers.
Results may indicate current infection, previous exposure, vaccination, autoimmune
reactivity, or passive antibody transfer and must be interpreted with timing and
clinical context.

\medterm{Serology For Brucellosis} Blood testing for antibodies to Brucella species, interpreted with exposure history and symptoms because cross-reactions and early false-negative results can occur.

\medterm{Seroma} A pocket of clear body fluid that forms near a surgical incision or injured tissue. It may cause swelling, pressure, or discomfort and often resolves, though large or painful collections may need drainage.

\medterm{Seronegative Arthritis} Inflammatory arthritis in which common blood tests for rheumatoid factor and anti-CCP antibodies are negative. It includes several diseases with different patterns, such as spine, skin, bowel, or eye involvement, so symptoms and examination guide diagnosis.

\medterm{Serosanguineous} Serosanguineous describes thin wound or body fluid containing both clear serum and a small amount of blood, often appearing pink or pale red.

\medterm{Serositis} Inflammation of the smooth lining around the lungs, heart, or abdominal organs. It can cause sharp pain that worsens with breathing or movement and may cause fluid to collect around the affected organ.

\medterm{Serotonin} A chemical messenger used by nerve cells and other tissues. It helps regulate mood, sleep, appetite, bowel function, and pain; too much activity from medicines can cause serotonin syndrome.

\synonyms
5-hydroxytryptamine; 5-HT

\medterm{Serotonin Blood Test} A blood test measuring serotonin, used selectively in evaluating suspected serotonin-producing neuroendocrine tumors; many medicines and specimen factors affect results.

\medterm{Serotonin Syndrome} A potentially life-threatening reaction caused by too much serotonin activity, usually after combining or taking too much of certain medicines. Agitation, sweating, diarrhea, shaking, fever, muscle stiffness, or confusion may occur.

\medterm{Serotonin-Norepinephrine Reuptake Inhibitors} Antidepressants that increase the activity of serotonin and norepinephrine, chemical messengers between nerve cells. They treat depression, some anxiety disorders, and certain nerve or long-lasting pain conditions; nausea, sweating, sleep changes, and withdrawal symptoms can occur.

\medterm{Serotype} A distinct version of a microorganism or antigen identified by how antibodies react to it.
\medterm{Serous} Thin, watery, and usually clear, resembling blood serum. Serous fluid can occur normally or collect in a wound, body cavity, or blister.
\medterm{Serous Borderline Ovarian Tumor} An ovarian tumor whose cells look abnormal but have not invaded nearby supporting tissue. It may stay in the ovary or form deposits on its surface; treatment and follow-up depend on its extent and microscopic features.

\medterm{Serous Cyst Adenomas Pancreas (Pancreatic Cysts)} Alternate index label for Pancreatic Cysts.

\medterm{Serous Fluid Cytology} Examination of cells in fluid collected from around the lungs, heart, or abdominal organs. It can look for cancer cells, infection, or inflammation; additional tests may be needed when the result is unclear.
\medterm{Serous Membranes} Thin, smooth linings that cover organs and reduce friction inside closed spaces, including the linings around the lungs, heart, and abdominal organs.
\medterm{Serous Tubal Intraepithelial Carcinoma} An early, noninvasive cancerous change in cells at the end of a fallopian tube near the ovary. It is thought to be the starting point for many high-grade cancers of the ovary, tube, or abdominal lining.

\medterm{Serpentine Supravenous Hyperpigmentation} A pattern of dark, wavy skin discoloration that follows the paths of surface veins, usually on the upper chest or neck. It can occur with vein inflammation, poor blood flow, some medicines, or cancer and requires assessment of the underlying cause.

\medterm{Serrated Pathway} One route by which some colon cancers develop from flat or saw-toothed polyps. These polyps may be easy to miss during colonoscopy and can acquire cell changes that increase cancer risk, so removal and follow-up are important.

\medterm{Serrated Polyposis Syndrome} A disorder characterized by multiple serrated lesions or polyps in the colon and an increased risk of colorectal cancer, requiring colonoscopic surveillance and removal of significant lesions.

\medterm{Serration} A row of small tooth-like projections or notches along an edge, used to describe anatomy, cells, instruments, or imaging appearances.

\medterm{Serratus Anterior Muscle} A muscle along the side of the chest that holds the shoulder blade against the ribs and helps raise the arm. Injury to its nerve can make the shoulder blade stick out and weaken pushing or overhead movements.

\medterm{Sertoli-Cell Tumor} A tumor arising from Sertoli cells, which support sperm development in the testicle. It is often benign but may produce estrogen or other hormones and cause breast enlargement, infertility, or a testicular mass.
\medterm{Sertoli-Leydig Cell Tumor} A rare ovarian tumor that may produce male hormones, causing increased facial hair, acne, a deeper voice, or irregular periods. Treatment depends on the tumor's stage and features.

\medterm{Serum} The clear liquid left after blood has clotted. It contains water, salts, hormones, antibodies, nutrients, and waste products but lacks fibrinogen and most clotting proteins, making it useful for laboratory tests.
\medterm{Serum Concentration} The amount of a substance measured in the liquid part of clotted blood. Results may reflect when the sample was taken, how the substance is absorbed and removed, and whether the level is expected to help or harm.

\medterm{Serum Free Hemoglobin Test} A blood test measuring hemoglobin not contained within red cells, used to support evaluation of intravascular hemolysis, transfusion reactions, or mechanical red-cell destruction.

\medterm{Serum Glutamic Oxaloacetic Transaminase} SGOT, now called AST, is an enzyme released by injured liver, heart, skeletal muscle, or other cells; its level must be interpreted with other findings.

\medterm{Serum Glutamic Pyruvic Transaminase} SGPT, now called ALT, is an enzyme concentrated in liver cells and commonly measured as a marker of hepatocellular injury.

\medterm{Serum Herpes Simplex Antibodies} A blood test for antibodies against herpes simplex virus types 1 or 2. IgG usually indicates previous infection, while IgM is less reliable; results do not show the infection site or prove that a current lesion is caused by herpes.

\medterm{Serum Iron Test} A blood test measuring circulating iron bound mainly to transferrin; it is interpreted with ferritin, transferrin saturation, and blood counts rather than used alone to diagnose iron deficiency.

\medterm{Serum Osmolal Gap} The difference between the measured concentration of dissolved particles in blood and the calculated value. A high gap may suggest an unmeasured substance, including a toxic alcohol, but other illnesses and laboratory factors can also affect it.

\medterm{Serum Phenylalanine Screening} A blood test measuring phenylalanine, used in newborn screening and follow-up of phenylketonuria or other disorders of phenylalanine metabolism.

\medterm{Serum Progesterone} A blood test measuring progesterone, a hormone made mainly after ovulation and during pregnancy. The result can help assess ovulation, ovarian function, fertility treatment, or early pregnancy, but interpretation depends strongly on timing and pregnancy status.

\medterm{Serum Protein Electrophoresis} A blood test that separates proteins in serum into groups. It can show abnormal antibody production or low protein levels and is often followed by immunofixation when a plasma-cell disorder is suspected.

\medterm{Serum Sickness} A delayed immune reaction, usually one to three weeks after exposure to certain medicines, injected proteins, or antiserum. It can cause fever, rash or hives, painful swollen joints, enlarged lymph nodes, and sometimes organ problems.
\medterm{Serum Therapy} Prevention or treatment using antibody-containing plasma, immunoglobulin, antitoxin, or antivenom. It provides immediate but usually temporary protection against a specific infection or toxin; allergic reactions and other infusion-related problems are possible.
\medterm{Serving Size} Serving size is the standardized amount used to describe the nutrients or calories in food and may differ from the amount a person actually eats.

\medterm{Sesamoid Bones} Small bones embedded within or beside tendons, where they protect the tendon and improve leverage. Examples include the kneecap and the small bones beneath the big toe.
\medterm{Sesamoiditis} Painful inflammation of tissues around a sesamoid bone, most often beneath the big toe. Pain usually worsens with walking, running, dancing, or pressure on the front of the foot.

\medterm{Sesquipedalian} Sesquipedalian means characterized by very long words; clear healthcare communication favors language that patients can understand.

\medterm{Sessile} Attached directly by a broad base rather than a narrow stalk. A sessile polyp can be harder to remove completely and may require careful examination because its shape can hide abnormal tissue.
\medterm{Set} In exercise, a group of consecutive repetitions performed before a rest period. In medicine, “set” can also mean a prepared group of instruments, tests, or equipment.

\medterm{Sever Disease} A painful irritation of the growth area at the back of the heel in growing children, often during sports. Heel pain worsens with running or jumping and usually improves with activity adjustment, heel support, calf stretching, and time.
\medterm{Severe Acute Respiratory Syndrome} A serious contagious respiratory illness caused by the SARS-associated coronavirus, SARS-CoV. It usually begins with fever and may cause chills, muscle aches, headache, diarrhea, dry cough, and breathing difficulty. Most affected people develop pneumonia, which can progress to low blood oxygen, respiratory failure, and death. The name usually refers to the disease identified in the 2002--2004 outbreak; SARS-CoV-2 causes COVID-19, a different disease.

\medterm{Severe Obesity} Obesity with a Body Mass Index of 40 kg/m\(^2\) or higher, or 35 kg/m\(^2\) or higher together with a serious obesity-related health problem. Body Mass Index does not fully measure health.

\synonyms
Severe Acute Respiratory Syndrome (Severe Acute Respiratory Syndrome (SARS))

\medterm{Severe Acute Respiratory Syndrome (Severe Acute Respiratory Syndrome)} Alternate index label for Severe Acute Respiratory Syndrome.

\medterm{Severe Brain Injury, Coma} Alternate index label; see \textit{Coma}.

\medterm{Severe Combined Immunodeficiency} A group of inherited conditions in which important immune cells do not develop or work properly. Affected babies can develop severe, repeated infections early in life and need urgent specialist care, protective measures, and immune-restoring treatment.

\medterm{Severe Condition (Sever Condition)} Alternate index label for Sever Condition.

\medterm{Severe Dengue} The modern World Health Organization (WHO) classification term for the critical, life-threatening form of dengue virus infection, typically developing around the time of fever defervescence (days 3 to 7). It is characterized by one or more of three major clinical manifestations: severe plasma leakage resulting in hypovolemic shock (dengue shock syndrome) or respiratory distress from fluid accumulation, severe bleeding (such as massive gastrointestinal hemorrhage), or severe organ impairment including acute liver failure, encephalitis, or myocarditis.

\synonyms
Dengue Shock Syndrome; Severe Dengue Virus Infection; Critical Dengue

\medterm{Sevoflurane} An inhaled anesthetic gas used to start and maintain general anesthesia. It acts in the brain to produce unconsciousness and reduces awareness of pain; it can lower breathing and blood pressure and may trigger malignant hyperthermia in susceptible people.

\medterm{Sex Change} A nontechnical term for gender transition or medical gender-affirming care; respectful, current terminology should be used.
\medterm{Sex Chromosomes} The X and Y chromosomes, which contribute to sex development and carry many other genes.
\medterm{Sex Education} Teaching about bodies, relationships, consent, contraception, sexually transmitted infections, and sexual health. Accurate, age-appropriate education supports safer decisions, respectful relationships, prevention, and access to appropriate healthcare.
\medterm{Sex Headaches} Headaches brought on by sexual activity, including a gradually increasing ache or a sudden severe headache at orgasm. A first or thunderclap headache requires urgent assessment to exclude bleeding or other dangerous causes.

\medterm{Sex Hormones} Hormones involved in sexual development, reproduction, menstrual cycles, pregnancy, and related body functions. Major groups include estrogens, progesterone, and androgens, which also affect bones, muscles, mood, and tissues.
\medterm{Sex Therapy} Structured psychological or behavioral therapy for sexual concerns or dysfunction.
\medterm{Sex-Linked Dominant} A pattern of inheritance in which a disease-causing gene variant lies on a sex chromosome and one altered copy can cause the condition. The chance of passing it on differs for mothers and fathers and for children of different sexes.

\medterm{Sex-Linked Inheritance} Inheritance of a trait caused by a gene on a sex chromosome, commonly the X chromosome.
\medterm{Sex-Linked Recessive} A pattern in which a disease-causing gene variant, usually on the X chromosome, commonly causes disease in males with one altered copy. Females usually need two altered copies, while those with one may be carriers or have milder symptoms.

\medterm{Sexual Abuse} Nonconsensual sexual contact, exploitation, or behavior; it is a safeguarding and health emergency when immediate danger or injury exists.
\medterm{Sexual Assault} Any sexual contact or sexual act without freely given consent, including unwanted touching, penetration, coercion, or sexual activity when a person cannot consent. It can cause injury, infection, pregnancy, and psychological trauma and may require medical and safeguarding support.

\synonyms
Rape (Sexual Assault)

\medterm{Sexual Assault - Prevention} Sexual-assault prevention includes consent education, recognizing coercion, safer environments, bystander intervention, and access to confidential support; responsibility always lies with the person who commits the assault.

\medterm{Sexual Dysfunction} A group of conditions involving persistent or recurrent difficulty with sexual desire, arousal, erection, ejaculation, orgasm, or pain that causes clinically significant distress. Causes may be physical, psychological, medication-related, hormonal, or related to life circumstances.
\medterm{Sexual Dysfunction, Female} Alternate index label; see \textit{Female sexual dysfunction}.
\medterm{Sexual Dysfunction Disorders} Alternate index label for Sexual Dysfunction.

\medterm{Sexual Intercourse (Coitus)} Sexual activity involving penetration, most commonly penile-vaginal, but sometimes penile-anal or other forms depending on the people involved. It may be for pleasure, intimacy, or reproduction and requires freely given consent.

\synonyms
Coitus

\medterm{Sexual Masochism Disorder} A mental-health diagnosis involving repeated sexual arousal from being humiliated, beaten, bound, or otherwise made to suffer, when the pattern causes significant distress or disrupts daily life, or involves a person who has not freely consented.

\medterm{Sexual Obsession} Alternate index label; see \textit{Compulsive sexual behavior}.

\medterm{Sexual Rehabilitation} Assessment and treatment to help restore or adapt sexual function after illness, injury, surgery, disability, or medicine effects. It may include education, counseling, pelvic-floor therapy, devices, or medicines.

\medterm{Sexual Sadism} A recurrent sexual interest in another person’s physical or psychological suffering or humiliation. The interest alone is not a mental disorder; acts without consent are abuse.

\medterm{Sexual Sadism Disorder} A mental-health diagnosis involving repeated sexual arousal from another person’s physical or emotional suffering, when it causes clinically significant distress or impairment or involves acting with someone who has not freely consented.

\medterm{Sexual Violence} Sexual violence is any sexual act, attempt, contact, or coercion without freely given consent and can cause physical injury, infection, pregnancy, trauma, and long-term health effects.

\medterm{Sexually Transmitted Disease Treatments} Alternate index label for Sexually Transmitted Infections.

\medterm{Sexually Transmitted Diseases} Alternate index label for Sexually Transmitted Infections.

\medterm{Sexually Transmitted Diseases (STDs In Women)} Alternate index label for Sexually Transmitted Infections.
\medterm{Sexually Transmitted Diseases (STDs)} Alternate index label for Sexually Transmitted Infections.

\medterm{Sexually Transmitted Diseases And Pregnancy (STDs)} Alternate index label for Sexually Transmitted Infections.

\medterm{Sexually Transmitted Infection} Alternate index label for Sexually Transmitted Infections.

\medterm{Sexually Transmitted Infections} Infections caused by bacteria, viruses, fungi, or parasites that spread mainly through sexual contact, including vaginal, anal, or oral sex and genital skin-to-skin contact. Some can also spread through blood exposure or from a pregnant person to a baby. Many cause no symptoms, especially early in the infection.

\synonyms
STIs
STDs

\medterm{Sezary Disease} Alternate index label; see \textit{Sezary syndrome}.

\medterm{Sezary Syndrome} An uncommon form of cutaneous T-cell lymphoma involving widespread red skin and abnormal lymphocytes in the blood.

\synonyms
Sezary Disease

\medterm{SéZary Syndrome} An uncommon form of skin lymphoma in which abnormal immune cells are found in the skin and blood. It usually causes widespread red, scaly, itchy skin and enlarged lymph nodes and requires specialist care.

\medterm{Sgarbossa Criteria} A set of electrocardiogram findings used to look for a heart attack when an existing conduction problem or pacemaker makes the tracing difficult to read. The findings can support, but cannot prove, an acute blocked heart artery.

\synonyms
Modified Sgarbossa Criteria; Smith-Modified Sgarbossa Rule

\medterm{Shadowing Artifact} An ultrasound artifact in which a strongly attenuating or reflecting structure produces a dark area behind it because little sound reaches the tissues beyond. Stones, bone, and gas commonly cause shadowing.

\medterm{Shaken Baby Syndrome} An older term for abusive head trauma in an infant or young child caused by violent shaking, with or without impact. It can cause brain and eye injuries, seizures, breathing problems, vomiting, unusual sleepiness, or death and requires emergency protection and care.

\medterm{Shaken Baby Syndrome (Abusive Head Trauma)} Alternate index label for Abusive Head Trauma.

\medterm{Shaken Impact Syndrome} A term for abusive head trauma involving shaking with or without impact, especially in infants.
\medterm{Shampoo - Swallowing} Accidental swallowing of a small amount of shampoo usually causes mouth or stomach irritation; larger or concentrated exposures may cause vomiting or breathing problems and warrant poison-control advice.

\medterm{Shared Decision-Making} A collaborative process in which clinicians and patients combine best available evidence
with patient values, preferences, goals, and circumstances to choose among reasonable
options.

\medterm{Shark Attack} A shark attack can cause severe lacerations, bleeding, tissue loss, drowning, and shock; emergency rescue, hemorrhage control, and trauma care are critical.

\medterm{Sharp Force Trauma} An injury caused by a sharp or pointed object, such as a knife, broken glass, or needle. It may produce a cut or a deeper puncture and can damage blood vessels, nerves, organs, or bones; deep or heavily bleeding wounds need emergency care.


\medterm{Sharps Container} A strong, leak-resistant container for used needles, lancets, blades, and other sharp medical items. It helps prevent injury and infection and should be clearly labeled, kept upright, and closed or replaced before it becomes too full.

\medterm{Sharps Injury} A percutaneous or mucocutaneous injury caused by a needle or other sharp instrument that
may expose a person to blood or body fluids. Immediate wound care, incident reporting,
risk assessment, and pathogen-specific post-exposure management are required.

\medterm{Shave Biopsy} Removal of a raised or superficial skin lesion by shaving off a thin layer with a sterile blade. The tissue is examined under a microscope; bleeding, infection, scarring, or an inadequate sample can occur.

\medterm{Shaving Cream Poisoning} Poisoning after swallowing shaving cream or inhaling its propellant. Small accidental amounts often cause mild stomach upset, but larger exposures can cause vomiting, choking, drowsiness, or breathing problems; seek product-specific advice.

\medterm{Sheehan Syndrome} Damage to the pituitary gland after severe blood loss or very low blood pressure during childbirth. It can cause failure to produce breast milk, extreme tiredness, low blood pressure, menstrual or fertility problems, and deficiencies of several hormones.

\medterm{Sheehan Syndrome (Postpartum Pituitary Necrosis)} Permanent injury to the pituitary gland caused by severe bleeding or very low blood pressure during or after childbirth. Because the gland makes several hormones, symptoms may include inability to breastfeed, low blood pressure, fatigue, and menstrual or fertility problems.

\synonyms
Postpartum Pituitary Necrosis

\medterm{Shellac Poisoning} Poisoning from swallowing or inhaling shellac, a coating product that may contain alcohol or other solvents. It can cause irritation, vomiting, dizziness, drowsiness, or lung injury after aspiration; do not induce vomiting.

\medterm{Shellfish Allergy} An immune reaction to proteins in crustaceans or mollusks that can cause hives, digestive or respiratory symptoms, or anaphylaxis.

\synonyms
Allergy, Shellfish

\medterm{Shellfish Poisoning} Illness caused by toxins accumulated in contaminated shellfish. Depending on the toxin, symptoms may include vomiting, diarrhea, tingling, weakness, confusion, breathing difficulty, or paralysis and require urgent medical advice.
\medterm{Shielding} Material placed between a person and a source of radiation to reduce exposure. The material depends on the radiation: lead or tungsten is used for many X-rays and gamma rays, while plastic or acrylic can be useful for beta radiation.

\medterm{Shift Work Sleep Disorder} A circadian rhythm sleep-wake disorder caused by work schedules that overlap the usual sleep period, producing insomnia, excessive sleepiness, or both.

\medterm{Shigella} A genus of bacteria causing shigellosis, an acute diarrheal illness that may include fever and bloody stools.
\medterm{Shigellosis} Shigellosis is intestinal infection caused by Shigella bacteria, producing diarrhea that may be bloody, fever, cramps, and dehydration.

\medterm{Shin Bone Fever} An older nonspecific term for painful inflammation or infection involving the shin or nearby tissues.

\medterm{Shin Splints} A common term for medial tibial stress syndrome, an overuse injury causing pain and tenderness along the inner border of the shinbone. It is associated with repetitive impact or a sudden increase in activity and involves irritation of tissues around the bone. It is distinct from a stress fracture and acute compartment syndrome, which are separate causes of lower-leg pain.

\synonyms
Medial Tibial Stress Syndrome


\medterm{Shin Spot} A shin spot is a nonspecific mark or lesion on the lower leg; its significance depends on appearance, duration, symptoms, injury, and associated disease.

\medterm{Shingles} Herpes zoster, reactivation of varicella-zoster virus in sensory nerves causing a painful one-sided blistering rash. Pain, burning, or tingling may begin before the rash. Eye or ear involvement can threaten vision or hearing, and persistent nerve pain after the rash heals is called postherpetic neuralgia; prompt assessment is important for facial or eye-area rashes.
\medterm{Shingles (Herpes Zoster)} Alternate index label for Herpes Zoster.


\medterm{Shingles And Pregnancy} Shingles is reactivation of varicella-zoster virus and usually causes a painful one-sided blistering rash. It is not the same as primary chickenpox exposure, but pregnancy requires prompt clinical advice about treatment, pain, and contact with people at risk.


\medterm{Shirodkar Cerclage} An operation that places a strong stitch high around the cervix to help prevent it opening too early during pregnancy. It is used in selected people at risk of late miscarriage or premature birth and can cause bleeding, infection, or injury to nearby organs.

\medterm{Shivering} Repeated, involuntary muscle contractions that generate heat, usually in response to cold or fever. Severe or unexplained shivering may occur with infection, medicine reactions, or other illness and should be assessed with the other symptoms.

\medterm{Shock} A life-threatening failure of circulation in which organs do not receive enough blood and oxygen. It can result from severe blood or fluid loss, heart failure, overwhelming infection, a severe allergic reaction, or an obstruction to blood flow. Signs may include confusion, faintness, cold or mottled skin, rapid breathing, little urine, or very low blood pressure and require emergency treatment.
\medterm{Shock Lung} An older term for acute respiratory distress syndrome (ARDS), a life-threatening form of lung inflammation in which fluid and injury make it difficult for oxygen to enter the blood.
\medterm{Shock Therapy} An imprecise term that may refer to electroconvulsive therapy or outdated shock treatments; the specific intervention should be named.
\medterm{Shock Wave Lithotripsy} A procedure that focuses repeated sound-pressure waves from outside the body onto a kidney or ureter stone to break it into smaller pieces that can pass in urine. Blood in the urine, pain, blockage, infection, or incomplete stone removal can occur.

\medterm{Shock, Secondary} Circulatory failure occurring as a consequence of another condition, such as infection, trauma, bleeding, heart failure, allergy, or endocrine crisis, with inadequate tissue perfusion.

\medterm{Shone Complex} A group of congenital left-sided obstructive heart lesions, classically including
supravalvular mitral ring, parachute mitral valve, subaortic stenosis, and coarctation
of the aorta. The combination and severity vary and may lead to heart failure or
systemic outflow obstruction.

\medterm{Short Arm Of A Chromosome} The shorter section of a chromosome, labeled p from the French word for “small.” Missing, extra, or rearranged DNA in this section can alter genes and cause disease.

\medterm{Short Bones} Small, roughly cube-shaped bones that provide stability and limited movement. The wrist and ankle contain most of the body's short bones.

\medterm{Short Bowel Syndrome} A condition in which the small intestine is shortened or damaged and cannot absorb enough nutrients from food. It most often occurs after surgery to remove part of the small intestine and may cause diarrhea, weight loss, and nutritional deficiencies.

\medterm{Short QT Syndrome} A rare inherited heart-rhythm disorder in which the heart's electrical recovery time is unusually short. It can cause atrial or ventricular arrhythmias, fainting, or sudden cardiac death; diagnosis requires specialist ECG and family-history assessment, and treatment may include an implantable defibrillator.

\synonyms
SQTS


\medterm{Short Philtrum} A short philtrum is reduced distance between the nose and upper lip and may be an isolated feature or part of a genetic, developmental, or prenatal-exposure syndrome.

\medterm{Short Stature} Height substantially below the expected range for age and sex, requiring assessment of growth pattern and cause.
\medterm{Short-Sight} Myopia, a refractive error in which distant objects are focused in front of the retina.
\medterm{Short-Term Memory} The ability to hold and use a small amount of information for seconds or minutes. It helps with following instructions, doing mental tasks, and remembering a recent conversation before information is stored more permanently or forgotten.

\synonyms
Working memory

\medterm{Shortness Of Breath} An uncomfortable awareness of difficult or insufficient breathing. See Dyspnea.

\medterm{Shotty} Shotty describes small, firm, discrete lymph nodes, often reactive after infection; persistent, enlarging, hard, or generalized nodes require evaluation.

\medterm{Shoulder} The area where the upper limb joins the body, including the shoulder blade, collarbone, and upper-arm bone. Its main joint allows a wide range of arm movement but is less stable and can be injured or dislocated.
\medterm{Shoulder Arthroscopy} A minimally invasive procedure using a camera and instruments through small incisions to diagnose and treat selected disorders of the shoulder joint and surrounding tissues.

\medterm{Shoulder Bursitis} Shoulder bursitis is inflammation of a fluid-filled sac around the shoulder, often associated with rotator-cuff disease or repetitive overhead activity. Pain with elevation and lying on the affected side is common; activity modification, therapy, medicines, injection, or treatment of the underlying cause may help.

\synonyms
Bursitis Shoulder (Shoulder Bursitis)

\medterm{Shoulder CT Scan} Computed tomography of the shoulder used to evaluate fractures, bone alignment, arthritis, tumors, infection, and complex bony anatomy.

\medterm{Shoulder Dislocation} Complete displacement of the upper arm bone from the shoulder socket, usually caused by a fall, collision,
or forceful twisting. It causes severe pain, deformity, and inability to move the arm; urgent reduction
and assessment for associated fractures or nerve and blood-vessel injury are required.

\medterm{Shoulder Dystocia} A childbirth emergency in which the baby’s head is born but a shoulder becomes stuck behind the mother’s pelvic bone. It requires immediate maneuvers because delay can cause lack of oxygen, fractures, or nerve injury in the baby and heavy bleeding in the mother.

\medterm{Shoulder Impingement} Pain associated with movement of the shoulder, often involving the rotator-cuff tendons or
subacromial tissues. Symptoms include pain with overhead activity and sometimes weakness;
assessment and treatment are individualized and may include exercise-based rehabilitation.

\medterm{Shoulder Instability} Excessive or recurrent movement of the shoulder joint caused by injury, laxity, or altered
joint anatomy. Associated lesions may include a Bankart tear, Hill-Sachs lesion, or injury to the joint capsule or
ligaments.

\medterm{Shoulder MRI Scan} Magnetic resonance imaging of the shoulder used to evaluate rotator-cuff tendons, labrum, cartilage, ligaments, muscles, bone marrow, and joint fluid.

\medterm{Shoulder Pain} Pain in or around the shoulder caused by injury, tendon or bursa disorders, arthritis, nerve disease, or
referred pain. Sudden weakness, deformity, chest symptoms, fever, or pain after major trauma
requires prompt evaluation.

\medterm{Shoulder Replacement} Surgery replacing damaged surfaces of the shoulder with artificial parts to relieve severe pain and improve movement. The implant may replace the joint in the usual or reverse configuration; infection, stiffness, dislocation, loosening, or nerve injury can occur.




\medterm{Shulmans Syndrome (Eosinophilic Fasciitis)} Alternate index label for Eosinophilic Fasciitis.

\medterm{Shunt} A passage that diverts blood, cerebrospinal fluid, urine, or another fluid from one channel to another.
\medterm{Shunt (Vascular)} A surgically created or implanted pathway that redirects blood or another fluid from one vessel or body space to another.

\synonyms
Vascular shunt; blood-vessel shunt

\medterm{Shwachman Syndrome} An inherited multisystem disorder usually involving exocrine pancreatic insufficiency,
bone-marrow dysfunction, and skeletal abnormalities. Children may have steatorrhea, poor
growth, recurrent infections, neutropenia, anemia, or increased risk of myeloid
malignancy.

\medterm{Shy-Drager Syndrome} Alternate index label; see \textit{Multiple system atrophy}.

\medterm{SI Units} The International System of Units used for standardized measurements, including meter, kilogram, second, mole, and pascal.
\medterm{SIADH} A condition in which the body releases too much antidiuretic hormone, causing it to retain water and dilute the blood sodium level. It may cause nausea, headache, confusion, seizures, or coma and requires treatment of the cause and sodium abnormality.

\medterm{Sialadenitis} Inflammation or infection of a salivary gland, causing painful swelling near the cheek or under the jaw, often worse during meals. Causes include bacteria, viruses, dehydration, or a blocked duct; persistent fever or pus needs medical care.

\medterm{Sialagogues} Substances or medicines that increase saliva production. They may help some people with dry mouth, but treatment also considers hydration, medicines, oral care, and the cause of reduced saliva.
\medterm{Sialogram} An imaging study of a salivary gland and its ducts after contrast is introduced through the duct opening, used to assess obstruction, stones, strictures, or chronic inflammation.

\medterm{Sialography} An imaging procedure that uses contrast material to show salivary glands and their ducts. It can help identify blockage, stones, narrowing, chronic inflammation, or tumors, although newer imaging is often preferred.

\medterm{Sialolithiasis} Formation of a stone in a salivary gland or its drainage tube, most often beneath the jaw. It can cause repeated swelling and pain during meals and may lead to blockage or infection.

\medterm{Salivary Gland Stones} Stones formed from mineral deposits in a salivary gland or its drainage tube. They can block saliva flow and cause swelling or pain, especially during meals; the medical term is sialolithiasis.

\medterm{Sialorrhea} Excess saliva or difficulty controlling and swallowing it, causing drooling. It commonly occurs with neurologic disease, swallowing difficulty, dental problems, or some medicines; treatment addresses the cause and may include swallowing therapy or selected medicines.
\medterm{Sialorrhoea} Excess saliva that pools in the mouth or dribbles out as drooling. It is often caused by difficulty controlling or swallowing saliva rather than making too much. It may occur with cerebral palsy, Parkinson disease, motor-neuron disease, or other conditions and can cause skin irritation or choking. Swallowing therapy and other treatments may help.
\medterm{Sialuria} A rare inherited disorder in which cells make and release too much sialic acid, a sugar used on cell surfaces. It may cause developmental delay, enlarged liver and spleen, distinctive facial features, and variable learning difficulties.

\medterm{Sibling} A person's brother or sister; family history from siblings can inform genetic and disease risk assessment.
\medterm{SIBO} Alternate index label; see \textit{Small intestinal bacterial overgrowth}.

\medterm{Sicca Syndrome} Persistent dryness of the eyes and mouth caused by reduced tears and saliva. It may result from Sjögren syndrome, medicines, dehydration, radiation, or other illnesses and can lead to eye irritation, mouth sores, tooth decay, or difficulty swallowing.


\medterm{Sick Sinus Syndrome} A problem with the heart’s natural pacemaker that causes the heartbeat to become too slow, pause, or alternate between fast and slow. It may cause dizziness, fainting, tiredness, palpitations, or reduced exercise ability.

\synonyms
Bradycardia-Tachycardia Syndrome

\medterm{Sickle Cell} An inherited group of blood disorders caused by abnormal hemoglobin. Red blood cells may become rigid and curved, causing anemia, severe pain, blocked blood flow, and damage to organs.

\medterm{Sickle Cell Anemia} An inherited blood disorder in which abnormal haemoglobin makes red blood cells stiff, sticky, and crescent-shaped. These cells break down easily and can block small blood vessels, causing anaemia, severe pain, stroke, lung problems, jaundice, and serious infections.

\medterm{Sickle Cell Anemia (Sickle Cell)} See \textit{Sickle Cell Anemia}.


\medterm{Sickle Cell Crisis} A sudden complication of sickle-cell disease, usually caused by sickled red cells blocking blood flow and causing severe pain. Other crises can damage the lungs, reduce blood-cell production, or trap blood in the spleen; urgent care may be needed.

\medterm{Sickle Cell Dactylitis} Painful swelling of the hands or feet caused by blockage of small blood vessels during a sickle-cell
episode, especially in infants and young children. It may be an early sign of sickle-cell disease
and requires appropriate clinical evaluation.

\medterm{Sickle Cell Disease} Alternate index label; see \textit{Sickle Cell}.

\medterm{Sickle Cell Test} A laboratory test detecting hemoglobin S or assessing sickling, used in screening and diagnosis of sickle-cell trait or sickle-cell disease; definitive testing identifies hemoglobin variants.

\medterm{Sickle Cell Trait} An inherited state in which a person has one gene for hemoglobin S and one usual hemoglobin gene. Most people have no symptoms and do not develop sickle cell disease, but rare problems can occur with severe dehydration, low oxygen, or extreme exertion.

\medterm{Sickle-Cell Anaemia} See \textit{Sickle Cell Anemia}.
\medterm{Sickness} A general term for illness or nausea; clinical assessment should identify the specific symptom or disease.
\medterm{Sickness Motion (Motion Sickness (Sea Sickness, Car Sickness))} Alternate index label for Motion Sickness (Sea Sickness, Car Sickness).

\medterm{Side Effect} An unintended effect that occurs while a medicine is used at a normal dose. It may be mild, such as nausea or drowsiness, or serious, such as bleeding, severe allergy, organ injury, or an abnormal heart rhythm.

\medterm{Side Effects} Unintended effects that occur while a medicine or treatment is used. They range from mild symptoms to serious harm; the risk depends on the treatment, dose, duration, other medicines, and the person’s health.

\medterm{Side-Effects} Unintended effects that occur while a medicine or treatment is used. They range from mild symptoms to serious harm; the risk depends on the treatment, dose, duration, other medicines, and the person’s health.
\medterm{Sideroblastic Anemia} Anemia caused by faulty hemoglobin production. Iron builds up inside developing red blood cells instead of being used properly; causes include inherited disorders, alcohol, some medicines or toxins, vitamin B6 deficiency, and bone-marrow disease.
\medterm{Siderosis} Accumulation of iron in body tissues, or lung disease from breathing iron-containing dust. It may cause no symptoms, but long-term dust exposure can produce cough, breathing difficulty, and permanent lung damage.
\medterm{SIDS} Sudden infant death syndrome, the sudden unexplained death of an infant younger than one year after thorough investigation.

\medterm{Sievert (Sv)} A unit used to estimate how much biological harm ionizing radiation could cause. It accounts for the type of radiation and the sensitivity of different tissues and is used when discussing radiation risk.

\synonyms
Sv

\medterm{Sight} The ability to see, involving the eyes and visual nervous system.
\medterm{Sigmoid} Sigmoid means S-shaped; the sigmoid colon is the curved section of large intestine between the descending colon and rectum.

\medterm{Sigmoid Colon} The curved, S-shaped section of the large intestine between the descending colon and rectum. It stores stool before bowel movements and is a common site for diverticular disease, twisting, and colorectal cancer.

\medterm{Sigmoid Volvulus} Twisting of the S-shaped lower colon around the tissue that supports it, blocking the bowel. It can cause severe abdominal swelling, pain, vomiting, and inability to pass stool or gas and needs urgent treatment because the blood supply may be cut off.

\medterm{Sigmoidoscopy} Endoscopic examination of the rectum and lower colon using a flexible or rigid sigmoidoscope. It can identify inflammation, bleeding, polyps, or tumors and may allow biopsy or treatment.

\medterm{Sign} An observable or measurable finding that may show disease or a change in health. Examples include a rash, fever, swelling, abnormal heartbeat, or high blood pressure; a sign differs from a symptom, which is felt and reported by the person.

\medterm{Sign Language} Sign language is a complete visual language using hand shape, movement, position, facial expression, and body movement rather than spoken sound.

\medterm{Signal Peptide} A short part of a newly made protein that acts like an address label, directing the protein to the cell membrane, outside the cell, or a particular internal compartment. It is often removed after delivery.

\medterm{Signal Transduction} The chain of events by which a hormone, nerve chemical, or medicine attaches to a cell and produces a response inside it. The response can change cell activity, movement, growth, survival, or gene use.

\medterm{Signal-To-Noise Ratio} A comparison between the useful information in a test or image and the unwanted background variation. A higher ratio usually produces a clearer image or more reliable measurement.
\medterm{Signature} A signature is a person’s identifying mark or electronic confirmation of approval, consent, authorship, or verification.

\medterm{Signature Strengths} Personal character strengths that a person values and naturally uses.

\medterm{Signet Ring Cell Carcinoma} A cancer whose cells contain mucus that pushes the cell nucleus to one side, giving a ring-like appearance under a microscope. It most often begins in the stomach or colon and may grow widely or aggressively.

\medterm{Signs Of An Asthma Attack} An asthma attack may cause worsening wheeze, cough, chest tightness, breathlessness, rapid breathing, difficulty speaking, or reduced peak flow; severe symptoms need emergency care.

\medterm{Sildenafil Citrate} The citrate salt of sildenafil, a phosphodiesterase-5 inhibitor used for erectile dysfunction and pulmonary arterial hypertension.
\medterm{Silent Heart Attack} A heart attack that causes no recognized pain or typical warning signs. It may be discovered later through an electrocardiogram or imaging and still damages heart muscle and increases future risks.

\synonyms
Unrecognized myocardial infarction

\medterm{Silent Ischemia} Reduced blood flow to heart muscle without typical chest pain or other recognized symptoms. It may still damage the heart and is detected through testing rather than symptoms alone.

\synonyms
Asymptomatic ischemia

\medterm{Silent Mutation} A DNA change that usually does not alter the amino acid in the protein made from the gene. It may still affect how the gene is copied, how long its RNA lasts, or how much protein is produced.

\medterm{Silent Sinus Syndrome} Silent sinus syndrome is painless inward collapse of the maxillary sinus and floor of the orbit, often causing facial asymmetry, enophthalmos, or a sunken eye. Imaging confirms the diagnosis, and endoscopic treatment restores sinus ventilation when indicated.

\medterm{Silent Thyroiditis} Silent thyroiditis is painless autoimmune inflammation of the thyroid that may briefly release excess hormone before causing temporary hypothyroidism and eventual recovery.

\medterm{Silicate} A salt or ester containing silicon and oxygen, found in minerals and industrial materials; inhalation of some silicate dusts can irritate the lungs or contribute to pneumoconiosis.

\medterm{Silicon Dioxide} A mineral found in soil, rock, dust, and industrial materials. Breathing crystalline forms over time can scar the lungs, cause silicosis, and increase the risk of lung cancer.

\medterm{Silicone} A family of synthetic silicon-containing polymers used in medical devices, implants, sealants, and consumer products; medical risks depend on the product and exposure.

\medterm{Silicone Elastomer} A flexible, rubber-like silicone material used in implants, tubing, dressings, and other devices; complications may include infection, leakage, allergy, or device failure.

\medterm{Silicone Implants} Medical implants containing silicone gel or elastomer used for breast reconstruction or augmentation; complications may include rupture, capsular contracture, infection, malposition, or need for revision.

\medterm{Silicone Mastitis} A granulomatous inflammatory reaction to silicone released from a breast implant or injected silicone. It may cause nodules, pain, lymphadenopathy, or imaging findings that mimic malignancy.

\medterm{Silicone Oil Hyperoleon} A complication after eye surgery in which silicone oil breaks into tiny droplets and moves into the front of the eye. It can raise eye pressure, irritate the eye, damage the cornea, and sometimes requires removal or washing of the oil.

\medterm{Silicones} Synthetic polymers used in medical devices, implants, dressings, and other products; adverse effects depend on formulation and site.
\medterm{Silicosis} A permanent lung-scarring disease caused by breathing crystalline silica dust. It can cause cough, breathlessness, reduced lung function, tuberculosis risk, and progressive respiratory disability after workplace or environmental exposure.
\medterm{Silk} Silk is a protein fiber used in textiles and some medical sutures or biomaterials; reactions and tissue response depend on the product and processing.

\medterm{Silk Road Disease (Behcet's Syndrome)} Alternate index label for Behcet's Syndrome.

\medterm{Siltuximab} A monoclonal antibody that binds interleukin-6 and is used to treat adults with multicentric Castleman disease who are negative for HIV and human herpesvirus 8.

\medterm{Silver} A metal used in some dressings, creams, and medical devices for antimicrobial activity; excessive exposure can cause permanent blue-gray skin discoloration called argyria.

\medterm{Silver Diamine Fluoride} A liquid applied to a tooth to slow or stop active tooth decay. It usually turns the treated decayed area black and may be useful when drilling or immediate restorative treatment is difficult.
\medterm{Silver Nitrate} A caustic antiseptic used topically to cauterize tissue, control minor bleeding, or treat some wounds; it can burn healthy skin and permanently stain tissue.

\medterm{Silver Poisoning} Silver poisoning, or argyria, follows excessive silver exposure and can cause permanent blue-gray skin discoloration and, rarely, organ or neurologic effects.

\medterm{Silver Sulfadiazine} A topical antimicrobial cream formerly used widely for burns; it may delay healing in some wounds and is not appropriate for everyone, including some newborns and people with sulfonamide allergy.

\medterm{Silver-Russell Syndrome} A growth disorder usually present from birth, marked by poor growth before and after birth, a small or asymmetrical body, and characteristic facial features. Most cases are linked to changes in gene activity rather than a change in the gene itself.

\medterm{Silymarin} A mixture of flavonolignans from milk thistle marketed as a dietary supplement; evidence for treating liver disease is limited and product quality varies.

\medterm{Simeprevir} An older hepatitis C virus NS3/4A protease inhibitor formerly used with other antivirals; it is no longer a usual first-line treatment because better direct-acting regimens are available.

\medterm{Simethicone} Simethicone is a non-absorbed antifoaming agent that coalesces gas bubbles in the gastrointestinal tract and is used for symptomatic relief of bloating; it does not treat the underlying cause of pain.

\medterm{Simian Crease} A simian crease is a single transverse crease across the palm, which may be a normal variant or occur with some chromosomal or developmental conditions.


\medterm{Simmonds Disease} An older term for severe pituitary-gland failure. It can cause extreme tiredness, weight loss, low blood pressure, loss of menstrual periods or sexual function, and deficiencies of several hormones.

\medterm{Simmonds' Disease} An older term for hypopituitarism, especially pituitary failure causing wasting and endocrine deficiency.
\medterm{Simple Febrile Seizure} A brief, generalized seizure caused by fever in a child, usually between six months and five years old. It lasts less than 15 minutes and does not recur during the same illness. The child should be assessed to find the cause of fever, but most children recover fully and do not develop epilepsy.

\synonyms
Benign Febrile Seizure; Typical Febrile Seizure

\medterm{Simple Goiter} A simple goiter is noncancerous enlargement of the thyroid without obvious thyroid-hormone excess or deficiency, caused by iodine imbalance, stimulation, or other factors.

\medterm{Simple Mastectomy} Simple mastectomy removes the breast tissue and nipple while leaving most chest muscles and usually the underarm lymph nodes.

\medterm{Simple Prostatectomy} Surgical removal of the enlarged central portion of the prostate to relieve bladder-outlet obstruction from benign prostatic hyperplasia; it does not remove the entire prostate.

\medterm{Simple Pulmonary Eosinophilia} A usually mild lung reaction in which eosinophils, a type of white blood cell, collect temporarily in the lungs. It may cause cough, fever, or wheezing and can follow medicines, parasites, or other exposures. Symptoms and chest shadows often clear on their own, but the cause should be assessed.

\medterm{Simple, Heart-Smart Substitutions} Easy food changes that may support heart health, such as choosing unsaturated oils instead of butter, whole grains instead of refined grains, and less processed or salty food.

\medterm{Simplexvirus} A genus of herpesviruses that includes herpes simplex virus types 1 and 2, which cause cold sores, genital herpes, and other infections.

\medterm{Simuliidae} The family of black flies; bites can cause itchy skin reactions, and some species transmit Onchocerca volvulus, the parasite that causes river blindness.

\medterm{Simultanagnosia} An inability to see more than one object, or the whole visual scene, at the same time even though basic eyesight may be adequate. It usually results from damage to brain areas that combine visual information.

\medterm{Simvastatin} A statin medicine that lowers LDL, often called bad cholesterol, by reducing cholesterol production in the liver. It helps lower the risk of heart attack and stroke; muscle pain, liver problems, and medicine interactions are possible.
\medterm{Sin Nombre Virus Infection} Alternate index label; see \textit{Hantavirus pulmonary syndrome}.

\medterm{Sincalide} A synthetic form of cholecystokinin used during gallbladder imaging to stimulate contraction and measure bile flow; it may cause abdominal pain, nausea, or flushing.

\medterm{Sindbis Virus} A mosquito-borne virus that can cause fever, rash, and joint pain. It occurs mainly in parts of Africa, Europe, Asia, and Australia.

\medterm{Sinding-Larsen-Johansson Syndrome} A painful overuse injury at the lower tip of the kneecap in growing children and adolescents. Pain usually worsens with running or jumping and improves as growth finishes and the irritated area heals.

\medterm{Sinecatechins} A topical green-tea catechin ointment used for external genital and perianal warts caused by human papillomavirus; it can cause local irritation and should not be used internally.

\medterm{Singer's Nodule} A small, usually noncancerous thickening on a vocal fold caused by repeated voice strain. It can cause lasting hoarseness, a rough voice, or reduced vocal range; voice therapy and reducing strain are often helpful.
\medterm{Single Nucleotide Polymorphism} A common inherited difference in one DNA building block. Most have little or no effect, but some are linked with disease risk, ancestry, or how a person responds to a medicine.

\medterm{Single Palmar Crease} A single transverse palmar crease is one crease crossing the palm instead of the usual two; it can be a normal variant or a feature of genetic syndromes.

\medterm{Single Photon Emission Computed Tomography} A nuclear medicine scan that uses a radioactive tracer and a rotating camera to create three-dimensional images of blood flow or organ function. It can assess the heart, bones, brain, kidneys, or other organs.

\medterm{Single Ventricle} A group of heart defects present at birth in which one lower chamber does most of the work of pumping blood to both the body and lungs. Babies usually need specialist care and staged operations.

\medterm{Single-Dose Toxicity} Harm caused by one exposure to a substance or medicine, assessed according to dose, route, timing,
and the exposed person's characteristics.

\medterm{Single-Fiber Electromyography} A specialized electrical test that records signals from individual muscle fibers. It measures small variations in the timing of nerve-to-muscle communication and can detect disorders such as myasthenia gravis, even when routine nerve tests are normal. A fine needle is used, and the test may be uncomfortable.

\synonyms
SFEMG; Single-Fiber EMG

\medterm{Single-Photon Absorptiometry} An older scan that used a small amount of radiation to estimate bone density, usually in an arm or other peripheral bone. More accurate bone-density methods have largely replaced it.

\synonyms
SPA

\medterm{Single-Photon Emission Computed Tomography} A nuclear-medicine scan that uses a radioactive tracer and a rotating camera to create three-dimensional images of blood flow or organ activity. It can assess the heart, brain, bones, kidneys, and other organs.

\medterm{Singlet Oxygen} A highly reactive form of oxygen involved in oxidative reactions, immune defenses, photodynamic therapy, and tissue injury from oxidative stress.

\medterm{Singultus} Medical term for hiccups; see \textit{Hiccups}.

\medterm{Sinoatrial Node} The heart's natural pacemaker in the right atrium that normally initiates each heartbeat.

\synonyms
Pacemaker, natural
\medterm{Sinuatrial Node} The sinoatrial node, a group of specialized cells in the right atrium that normally initiates the heartbeat and sets the cardiac rhythm.

\medterm{Sinus} A sinus is a cavity, channel, or tract; in anatomy it may mean a paranasal air space, venous channel, or abnormal passage from a diseased area to the surface.
\medterm{Sinus Arrhythmia} Sinus arrhythmia is normal variation in sinus-heart rate, commonly speeding during inhalation and slowing during exhalation, especially in young people.

\medterm{Sinus Bradycardia} Sinus bradycardia is a sinus rhythm slower than the expected rate; it can be normal during sleep or athletic conditioning but may also reflect drugs or disease.

\medterm{Sinus CT Scan} Computed tomography of the paranasal sinuses used to evaluate chronic or complicated sinusitis, anatomic obstruction, trauma, inflammatory disease, or tumors.

\medterm{Sinus Headache} Pain attributed to a sinus headache is often migraine, but sinus inflammation can cause facial pressure, congestion, reduced smell, and pain that worsens with bending. Fever, purulent discharge, persistent symptoms, or eye or neurologic signs guide evaluation.

\medterm{Sinus Histiocytosis} An increase in cleanup immune cells within the open spaces of a lymph node. It is often a response to nearby inflammation, infection, a tumor, or tissue breakdown and is interpreted with the rest of the biopsy.

\medterm{Sinus Lift} A dental procedure that adds bone beneath the lining of the maxillary sinus in the upper jaw. It may create enough support for a dental implant when the jawbone has become too thin.
\medterm{Sinus MRI Scan} Magnetic resonance imaging of the paranasal sinuses and surrounding soft tissues used to evaluate tumors, invasive infection, orbital or intracranial extension, and selected inflammatory disease.

\medterm{Sinus Node} The heart’s natural pacemaker, located in the right atrium.

\synonyms
Sinoatrial node

\medterm{Sinus Pain} Pressure or pain around the forehead, cheeks, nose, or eyes caused by sinus inflammation, infection, migraine,
or other conditions. Severe headache, eye swelling, vision change, high fever, or neurologic
symptoms requires urgent assessment.

\medterm{Sinus Rhomboidalis} A diamond-shaped depression forming the floor of the fourth brain ventricle, near structures controlling movement and vital functions.

\medterm{Sinus Rhythm} A heart rhythm started by electrical impulses from the sinoatrial node, the natural pacemaker. It is usually regular, although sinus rhythm can be too fast, too slow, or affected by illness.

\medterm{Sinus Tachycardia} A heart rhythm originating in the sinoatrial node with a rate faster than normal, commonly occurring as a physiologic response to exercise, fever, pain, dehydration, or stress.
\medterm{Sinus X-Ray} Plain radiographic imaging of the paranasal sinuses; it has limited sensitivity and is often replaced by CT when detailed evaluation is needed.

\medterm{Sinuses} Air-filled cavities within the skull bones lined by respiratory mucosa, connected to the
nasal cavity via ostia.

\medterm{Sinusitis} Inflammation of the lining of one or more paranasal sinuses. It may follow a viral respiratory infection or be associated with allergy, nasal inflammation, or other causes. Common features include nasal obstruction, discharge, facial pressure or pain, headache, and reduced smell. Acute and chronic forms differ in duration and pattern.
\medterm{Sinusitis And Sinus Problems} Alternate index label for Sinusitis.


\medterm{Sinusitis, Acute} Alternate index label; see \textit{Acute sinusitis}.

\medterm{Sinusitis, Chronic} Alternate index label; see \textit{Chronic sinusitis}.

\medterm{Sinusoid} A small, wide blood vessel with unusually open walls, found especially in the liver, spleen, and bone marrow. Its structure allows blood to exchange cells, proteins, nutrients, and waste with nearby tissue.

\medterm{Sinusoidal Obstruction Syndrome (Hepatic Veno-Occlusive Disease)} Liver injury caused by damage and blockage of the small veins and sinusoids inside the liver. It most often occurs after certain cancer treatments or blood-forming stem-cell transplantation. It may cause painful liver enlargement, jaundice, rapid weight gain, fluid in the abdomen, and portal hypertension.

\medterm{Sinusoidal Fetal Heart Rate Pattern} A rare, smooth, regular wave pattern on an unborn baby’s heart-rate tracing with little normal variation. It can indicate severe anemia or low oxygen and requires immediate assessment and monitoring.


\medterm{Sipuleucel-T} A personalized cancer immunotherapy made from a patient’s immune cells and used for selected advanced prostate cancers. The cells are trained to recognize a prostate-cancer target before being returned to the patient.
\medterm{Sirenomelia} A rare congenital malformation in which the lower limbs are partially or completely
fused. It is usually associated with severe abnormalities of the genitourinary,
gastrointestinal, vascular, and spinal systems and is often fatal in the neonatal
period.

\medterm{Sirolimus} Sirolimus is an mTOR inhibitor used for immunosuppression after transplantation and selected vascular or lymphatic disorders; mouth ulcers, hyperlipidaemia, cytopenias, infection, impaired wound healing, and pneumonitis can occur.

\medterm{Sisomicin} An aminoglycoside antibiotic related to gentamicin that has activity against some aerobic gram-negative bacteria; use is limited by kidney and hearing toxicity.

\medterm{Sister Mary Joseph Nodule} A metastatic tumor deposit in or around the umbilicus, usually from an intra-abdominal or pelvic malignancy.

\medterm{Sister Mary Joseph's Nodule} A firm, hard nodule in the umbilicus resulting from metastatic malignancy, most commonly
from intra-abdominal or pelvic cancers including gastric, ovarian, and colorectal
carcinoma. It indicates advanced, disseminated disease with a poor prognosis. Diagnosis
is by fine needle aspiration biopsy.

\medterm{Sitagliptin Phosphate} Sitagliptin phosphate is a DPP-4 inhibitor that prolongs incretin activity and improves glucose-dependent insulin secretion in type 2 diabetes; dose adjustment in kidney disease and rare pancreatitis or severe joint pain are important considerations.


\medterm{Sitosterolemia} A rare inherited disorder in which the body absorbs and stores too much plant sterol, a substance found in plant foods. The buildup can cause early artery disease, skin lumps, anemia, easy bleeding, or low platelets.


\medterm{Sitting Position} An upright or semiupright position in which the patient's torso is elevated and the hips and knees are flexed; it is used for selected shoulder and neurosurgical procedures.

\medterm{Situational Syncope} A reflex faint triggered by a specific situation such as coughing, urination, defecation, swallowing, laughter, or prolonged standing, caused by transient reduced cerebral perfusion.

\medterm{Situs Ambiguus} A birth difference in which internal organs are arranged irregularly rather than in the usual left-right pattern. It is often associated with congenital heart defects and an absent, small, or abnormally functioning spleen.

\medterm{Situs Inversus} A congenital condition in which the thoracic and abdominal organs are reversed or
mirrored from their normal positions. Complete situs inversus (situs inversus totalis)
involves all organs.

\medterm{Situs Inversus Totalis} A congenital condition in which the major visceral organs are mirrored from their normal
positions, with the liver on the left, heart on the right, and other viscera similarly
reversed.

\medterm{Sitz Bath} Immersion of the perineal and anal region in warm or prescribed water for symptom relief in conditions such as hemorrhoids, anal fissures, or postpartum discomfort.

\medterm{Six Fingers Or Toes} Six fingers or toes is polydactyly, a congenital extra-digit difference that may occur alone or as part of a genetic syndrome.

\medterm{Six-Minute Walk Test} A simple test measuring how far a person can walk in six minutes on a flat path. Clinicians record distance, symptoms, heart rate, and oxygen level to assess exercise ability and monitor heart or lung disease. Rest breaks are allowed, and the result is interpreted with other findings.

\synonyms
6MWT; Six-Minute Walk; 6-Minute Walk Test

\medterm{Sixth Disease (Roseola)} Alternate index label for Roseola.

\medterm{Size Perception} The brain’s judgment of how large or small an object or body part is. Vision, touch, distance, and context all contribute; eye or brain disorders can make objects appear unusually large or small.

\medterm{Sjogren Syndrome} Older name for Sjögren disease, an autoimmune disease that commonly causes dry eyes and dry mouth.

\medterm{Sjögren Disease} The current preferred name for Sjögren syndrome, an autoimmune disease that commonly causes dry eyes and dry mouth and may affect other organs.

\medterm{Sjogren-Larsson Syndrome} An inherited disorder causing dry, scaly skin, muscle stiffness, and problems with movement or development. It results from changes in a gene involved in breaking down certain fats.
\medterm{SjöGren Syndrome} Older spelling and name for Sjögren disease.

\medterm{SjöGren's Disease} Alternate index label; see \textit{Sjögren disease}.

\medterm{SjöGren's Syndrome} Older name for Sjögren disease.
\synonyms
SjöGren's Disease

\medterm{SJS} Alternate index label; see \textit{Stevens-Johnson syndrome}.

\medterm{Skeletal Development} Skeletal development is formation, growth, remodeling, and maturation of bones and cartilage from embryonic life through adulthood.

\medterm{Skeletal Dysplasia} An inherited disorder of bone or cartilage development that may cause short or uneven growth, unusually shaped limbs, spine problems, or other differences in the skeleton. Effects vary widely between conditions.

\medterm{Skeletal Limb Abnormalities} Skeletal limb abnormalities are differences in the formation, length, alignment, or structure of an arm or leg caused by genetic, developmental, vascular, injury, or other factors.

\medterm{Skeletal Muscle} Striated muscle attached to bones and generally under voluntary control.
\medterm{Skeletal Muscle Tissue} Voluntary striated muscle attached to bones or other structures, composed of contractile fibers that generate movement, posture, heat, and joint stability.

\medterm{Skeletal Muscles} Muscles attached to bones by tendons that produce movement, maintain posture, and generate heat. They are usually under voluntary control but also support automatic actions such as breathing.

\medterm{Skeletal Myocytes} Skeletal myocytes are striated muscle cells responsible for voluntary movement, posture, heat production, and stabilization of joints.

\medterm{Skeletal Survey} A set of X-rays of most of a child's bones, used especially when multiple injuries or a medical condition affecting bone development is suspected. It can show fractures, bone abnormalities, and signs of healing and must be interpreted with the full clinical assessment.

\medterm{Skeletal System} The skeletal system consists of bones, cartilage, joints, and supporting ligaments that provide structure, protect organs, enable movement, and store minerals.

\medterm{Skeleton} The body’s framework of bones and cartilage that supports movement and protects organs. It also stores minerals, contains marrow that makes blood cells, and provides attachment points for muscles and ligaments.
\medterm{Skeleton, Bones Of The} The body’s framework of bones and cartilage, providing support, protection, leverage for movement, mineral storage, and marrow-based blood-cell production.

\medterm{Skeletonization} Final stage of postmortem decomposition where soft tissues, organs, and skin have
completely decayed or been scavenged, leaving only clean osseous skeletal remains.


\medterm{Skene Gland Cyst} A fluid-filled sac that forms when a Skene gland, also called a paraurethral gland, becomes blocked. These glands are located beside the opening of the urethra. Small cysts may cause no symptoms. Larger or infected cysts may cause pain, painful urination, difficulty passing urine, pain during sex, recurrent urinary infections, or an abscess. Diagnosis is usually made by pelvic examination. Symptomatic cysts may need surgical removal or marsupialization.

\medterm{Skilled Nursing Facility} A facility where people receive 24-hour nursing care and rehabilitation after illness, injury, surgery, or a hospital stay. Care may include wound care, medicines, injections, therapy, and help with daily activities until the person can return home or move to another setting.

\medterm{Skilled Nursing Or Rehabilitation Facilities} Facilities that provide 24-hour nursing care, rehabilitation, medication management, and assistance with daily activities after illness, injury, surgery, or functional decline.

\medterm{Skin} The body's outer covering, consisting mainly of the epidermis and dermis with appendages such as hair, nails, and glands. It protects against injury, infection, ultraviolet light, and water loss; detects touch, pressure, pain, and temperature; helps regulate body temperature; and contributes to immune defense and vitamin D production.
\medterm{Skin Lesion} An abnormal area of skin that may result from infection, inflammation, injury, allergy, or a benign or malignant growth.

\medterm{Skin - Clammy} Skin that feels cool, damp, or unusually sweaty. It may occur with heat, anxiety, pain, low blood sugar, shock, or other illness.

\medterm{Skin Abscess} A painful, swollen pocket of pus under the skin, usually caused by a bacterial infection. It may need drainage and antibiotics.

\medterm{Skin Absorption} Movement of a substance through the skin into underlying tissues or the bloodstream. Absorption depends
on the substance, skin condition, contact time, and exposed area.

\medterm{Skin Aging} Natural aging and long-term sun exposure gradually thin and dry the skin, reduce elasticity, slow healing, and cause wrinkles or dark spots. Limiting ultraviolet exposure, using protective clothing, and applying broad-spectrum sunscreen can reduce further damage.

\medterm{Skin Aging And Appearance} Skin aging is a complex biological process influenced by both intrinsic factors, such as
genetics and natural aging, and extrinsic factors, particularly ultraviolet radiation
exposure.

\medterm{Skin Allergies And Contact Dermatitis} Inflammatory skin reactions triggered by exposure to allergens or irritants that come
into direct contact with the skin.

\medterm{Skin Allergy Testing} Tests in which small amounts of suspected allergens are placed on or just under the skin and the reaction is observed. They can show sensitization to foods, medicines, venoms, or environmental allergens but must be interpreted with symptoms and exposure history.


\medterm{Skin And Aging} Skin aging involves intrinsic and extrinsic changes. Intrinsic aging causes thinning,
wrinkling, and fragility. Extrinsic (photo)aging causes wrinkles, telangiectasias,
lentigines, and actinic keratoses. Collagen loss, elastin degradation, and decreased
cell turnover contribute to aging skin.

\medterm{Skin And Environment} Environmental factors significantly affect skin health. Pollution causes oxidative
stress, inflammation, and accelerated aging. Particulate matter penetrates skin,
generating free radicals. Blue light from screens causes hyperpigmentation, particularly
in darker skin tones. Indoor heating and air conditioning cause dryness.

\medterm{Skin And Hair Changes During Pregnancy} Pregnancy can alter pigmentation, hair growth or shedding, acne, stretch marks, and vascular features through hormonal and circulatory changes; concerning lesions need assessment.

\medterm{Skin Appendages} Structures that grow from or lie within the skin, including hair follicles, oil glands, and sweat glands. They help lubricate and cool the body and can be affected by acne, infection, inflammation, or inherited disorders.

\medterm{Skin Barrier Function} The outer skin barrier limits water loss and helps keep out irritants, allergens, and germs. Its tightly packed surface cells and natural oils maintain protection; injury or inflammation can make skin dry, cracked, and more vulnerable.
\medterm{Skin Barrier} Alternate term; see \textit{Skin Barrier Function}.

\medterm{Skin Biopsy} Removal of a small skin sample for examination under a microscope. It can help diagnose skin cancer, infection, inflammatory disease, blistering disease, or an unexplained rash; the method depends on the lesion.

\medterm{Skin Biopsy Methods} The sampling method is chosen according to the lesion’s appearance and depth. A punch takes a round full-thickness sample, a shave removes a superficial layer, and an excision removes the whole lesion when complete removal is appropriate.

\medterm{Skin Biopsy Techniques} Skin biopsy is selected according to the suspected diagnosis and lesion depth. Punch
biopsy samples full-thickness skin; shave biopsy samples superficial lesions; excisional
biopsy removes an entire small lesion; and incisional biopsy samples part of a large
lesion.

\medterm{Skin Blushing/Flushing} Temporary reddening or warming of the skin caused by widening of small blood vessels. It may result from heat, emotion, exercise, fever, medicines, or hormone changes.

\medterm{Skin Cancer} Cancer that begins in cells of the skin. Common types include basal cell carcinoma, squamous cell carcinoma, and melanoma; less common types arise from other skin structures. Risk is influenced by ultraviolet exposure, immune suppression, genetic susceptibility, and certain preexisting lesions. A skin cancer may appear as a changing, bleeding, growing, scaly, or nonhealing spot.
\synonyms
Cancer, Skin

\medterm{Skin Cancer Epidemiology} Skin cancers are common and include basal-cell carcinoma, squamous-cell carcinoma, and
melanoma. Incidence and mortality vary by tumor type, geography, skin phenotype, ultraviolet
exposure, immune status, and access to diagnosis and treatment.

\medterm{Skin Cancer Screening} Examination of the skin for new or changing spots, especially in people at increased risk. Clinicians may use a full-body examination, a magnifying instrument, or a biopsy; routine screening recommendations vary by individual risk.

\medterm{Skin Cancer Staging} Skin cancer staging guides treatment and prognosis. Melanoma: Breslow thickness (most
important prognostic factor), ulceration, mitotic rate, sentinel lymph node status (AJCC
8th edition).

\medterm{Skin Cancer Surgery} Surgery to remove skin cancer with a margin of apparently healthy tissue. The method depends on the cancer type, size, location, and spread and may include standard excision or Mohs surgery, which checks the edges during removal.

\medterm{Skin Culture} Laboratory testing of a swab, scraping, fluid, or tissue sample from the skin to identify bacteria, fungi, or other organisms causing infection.

\medterm{Skin Cancer, Melanoma} Alternate index label; see \textit{Melanoma}.

\medterm{Skin Carcinoma} Cancer arising from skin cells, most commonly basal-cell or squamous-cell carcinoma. It may appear as a changing sore, lump, scaly patch, or wound that does not heal and needs examination.

\medterm{Skin Care} Measures that protect and maintain the skin, including gentle cleansing, moisturization, sun protection, wound care, avoidance of irritants, and treatment of identified disease.

\medterm{Skin Care And Incontinence} Skin protection for people with urine or stool leakage, including gentle cleansing, thorough drying, absorbent products, barrier creams, regular changing, and treatment of the underlying cause.

\medterm{Skin Color} The visible color of skin, determined mainly by melanin, blood flow, oxygenation, and other pigments.

\medterm{Skin Erosion} A shallow loss of the epidermis that leaves a moist or crusted surface but does not extend deeply into the dermis, often caused by friction, scratching, blister rupture, or inflammation.

\medterm{Skin Scraping} Removal of a small amount of surface skin scale or material for laboratory examination, commonly to look for fungi, mites, or other causes of a rash.

\medterm{Skin Examination} Systematic inspection and palpation of the skin, hair, nails, and mucous membranes to assess color, lesions, texture, distribution, vascular signs, infection, and cancer warning features.

\medterm{Skin Findings In Newborns} Common skin changes in newborns, such as milia, newborn rash, peeling, birthmarks, or bluish hands and feet. Most are harmless, but fever, blisters, widespread redness, or a sick appearance needs prompt assessment.

\medterm{Skin Flaps} Pieces of skin and underlying tissue moved to cover a wound while keeping or reconnecting their blood supply. Flaps may come from nearby tissue or another part of the body and are used after injury, cancer removal, or reconstructive surgery.


\medterm{Skin Graft} Skin taken from one area of the body and placed over a wound, burn, ulcer, or surgical defect. A graft must obtain blood from the wound surface to survive and may fail, scar, change color, or require further surgery.

\medterm{Skin Graft, Pedicle} A piece of skin moved to cover a wound while remaining attached to its original blood supply through a tissue bridge.

\medterm{Skin Immune System} The cells and protective substances in skin that recognize germs, remove damaged cells, and control inflammation. It works with the skin barrier to prevent infection and support healing after injury.

\medterm{Skin Induration} Unusual hardness or firmness of the skin or tissue beneath it. Causes include scar formation, inflammation, swelling, infection, an infiltrating disease, or a growth; examination determines the cause.

\medterm{Skin Infection} Infection of the skin or underlying soft tissue, commonly caused by bacteria, fungi, or viruses. Redness, warmth,
swelling, pain, drainage, or fever may occur; rapidly spreading redness, severe pain, blisters, or systemic
illness requires prompt medical assessment.

\medterm{Skin Integrity} The condition of skin and underlying tissue when it remains intact and protected from injury. Prolonged pressure, friction, moisture, poor blood flow, and illness can cause breakdown, especially over bony areas.

\medterm{Skin Irritation} Inflammation or discomfort caused by direct contact with an irritant, often producing redness, burning,
itching, dryness, or scaling.

\medterm{Skin Lesion Aspiration} Removal of fluid or cells from a skin lump with a needle for testing or drainage.

\medterm{Skin Lesion Biopsy} Removal of part or all of a skin lesion for microscopic examination to diagnose infection, inflammatory disease, benign tumor, or skin cancer.

\medterm{Skin Lesion KOH Exam} Microscopic examination of skin scrapings treated with potassium hydroxide to detect fungal hyphae or yeast in suspected superficial fungal infection.

\medterm{Skin Lesion Of Blastomycosis} A blastomycosis skin lesion results from dissemination of Blastomyces fungal infection and may appear as verrucous or ulcerated lesions with raised borders.

\medterm{Skin Lesion Removal} Removal of an unwanted, suspicious, or diseased area of skin by shaving, cutting, freezing, scraping, or another method. Tissue may be sent for laboratory examination; bleeding, infection, scarring, and incomplete removal are possible.


\medterm{Skin Lumps} Skin lumps are localized bumps or masses caused by cysts, lipomas, inflamed follicles, infections, benign tumors, or less commonly cancer; persistence or change warrants assessment.

\medterm{Skin Mastocytoma} A localized collection of mast cells in the skin, usually a yellow-brown or reddish plaque or nodule in a child that may swell, itch, or redden when rubbed.

\medterm{Skin Microbiome} The community of bacteria, fungi, viruses, and other microorganisms normally living on the skin. They interact with the skin and immune system; changes in this community may contribute to some skin diseases.

\medterm{Skin Necrosis} Death of skin tissue from ischemia, infection, pressure, toxins, trauma, vasculitis, or severe inflammation, producing discoloration, blisters, ulcers, or black eschar.

\medterm{Skin Neoplasm} A skin neoplasm is an abnormal growth in the skin that may be benign, precancerous, or malignant.

\medterm{Skin Nodules} Firm, raised lumps in or under the skin. They may result from inflammation, infection, injury, cysts, or benign or malignant growths.

\medterm{Skin Of Color} A broad term for people with more deeply pigmented skin. Some conditions may appear differently, and dark marks or raised scars may develop more readily after inflammation or injury; skin cancer can occur in every skin tone.

\medterm{Skin Or Nail Culture} Laboratory culture of skin, hair, or nail material to identify bacteria or fungi when infection is suspected; results require clinical correlation.

\medterm{Skin Pigmentation} Skin color is produced mainly by melanin made by pigment-producing cells. Sun exposure, hormones, inflammation, medicines, and illness can make areas darker or lighter; new or changing discoloration may need examination.


\medterm{Skin Resurfacing} Procedures that improve skin texture, scars, wrinkles, or uneven color by removing or changing surface skin layers. Methods include lasers, chemical peels, and abrasion; risks include infection, pigment changes, and scarring.

\medterm{Skin Self-Exam} A skin self-exam is a regular check of the skin for new or changing moles, sores, lumps, or patches that may need medical assessment.

\medterm{Skin Tag} A skin tag is a benign soft growth of connective tissue, commonly in body folds and associated with friction, age, pregnancy, obesity, or insulin resistance. Removal may be performed for irritation or appearance, but a changing or atypical lesion should be diagnosed first.

\synonyms
Acrochordon (Skin Tag)

\medterm{Skin Tags} Soft, usually harmless skin growths that commonly occur in body folds such as the neck, armpits, groin, or beneath the breasts. They may become irritated or bleed; a changing growth should be examined before removal.

\medterm{Skin Test} Alternate index label for diagnostic or immunologic skin tests, including prick, patch, and tuberculin testing.
\medterm{Skin Turgor} The skin's ability to return to its normal position after being gently pinched. Reduced turgor may occur with dehydration, although it is less reliable in older adults.

\medterm{Skin Writing} A skin condition in which light scratching or pressure causes temporary raised red lines or welts.

\medterm{Skin, Diseases Of} Conditions affecting skin, hair, nails, glands, or subcutaneous tissue, including infection, inflammation, neoplasia, and genetic disorders.
\medterm{Skin, Itchy} Alternate index label; see \textit{Itchy skin (pruritus)}.

\medterm{Skin-Grafting} Surgery that moves skin from a donor area to cover a wound, burn, ulcer, or tissue defect. The graft needs blood supply from the wound bed and may fail, scar, or change color.
\medterm{Skipped Heartbeats} Alternate index label; see \textit{Heart palpitations}.

\medterm{Skull} The bony framework of the head that protects the brain and supports the face, eyes, ears, and jaw.
\medterm{Skull Fracture} A break in one or more skull bones caused by trauma. It may occur with bleeding inside the skull, brain injury, or leakage of clear fluid from the nose or ear and requires medical assessment.

\medterm{Skull X-Ray} A skull X-ray uses low-dose radiation to show skull bones and may help assess selected fractures, foreign bodies, or bone abnormalities.

\medterm{Slanted Ear} An ear whose upper edge slopes unusually compared with typical ear orientation; it may be isolated or part of a genetic condition.

\medterm{SLAP Lesion} A tear in the upper rim of the shoulder socket, sometimes where the biceps tendon attaches. It may follow a fall, lifting injury, or repeated overhead activity and can cause pain, clicking, weakness, or instability.

\medterm{SLAP Tears} Tears of the superior glenoid labrum extending from anterior to posterior, sometimes involving the
biceps tendon attachment. They may cause pain, clicking, or instability, especially with overhead activity.

\medterm{Slapped Cheek Disease} Alternate index label; see \textit{Parvovirus infection}.

\medterm{Slapped Cheek Syndrome} Erythema infectiosum, a parvovirus B19 infection causing a slapped-cheek rash and sometimes joint or hematologic complications.
\medterm{SLE} Systemic lupus erythematosus, a chronic autoimmune disease that can affect skin, joints, kidneys, blood, nervous system, and other organs.
\medterm{Sleep} A naturally recurring state in which awareness and responses to the surroundings decrease while the brain and body follow organized patterns. Sleep supports learning, memory, emotional health, physical recovery, immune function, and metabolism.
\synonyms
Disorders Sleep (Sleep)

\medterm{Sleep (Geriatric)} Sleep patterns and sleep problems in older adults, including insomnia, sleep apnea, circadian changes,
and medication-related disturbances.

\synonyms
Sleep care for older adults; geriatric sleep problems

\medterm{Sleep And Sleep Disorders In Children And Teens} Children and teens need age-appropriate sleep, and disorders may involve insomnia, circadian delay, sleep apnea, parasomnias, restless legs, or insufficient opportunity. Snoring, pauses in breathing, unusual daytime sleepiness, mood change, or poor school performance warrants evaluation.

\medterm{Sleep Apnea} Repeated pauses or reductions in breathing during sleep. It can result from the airway closing, the brain failing to signal breathing, or both, and may cause loud snoring, gasping, poor sleep, daytime sleepiness, and low oxygen.

\medterm{Sleep Apnea (Dental)} Alternate index label for Oral Appliance Therapy.

\medterm{Sleep Apnea (Obstructive Sleep Apnea - OSA)} Alternate index label for Obstructive Sleep Apnea.


\medterm{Sleep Apnea, Central} Alternate index label; see \textit{Central sleep apnea}.

\medterm{Sleep Apnea, Obstructive} Alternate index label; see \textit{Obstructive sleep apnea}.

\medterm{Sleep Apnoea} Alternate index label for Sleep Apnea.

\medterm{Sleep Architecture} The pattern and order of sleep stages during a night or other sleep period. It includes lighter sleep, deep sleep, and dream-related REM sleep; illness, medicines, age, stress, and sleep disorders can change how much time is spent in each stage.

\medterm{Sleep Bruxism} Repeated grinding or clenching of the teeth during sleep, which may cause tooth wear, jaw pain, or headaches.

\medterm{Sleep Debt} The accumulated difference between the amount of sleep a person needs and the amount obtained over time; it may impair alertness, mood, cognition, and performance.

\medterm{Sleep Diary} A daily record of sleep and wake times, awakenings, naps, medicines, caffeine, and related factors used to assess sleep patterns and treatment response.

\medterm{Sleep Disorder Periodic Limb Movement (Periodic Limb Movement Disorder)} Alternate index label for Periodic Limb Movement Disorder.

\medterm{Sleep Disorders} A group of conditions affecting sleep quality, timing, duration, breathing, movement, or behavior. Major groups include insomnia, sleep-related breathing disorders, central disorders of hypersomnolence, circadian rhythm sleep-wake disorders, parasomnias, and sleep-related movement disorders. Examples include sleep apnea, narcolepsy, restless legs syndrome, and sleep paralysis.

\medterm{Sleep Disorders In Older Adults} Sleep problems in later life, including insomnia, sleep apnea, restless legs, and circadian changes; medication effects, pain, depression, and medical illness should be assessed.

\medterm{Sleep Efficiency} The percentage of time in bed that a person is actually asleep. It is calculated by dividing total sleep time by time in bed; a low value may occur with insomnia, repeated awakenings, or other sleep disorders.

\medterm{Sleep Fragmentation} Repeated brief awakenings or partial arousals that interrupt otherwise continuous sleep. It can reduce restorative sleep and occur with sleep apnea, pain, noise, movement disorders, medicines, or illness.

\medterm{Sleep Hygiene} Habits and surroundings that support regular, restorative sleep. Helpful measures include a consistent schedule, a dark and quiet bedroom, relaxing before bed, and limiting caffeine, alcohol, nicotine, and stimulating screen use near bedtime.

\medterm{Sleep Medications And Supplements} Pharmacological agents used to treat insomnia and other sleep disorders when behavioral
interventions alone are insufficient.

\medterm{Sleep Onset Latency} The time required to fall asleep after attempting to sleep; prolonged latency may occur with insomnia, anxiety, circadian misalignment, stimulants, or another sleep disorder.

\medterm{Sleep Paralysis} Temporary inability to move or speak when falling asleep or waking, while awareness is partly or fully preserved.
Episodes usually last seconds to minutes and may include frightening hallucinations. Recurrent episodes
can be associated with sleep deprivation, irregular sleep, narcolepsy, or other sleep disorders.

\medterm{Sleep Related Breathing Disorders} Sleep-related breathing disorders cause abnormal ventilation or oxygenation during sleep, including obstructive and central sleep apnea and sleep-related hypoventilation. Snoring, witnessed pauses, gasping, morning headache, or daytime sleepiness may prompt sleep testing and treatment.

\synonyms
Breathing Disorders Sleep Related (Sleep Related Breathing Disorders)

\medterm{Sleep Spindles} Short bursts of regular brain activity that occur mainly during stage 2 non-REM sleep. They are seen on a sleep recording and are thought to help protect sleep and process newly learned information.

\medterm{Sleep Studies (Polysomnography)} Overnight tests that record brain activity, eye movements, heart rate, breathing, blood oxygen, body position, and limb movements. They help identify sleep apnea, unusual movements, seizures during sleep, and other sleep disorders.

\synonyms
Polysomnography

\medterm{Sleep, Disorders Of} Health conditions that disturb how much a person sleeps, when they sleep, or how restful sleep is. They may involve insomnia, breathing pauses, unusual movements, excessive sleepiness, unusual nighttime behavior, or a body clock that is out of step.
\medterm{Sleep-Related Hypoventilation} Breathing that is too shallow during sleep to remove enough carbon dioxide or maintain normal oxygen. It can result from obesity, lung or muscle disease, chest-wall problems, brain disorders, or sedating medicines.

\medterm{Sleep-Related Hypoxemia} Abnormally low blood oxygen during sleep that may occur with sleep-disordered breathing, lung disease, neuromuscular weakness, hypoventilation, or other medical conditions.

\medterm{Sleepiness, Daytime} Alternate index label for Excessive Daytime Sleepiness.

\medterm{Sleeping Sickness} An infection caused by Trypanosoma parasites and spread by tsetse flies in sub-Saharan Africa. It may begin with fever and swollen nodes, then progress to sleep disturbance, confusion, and other nervous-system problems.
\medterm{Sleeplessness} Alternate lay term; see \textit{Insomnia Disorder}.

\medterm{Sleepwalking} Recurrent walking or other complex behavior during sleep, usually with little memory afterward. It is more common in children and may be triggered by sleep deprivation, stress, fever, medicines, or other sleep disorders; injury prevention is important.

\medterm{Sleepwalking Disorder} A sleep disorder involving repeated episodes of walking or other complex behavior during sleep, usually with little memory afterward, that cause injury, distress, sleep disruption, or impaired daytime functioning.

\medterm{Sleepwalking And Children} Sleepwalking in children is usually benign and often resolves with age; regular sleep, hazard prevention, and evaluation of frequent, violent, injurious, or atypical episodes are important.

\medterm{Sleeve Gastrectomy} An operation that removes most of the stomach, leaving a narrow sleeve-shaped section. It limits how much food the stomach can hold and changes appetite hormones; bleeding, leakage, reflux, and nutritional problems can occur.
\medterm{Sleeve Lobectomy} An operation that removes a lung lobe and the nearby affected part of the main breathing tube, then reconnects the remaining airway. It can remove a centrally located tumor while preserving more healthy lung than removing the whole lung. Possible complications include air leakage, narrowing of the airway, infection, and bleeding.

\synonyms
Bronchial Sleeve Resection; Sleeve Resection of the Lung

\medterm{Sliding Filament} The process by which muscles contract: two protein fibers slide past each other, shortening the muscle’s working units. Calcium and energy from ATP allow the fibers to attach, pull, release, and repeat the cycle.

\medterm{Sliding Hiatal Hernia} A sliding hiatal hernia occurs when the stomach and gastroesophageal junction move upward through the diaphragm, often contributing to reflux.

\medterm{Sling} A supportive device or surgically placed strip used to support an injured limb or urinary structures.

\medterm{Slipped Capital Femoral Epiphysis} A hip problem in growing children or teenagers in which the upper end of the thigh bone slips through its growth area. It may cause hip, groin, thigh, or knee pain and a limp and needs urgent assessment to prevent further slipping.

\medterm{Slipped Disk} A spinal disk whose tough outer layer tears or weakens, allowing its inner material to bulge and irritate nearby nerves. It may cause back or neck pain, numbness, weakness, or sciatica.
\synonyms
Herniated disk

\medterm{Slipped Vertebra} Forward movement of one spinal bone compared with the bone below it, also called spondylolisthesis. It may cause back pain, stiffness, hamstring tightness, or pressure on a nerve, although some people have no symptoms.

\synonyms
Spondylolisthesis

\medterm{Slipping Rib Syndrome} Pain caused by excessive movement of the lower ribs, usually from loosened cartilage. The pain may be sharp and can sometimes be reproduced by pressing or moving the rib.

\medterm{Slit Lamp} An eye-examination microscope that uses magnification and a narrow beam of light to inspect the eyelids, cornea, lens, and other eye structures. It helps detect infection, injury, cataract, and glaucoma-related changes.
\synonyms
Biomicroscope

\medterm{Slit Lamp Biomicroscopy} A detailed eye examination using a bright, narrow beam of light and a microscope to view the front of the eye. Additional lenses can help examine the retina and other structures farther back.

\medterm{Slit-Lamp Biomicroscope} Alternate index label for Slit Lamp.

\medterm{Slit-Lamp Exam} Alternate index label for Slit Lamp Biomicroscopy.

\medterm{SLL} Alternate index label; see \textit{Small lymphocytic lymphoma}.

\medterm{Slough} Soft, dead, or detached tissue covering or separating from a wound, ulcer, or inflamed surface. It can delay healing and may need gentle removal after blood flow, infection, and the wound’s depth are assessed.
\medterm{Slow Heart Rate - Bradycardia} A heart rate slower than expected, often below 60 beats per minute in an adult. It may be normal in some people but can cause dizziness, fainting, tiredness, or shortness of breath when it reduces blood flow.

\medterm{Slow Heartbeat} Alternate index label; see \textit{Bradycardia}.

\medterm{Slow-Twitch Fiber} A type of muscle fiber suited to steady activity over a long time. It uses oxygen efficiently and resists tiredness, helping with posture, walking, and endurance activities rather than short bursts of maximum force.

\synonyms
Type I muscle fiber

\medterm{Slow-Wave Sleep} The deepest stage of non-REM sleep, marked by slow brain waves and reduced response to sounds or touch. It supports physical recovery, memory, and hormone release and is often reduced by ageing, illness, alcohol, or sleep disruption.

\synonyms
Deep sleep

\medterm{Small Airway Disease} Narrowing, inflammation, or blockage of the smallest breathing tubes in the lungs. It may cause cough, wheezing, breathlessness, or trapped air and can occur with asthma, COPD, infection, or other lung disease.

\medterm{Small Bowel Adenocarcinoma} A malignant gland-forming epithelial tumor of the small intestine. It is uncommon and may cause occult bleeding, obstruction, abdominal pain, or weight loss; risk is increased in selected inherited and inflammatory disorders.

\medterm{Small Bowel Bacterial Overgrowth} Excess bacteria in the small intestine, where fewer bacteria are normally present. It can cause bloating, abdominal discomfort, diarrhea, weight loss, and poor absorption of nutrients; some cases involve methane-producing organisms instead.

\medterm{Small Intestinal Fungal Overgrowth} A proposed condition in which excess fungi are found in the small intestine and may be associated with bloating, nausea, diarrhea, or abdominal discomfort. Symptoms are nonspecific, testing is not standardized, and other causes should be assessed before attributing symptoms to it.

\medterm{Small Bowel Follow-Through} A series of X-rays taken after swallowing contrast, usually barium, to follow its movement through the small intestine. It can show narrowing, obstruction, inflammation, or abnormal transit.

\medterm{Small Bowel Bleeding} Bleeding from the small intestine. It may cause black or maroon stool, hidden blood loss, or iron-deficiency anemia and may require capsule endoscopy, enteroscopy, imaging, or other tests to find the source.

\medterm{Small Bowel Obstruction} Partial or complete blockage of the small intestine. Common causes include adhesions,
hernias, tumors, and inflammatory narrowing. Symptoms may include abdominal pain,
vomiting, distension, and inability to pass stool or gas.

\medterm{Small Bowel Prolapse} Alternate index label; see \textit{Small bowel prolapse (enterocele)}.

\medterm{Small Bowel Resection} Surgical removal of a diseased segment of the small intestine followed by reconnection or stoma formation when needed.


\medterm{Small Bowel Series} A series of X-ray images of the small intestine after swallowing contrast material, usually barium. It can show narrowing, blockage, abnormal movement, inflammation, tumors, or changes from Crohn disease.

\medterm{Small Bowel Tissue Smear/Biopsy} Examination of cells or a tissue sample from the small intestine to investigate poor nutrient absorption, celiac disease, inflammation, infection, or a tumor.

\medterm{Small Calorie (Calorie)} A small calorie is a unit of heat energy; food labels usually use a kilocalorie, written Calorie with a capital C, equal to 1,000 small calories.

\medterm{Small Cell Carcinoma} A fast-growing cancer made of small cells, most often beginning in the lung and strongly linked to smoking. It may shrink at first with treatment but often returns or spreads early, including to the brain.

\medterm{Small Cell Carcinoma Of The Ovary, Hypercalcemic Type} A rare, highly aggressive ovarian tumor in young patients, often associated with alterations in SMARCA4 and sometimes presenting with hypercalcemia.

\medterm{Small Cell Lung Cancer} An aggressive lung cancer strongly associated with smoking and early spread. It often responds initially to chemotherapy or radiation, but recurrence is common and treatment depends on the stage and overall health.

\medterm{Small Eye} An eye that is smaller than usual because it did not develop fully before birth. Vision may be reduced, especially when other eye structures are also affected, and assessment can identify associated conditions.

\medterm{Small Fiber Neuropathy} Damage to tiny nerves that carry pain, temperature, and automatic body signals. It can cause burning feet, tingling, poor temperature sensation, abnormal sweating, digestive symptoms, or heart-rate changes; diabetes is one possible cause.

\medterm{Small For Gestational Age} Birth weight or size below the expected range for gestational age, caused by constitutional small size, placental insufficiency, infection, genetic disease, or maternal factors.

\medterm{Small Intestinal Bacterial Overgrowth} An abnormal increase or change in bacteria within the small intestine, often causing bloating, abdominal
pain, nausea, diarrhea, fullness, weight loss, or nutrient deficiency. Causes include impaired intestinal movement,
structural problems, surgery, and some systemic diseases; persistent symptoms require medical evaluation.

\medterm{Small Intestinal Bacterial Overgrowth (Small Intestinal Bacterial Overgrowth)} Alternate index label; see \textit{Small Intestinal Bacterial Overgrowth}.

\medterm{Small Intestinal Ischemia And Infarction} Inadequate blood supply to the small intestine that can progress to irreversible tissue death, perforation, sepsis, and shock.

\medterm{Small Intestine} The long, coiled part of the digestive tract between the stomach and large intestine. It receives bile and pancreatic juices, completes digestion, and absorbs most nutrients and much of the body's water.

\medterm{Small Intestine Aspirate And Culture} Collection and culture of fluid aspirated from the small intestine to evaluate suspected bacterial overgrowth or selected infections; the procedure is invasive and results require clinical correlation.

\medterm{Small Intestine Prolapse} Alternate index label; see \textit{Small bowel prolapse (enterocele)}.

\medterm{Small Round Blue Cell Tumor} A tumor whose cells look small, round, and blue under a microscope. This appearance occurs in several cancers, especially in children, so additional tissue tests are needed to identify the exact type.

\medterm{Small Vessel Disease} Disease affecting small blood vessels, which can impair blood flow and contribute to brain, heart, kidney, or other organ problems.

\synonyms
Coronary Microvascular Disease

\medterm{Small-Bowel Transplantation} Surgery replacing a severely diseased or nonfunctioning small intestine with donor bowel. It may help people unable to receive enough nutrition through the gut but requires lifelong rejection-preventing medicines and monitoring.
\medterm{Small-Bowel Volvulus} Twisting of the small intestine that can block its contents and cut off its blood supply. It may cause severe abdominal pain, vomiting, swelling, and inability to pass stool or gas and requires urgent treatment.

\medterm{Small-Cell Carcinoma} A fast-growing cancer, most often starting in the lung and strongly linked with smoking. It may shrink quickly with chemotherapy or radiation but often returns or spreads early, including to the brain.
\medterm{Smallpox} A severe contagious disease caused by variola virus, producing fever and rash; global eradication was certified in 1980.

\medterm{SMART} A way to set goals that are specific, measurable, achievable, relevant, and time-limited. SMART goals state exactly what is wanted, how progress will be checked, and when the goal should be reached.

\synonyms
SMART goals

\medterm{Smashed Fingers} A smashed finger may injure skin, nail, bone, tendon, or nerve; deformity, severe pain, numbness, impaired movement, or blood under the nail may require urgent evaluation.

\medterm{Smear Of Duodenal Fluid Aspirate} A smear of duodenal fluid aspirate examines fluid from the first part of the small intestine for selected parasites or other microscopic findings.

\medterm{Smegma} A collection of shed cells, oils, and moisture around the genital folds, which may promote irritation if hygiene is poor.
\medterm{Smell} The sense of detecting airborne chemicals through receptors in the nasal lining and processing the signals in the brain. Reduced or lost smell can follow infection, nasal disease, head injury, medicines, or neurologic disease.
\medterm{Smell - Impaired} Reduced or absent ability to detect odors, called hyposmia or anosmia, caused by nasal obstruction, viral injury, head trauma, neurodegeneration, medications, or other disease.

\medterm{SMILE} A laser eye operation that corrects short-sightedness and some short-sighted astigmatism. The laser creates a thin piece of tissue inside the cornea, which is removed through a small opening without making a larger corneal flap.

\medterm{Smile Design} Planning smile based on facial features, gum display, tooth proportions. Digital smile design software available.
\medterm{Smith-Lemli-Opitz Syndrome} A rare inherited disorder affecting cholesterol production before birth. It may cause growth and developmental problems, distinctive facial features, heart or limb abnormalities, feeding difficulty, and genital differences.

\medterm{Smith-Magenis Syndrome} A genetic neurodevelopmental disorder, usually involving deletion of 17p11.2 or an RAI1
variant. It causes developmental delay, intellectual disability, distinctive facial
features, sleep disturbance, behavioral dysregulation, and variable congenital
abnormalities.

\medterm{Smoke Evacuator} A device that captures and filters smoke produced during surgery by electrical, laser, or ultrasonic instruments. Removing the plume reduces staff and patient exposure to particles, irritating gases, and potentially infectious material.

\medterm{Smoke Inhalation} Injury from breathing smoke, hot gases, or combustion products, potentially causing airway burns, poisoning, and lung injury.
\medterm{Smoke Inhalation Injury} Pulmonary injury caused by inhalation of toxic products of combustion.

\medterm{Smoke Secondhand (Secondhand Smoke)} Alternate index label for Secondhand Smoke.

\medterm{Smokeless Tobacco} Smokeless tobacco delivers nicotine and causes dependence, oral lesions, gum recession, tooth loss, cardiovascular risk, and cancers of the mouth, throat, and pancreas. Quitting reduces harm, and counseling plus approved medicines improve cessation success.

\synonyms
Tobacco Chewing (Smokeless Tobacco)

\medterm{Smokers Keratosis (Leukoplakia)} Alternate index label for Leukoplakia.

\medterm{Smokers Lung (Smoker's Lung: Pathology Photo Essay)} Alternate index label for Smoker's Lung: Pathology Photo Essay.

\medterm{Smoking} Inhalation of smoke from burning tobacco or other substances, associated with addiction, cancer, cardiovascular disease, and lung disease.
\medterm{Smoking And Asthma} Smoking irritates and inflames airways, worsens asthma control, increases attacks, and reduces treatment response; stopping smoking and avoiding secondhand smoke improve outcomes.

\medterm{Smoking And COPD} Smoking is the major preventable cause of COPD and accelerates lung-function loss; quitting reduces progression and exacerbations even after COPD is established.

\medterm{Smoking And Surgery} Smoking before or after an operation increases the risk of breathing problems, wound infection, poor healing, blood clots, heart complications, and failure of repaired tissue. Stopping as early as possible improves recovery, even if surgery is soon.

\medterm{Smoking Cessation} The process of stopping tobacco or nicotine use. It may involve counseling, behavioral support, medicines, or nicotine
replacement.

\medterm{Smoking Cessation Medications} Medicines that reduce nicotine withdrawal and cravings, including nicotine replacement, bupropion, and varenicline, used with behavioral support and individualized for health risks and preferences.

\medterm{Smoldering Multiple Myeloma} An asymptomatic plasma-cell disorder with increased clonal plasma cells or monoclonal protein but no myeloma-defining organ damage; it requires risk-based monitoring for progression to active multiple myeloma.

\synonyms
Voluntary muscles

\medterm{Smooth Muscle} Muscle that works automatically in the walls of hollow organs, blood vessels, and airways. It controls processes such as bowel movement, bladder emptying, blood-vessel width, airway narrowing, and childbirth contractions.
\medterm{Smothering} Mechanical obstruction of the external nose and mouth that prevents effective
ventilation. It may be caused by a covering, compression, or an obstructing soft
material; findings are often nonspecific and must be interpreted with the scene,
history, and complete examination.

\medterm{Snake Bite} Snake bites can cause puncture wounds, tissue injury, bleeding, paralysis, or venom-related effects on clotting, nerves, muscles, kidneys, or the heart. Call emergency services, keep the person still, and do not cut, suck, ice, or tourniquet the wound.

\synonyms
Bite Snake (Snake Bite)

\medterm{Snake Bites} Wounds caused by a snake's teeth; venomous bites may cause pain, swelling, bleeding, tissue damage, paralysis, abnormal clotting, or shock. Seek emergency care immediately, keep the person still, and do not cut, suck, ice, or tourniquet the wound.

\medterm{Snake-Bite} A wound from a snake that may cause local injury, envenomation, infection, or systemic toxicity.
\medterm{Snapping Finger} Trigger finger, painful catching or locking of a finger caused by stenosing inflammation around a flexor tendon.
\medterm{Snapping Hip Syndrome} A clicking or snapping sensation in the hip, sometimes with pain, caused by a tendon moving over bone or by a problem inside the joint. Activity adjustment and exercises often help; persistent pain may need imaging or other treatment.

\medterm{Sneddon Syndrome} A rare disorder in which small arteries become blocked, causing a net-like purplish skin pattern and repeated strokes or stroke-like events. Headache, high blood pressure, memory problems, or heart-valve changes may also occur.

\medterm{Sneezing} A forceful reflex expulsion of air through the nose and mouth, often triggered by nasal irritation.
\medterm{Snellen Chart} A chart of letters or symbols used to measure visual acuity at a specified distance.
\medterm{Snoring} Noisy breathing during sleep caused by vibration of relaxed or narrowed tissues in the throat. It can be harmless, but loud frequent snoring with pauses in breathing, choking, morning headaches, or daytime sleepiness may indicate sleep apnoea.
\medterm{Snoring And Sleep Apnea} Snoring is the sound produced by vibrating soft tissues in the upper airway during
sleep, often resulting from partial airway obstruction. Sleep apnea is a more serious
condition characterized by repeated episodes of complete or partial airway collapse
during sleep, leading to intermittent hypoxia, frequent arousal, and fragmented sleep.


\medterm{Snort} A snort is a forceful nasal or respiratory sound; it may occur with congestion, sleep-disordered breathing, irritation, or inhalation of a substance.

\medterm{Snot} Alternate informal term; see \textit{Mucus}.

\medterm{Snow Blindness} A painful temporary burn of the eye’s clear front surface caused by intense ultraviolet light, often reflected from snow, ice, water, or sand. Symptoms may include gritty pain, tearing, redness, light sensitivity, and blurred vision.
\medterm{SNRIs} A class of antidepressants that increase the activity of serotonin and norepinephrine, chemical messengers between nerve cells. They treat depression, some anxiety disorders, and certain nerve or long-lasting pain conditions.

\medterm{Snuff} Snuff is finely ground tobacco inhaled through the nose or placed in the mouth and can cause nicotine dependence, oral disease, and cardiovascular harm.

\medterm{Snuffles} A runny or congested nose; in newborns, persistent snuffles can be a sign of congenital syphilis.
\medterm{Social And Psychological Causes Of Alcoholism} Alcohol-use disorder reflects interactions among genetics, brain reward pathways, stress, trauma, mental illness, social environment, and alcohol availability—not a moral failure. Treatment may include withdrawal care, counseling, peer support, and medicines that reduce drinking or relapse.

\medterm{Social Anxiety Disorder} Intense, persistent fear of being judged, embarrassed, or rejected in social or performance situations. It can lead to avoidance, physical anxiety symptoms, and difficulty with school, work, relationships, or daily activities.

\medterm{Social Anxiety Disorder (Social Phobia)} Alternate index label for Social Anxiety Disorder.

\synonyms
Phobia, Social

\medterm{Social Isolation} Limited contact or interaction with other people. Persistent isolation may adversely affect mental and
physical health.

\medterm{Social Worker} A professional who helps patients and families manage practical, emotional, and social problems related to illness. They may arrange community resources, support discharge planning, assist with benefits, and provide counseling or referrals.

\medterm{Sociophobia} An older term for social anxiety disorder, involving excessive fear of scrutiny, embarrassment, humiliation, or negative evaluation in social or performance situations.

\medterm{Sodium} A mineral and body salt helping control fluid balance, nerves, muscles, and blood pressure; abnormal levels can be harmful.
\medterm{Sodium Bicarbonate Therapy} Treatment using sodium bicarbonate to make overly acidic blood less acidic in selected conditions, such as some cases of kidney disease. It can increase sodium and fluid load and requires monitoring.

\medterm{Sodium Bisulfate Poisoning} Poisoning from sodium bisulfate, an acidic chemical used in some cleaners and pool products. It can burn the skin, eyes, mouth, and digestive tract; fumes may irritate breathing passages, and ingestion requires urgent advice.

\medterm{Sodium Blood Test} A blood test measuring serum sodium, a major extracellular electrolyte involved in water balance and nerve and muscle function; results are interpreted with fluid status and other electrolytes.

\medterm{Sodium Carbonate Poisoning} Poisoning from sodium carbonate, an alkaline chemical found in some cleaning products. It may irritate or burn the eyes, skin, mouth, and digestive tract; do not induce vomiting after swallowing it.

\medterm{Sodium Chloride} The salt sodium chloride; sterile solutions are used for fluid replacement, irrigation, and drug dilution.
\medterm{Sodium Cromoglycate} A mast-cell stabilizer used in selected allergy and asthma treatments; also called cromolyn.
\medterm{Sodium Diatrizoate} A water-soluble iodinated contrast agent used for selected radiographic examinations. It can cause allergic-like reactions, kidney stress, or aspiration complications, so its use depends on the procedure and patient risks.
\medterm{Sodium Hydroxide Poisoning} Severe poisoning from sodium hydroxide, a strong alkali in some drain and oven cleaners. It can rapidly burn the mouth, throat, esophagus, stomach, skin, and eyes; emergency assessment is required and vomiting must not be induced.

\medterm{Sodium Hypochlorite} The active chemical in many bleach and disinfectant products. It kills many germs when used correctly, but concentrated liquid can burn skin, eyes, and the digestive tract; mixing it with acids or ammonia can release dangerous gas.
\medterm{Sodium Hypochlorite Poisoning} Poisoning from bleach containing sodium hypochlorite. Swallowing may irritate or burn the mouth and digestive tract; mixing bleach with acids or ammonia can release dangerous gases causing cough, chest pain, or breathing failure.

\medterm{Sodium In Diet} Dietary sodium helps regulate fluid balance and nerve and muscle function, but excess intake can raise blood pressure; processed foods are a major source and intake should follow individual health advice.

\medterm{Sodium Low Levels In The Blood (Hyponatremia)} Alternate index label for Hyponatremia.

\medterm{Sodium Polystyrene Sulfonate} A resin that binds potassium in the intestine and has been used to lower high blood potassium. It can cause constipation and serious intestinal injury, so its use requires careful medical assessment.

\medterm{Sodium Urine Test} A urine test measuring sodium excretion, used with blood tests and clinical findings to assess volume status, kidney handling of sodium, and causes of hyponatremia or hypertension.

\medterm{Sodium Valproate} An anticonvulsant and mood stabilizer; it carries important pregnancy and reproductive risks and requires specialist prescribing.
\medterm{Soft Markers For Aneuploidy} Ultrasound findings during pregnancy that are common in healthy fetuses but occur somewhat more often when there is an extra or missing chromosome. They do not diagnose a chromosome condition and are interpreted with other screening results.

\medterm{Soft Tissue Augmentation} Restoration or enhancement of volume and contour in skin or underlying soft tissue using fillers, fat transfer, implants, or other reconstructive techniques.

\medterm{Soft Tissue Injury} Damage to muscles, tendons, ligaments, fascia, skin, or other non-bony tissues, often caused by
trauma or overuse.

\medterm{Soft Tissue Rheumatism} An older broad term for pain or inflammation affecting muscles, tendons, ligaments, or bursae. Modern diagnosis usually names the specific condition, such as tendinitis, bursitis, or a muscle strain.

\medterm{Soft Tissue Sarcoma} A group of malignant tumors arising in muscle, fat, fibrous tissue, blood vessels, or
other connective tissues. Treatment depends on the subtype, location, stage, and relationship to nearby structures.

\synonyms
Cancer, Soft Tissue Sarcoma
Sarcoma, Soft Tissue

\medterm{Soft Tissue Tumors} Growths arising in muscle, fat, fibrous tissue, blood vessels, nerves, or other supporting tissues. Many are harmless, but some are cancerous; a growing, deep, firm, painful, or persistent lump needs assessment.

\medterm{Soft-Tissue Calcinosis} Buildup of calcium deposits in skin, muscles, tendons, or other soft tissues. The deposits may form hard lumps, cause pain or limited movement, break through the skin, or occur because of kidney disease, abnormal calcium levels, inflammation, or tissue damage.

\medterm{Soil Contamination} The presence of harmful substances in soil, such as heavy metals, pesticides, fuel, or industrial waste. People may be exposed through food, water, dust, or skin; testing, removal, covering, or treatment of the soil can reduce risk.

\medterm{Solar Keratosis} Alternate index label; see \textit{Actinic keratosis}.

\medterm{Solar Keratosis (Actinic Keratosis)} Alternate index label for Actinic Keratosis.

\medterm{Solar Lentigines} Alternate index label; see \textit{Age spots (liver spots)}.

\medterm{Solar Plexus} A network of involuntary nerves in the upper abdomen, centered around the celiac plexus. It helps regulate abdominal organs, and irritation or injury may contribute to upper-abdominal pain or altered organ function.
\medterm{Solar Urticaria} A rare allergic-type reaction to sunlight or other light that causes itchy red raised patches within minutes on exposed skin. The rash usually improves after leaving the light, but severe widespread symptoms need medical care.

\medterm{Solder Poisoning} Poisoning from swallowing solder or inhaling fumes or dust during heating. Products may contain lead, tin, antimony, or other metals and can cause stomach symptoms, anemia, nerve injury, or kidney damage; exposure needs evaluation.

\medterm{Sole Sweating Excessive (Hyperhidrosis)} Alternate index label for Hyperhidrosis.

\medterm{Soleus Muscle} A deep calf muscle that joins the Achilles tendon and helps point the foot downward. It is important for standing, walking, and running; injury or overuse can cause calf pain and weakness.

\medterm{Solitary Fibrous Tumor} A solitary fibrous tumor is a usually slow-growing growth of connective tissue that can arise in the pleura, meninges, or other sites and may rarely behave aggressively.

\medterm{Solitary Pulmonary Nodule} A single small spot in the lung surrounded by otherwise normal lung tissue. It may be a scar, infection, harmless growth, or early cancer; its size, appearance, and change over time determine whether follow-up or biopsy is needed.


\medterm{Solution} A homogeneous mixture in which one or more substances are dissolved in a solvent; medical solutions include saline and antiseptics.
\medterm{Solvent Abuse (Misuse)} Harmful inhalation or ingestion of volatile solvents, which can cause neurologic, cardiac, liver, kidney, or respiratory injury.
\medterm{Solvent Exposure} Contact with organic solvents by inhalation, ingestion, or skin absorption. Acute exposure may cause
irritation or central nervous system effects; repeated exposure may damage organs.

\medterm{Somatic} Relating to the body, especially the body wall and musculoskeletal system, or to physical symptoms.
\medterm{Somatic Symptom And Related Disorders} A group of mental disorders involving distressing physical symptoms, excessive health-related thoughts or behaviors, illness fears, functional neurologic symptoms, or intentionally produced illness. Examples include somatic symptom disorder, illness anxiety disorder, functional neurologic disorder, and factitious disorder. Malingering is separate because it involves intentional symptom production for an external benefit.
\medterm{Somatic Mutation} A genetic change acquired by a body cell during life rather than inherited from a parent. It is passed to descendants of that cell and may contribute to cancer, but it usually is not passed to children.

\medterm{Somatic Symptom Disorder} A mental disorder in which one or more physical symptoms cause significant distress or disruption of daily life, together with excessive thoughts, feelings, or behaviors related to the symptoms. The symptoms are real and may occur with or without an identified medical condition; their severity is not determined solely by whether a medical cause is found.

\medterm{Somatic Variant} A change in DNA acquired by a body cell after conception rather than inherited from a parent. It usually cannot be passed to children but may affect only some tissues or contribute to cancer.

\medterm{Somatization} The experience or expression of psychological distress through physical symptoms such as pain, fatigue, or stomach problems. The symptoms are real and not intentional; assessment should consider both physical and psychological causes.

\medterm{Somatostatin} A hormone made in the brain, pancreas, and digestive tract that slows growth hormone, insulin, glucagon, and digestive-hormone release. It also reduces movement and secretion within parts of the digestive system.
\medterm{Somatostatin Analog} A medicine that imitates somatostatin, a hormone that slows the release of several hormones and digestive secretions. It is used for selected hormone-producing tumors, acromegaly, severe diarrhea, and some bleeding conditions.

\medterm{Somatostatinoma} A rare hormone-producing tumor usually arising in the pancreas or small intestine. It may cause diabetes, diarrhea or greasy stools, gallstones, weight loss, or no symptoms; treatment depends on whether it has spread and may include surgery or cancer medicines.

\medterm{Somatotropic Pituitary Adenoma} A usually noncancerous pituitary tumor that makes too much growth hormone. In adults it causes acromegaly, with enlargement of the hands, feet, jaw, or other tissues; in children it can cause excessive height before growth ends.

\medterm{Somatotropin} Growth hormone, produced by the pituitary gland. It supports childhood growth and influences protein, fat, and carbohydrate metabolism; excess or deficiency can cause characteristic growth and metabolic disorders.
\medterm{Somatotype} A descriptive body-build classification; it is not a disease diagnosis.
\medterm{Somnambulism} Sleepwalking or performing other activities while only partly awake, usually during deep sleep. Episodes are more common in children and may be triggered by lack of sleep, stress, illness, or medicines; safety measures are important, especially if wandering or injury occurs.
\medterm{Somniloquy (Sleep-Talking)} Talking, laughing, or making sounds during sleep without knowing it. It is usually harmless, but frequent or disruptive episodes may occur with sleepwalking, nightmares, breathing problems during sleep, medicines, or other sleep disorders.

\synonyms
Sleep talking

\medterm{Somnolence} Unusual sleepiness or a tendency to fall asleep when a person should be awake. It may result from too little sleep, medicines, alcohol or drugs, sleep disorders, metabolic illness, or neurologic disease.

\medterm{Somnolent} Somnolent means unusually sleepy or drowsy, which may result from sleep deprivation, medicines, intoxication, metabolic illness, or neurologic disease.

\medterm{Sonographer} A healthcare professional trained to obtain diagnostic ultrasound images. They position the patient, operate the scanner, record appropriate views and measurements, and document findings for interpretation by a qualified clinician.
\medterm{Soporifics} Medicines or substances that cause sleepiness or promote sleep. They may impair alertness, coordination, and breathing, especially when combined with alcohol or other sedatives.
\medterm{Sorbitol} A sugar alcohol used as a sweetener and in some medicines. Larger amounts draw water into the bowel and act as a laxative, but may cause gas, cramps, diarrhea, or dehydration.
\medterm{Sore} Painful or tender, or describing an open area on skin or a moist body lining. A sore may result from injury, infection, inflammation, pressure, allergy, or another medical condition.
\medterm{Sore Throat} Pain, scratchiness, or irritation in the throat. Common causes include viral infections, bacterial infections, allergies, dry air, smoke, voice strain, and acid reflux; persistent, severe, or breathing-related symptoms need medical assessment.
\medterm{Sore Throat (Pharyngitis)} Alternate index label for Pharyngitis.


\medterm{Sores Canker (Canker Sores)} Alternate index label for Canker Sores.


\medterm{Sotalol} A medicine with beta-blocking and class III antiarrhythmic effects, used for selected cardiac rhythm disorders.
\medterm{Sotos Syndrome} A genetic disorder that causes unusually rapid growth in childhood, a large head, distinctive facial features, and delayed development or learning difficulties. Coordination, behavior, seizures, heart problems, or kidney problems may also occur. Growth usually slows with age. A change in the \textit{NSD1} gene causes most cases.

\synonyms
Cerebral Gigantism; NSD1-Related Overgrowth Syndrome

\medterm{Sound} Mechanical energy transmitted as waves and perceived by hearing; clinical sound assessment includes speech and respiratory or cardiac sounds.
\medterm{Source Control} A procedure or intervention that eliminates, drains, debrides, or otherwise controls an
anatomic focus of infection or contamination. Examples include drainage of an abscess,
removal of infected tissue or device, and repair of a perforation.

\medterm{Soy} Soy is a legume providing protein and isoflavones; it is nutritious for many people but can cause allergy and may interact with some medicines or thyroid treatment timing.

\medterm{Soy Protein Intolerance} Adverse gastrointestinal reaction to soy proteins, most often seen in infants and young
children. It may cause vomiting, diarrhoea, abdominal pain, blood or mucus in the stool,
feeding difficulty, or poor growth.

\medterm{Soya} Soybean and products made from it; soy is a food source and can cause allergy in susceptible people.
\medterm{Space Medicine} The medical specialty concerned with human health and performance during aviation and space travel. It addresses altered gravity, radiation, isolation, motion sickness, sleep, emergencies, exercise, and returning safely to Earth.
\medterm{Spacer} A chamber attached to a metered-dose inhaler that holds the released medicine briefly, making it easier to breathe into the lungs and reducing medicine left in the mouth or throat.
\synonyms
Valved holding chamber

\medterm{Spanish Fly} A historical name for cantharidin-containing blister beetles; ingestion is poisonous and it is not a safe aphrodisiac.
\medterm{Spasm} An involuntary, often painful tightening of one muscle or a group of muscles. Spasms may result from fatigue, dehydration, nerve disease, injury, medicines, inflammation, or problems affecting the brain or spinal cord.
\medterm{Spasmodic Dysphonia} Spasmodic dysphonia is a focal laryngeal dystonia causing involuntary voice-box spasms and strained, breathy, or interrupted speech.

\medterm{Spasmodic Torticollis} A form of neck dystonia with involuntary muscle contractions causing twisting, tilting, or abnormal neck posture.
\medterm{Spasmus Nutans} An uncommon early-childhood syndrome of rapid shimmering eye movements, head nodding, and torticollis that usually resolves but requires evaluation to exclude neurologic or ocular disease.

\medterm{Spastic} Describes muscles that are unusually tight and resist being stretched, especially when moved quickly. It may cause stiffness, jerking, abnormal posture, and difficulty controlling movement after nervous-system injury.
\medterm{Spastic Colon (Irritable Bowel Syndrome)} Alternate index label for Irritable Bowel Syndrome.

\medterm{Spasticity} Velocity-dependent increase in muscle tone caused by an upper motor-neuron lesion, producing stiffness, spasms, abnormal reflexes, and restricted movement.

\medterm{Spasticity Management} Assessment and treatment of abnormal muscle stiffness and spasms. Measures may include
stretching, therapy, positioning, medicines, injections, or surgery.

\medterm{Spatial Resolution} The ability of an imaging system to show two nearby small objects as separate. Better resolution shows finer detail, but the result depends on the scanner, movement, body part, scan settings, and the type of imaging used.

\medterm{Special Care Baby Unit} A ward within a hospital that provides specialist care for newborn babies who need
additional support but not intensive care.

\medterm{Special Stain} A histochemical stain that highlights a particular tissue component, organism, or
deposited substance not adequately demonstrated by routine hematoxylin and eosin.
Examples include periodic acid-Schiff, Ziehl-Neelsen, Congo red, and Gomori methenamine
silver stains.

\medterm{Specific Absorbed Fraction} The portion of energy released by radiation in one body area that is absorbed by a particular organ or tissue. It helps calculate the radiation dose during nuclear-medicine treatment or testing.

\medterm{Specific Activity} The amount of radioactivity in a given amount of material. It matters when preparing radioactive medicines because enough activity is needed for imaging or treatment while limiting the amount of substance and unnecessary exposure.

\medterm{Specific Gravity} A comparison of urine density with water that gives an estimate of how concentrated the urine is. It may rise with dehydration, sugar, or protein in urine and fall when the kidneys cannot concentrate urine or when excess fluid is present.

\medterm{Specific Learning Disorder} A developmental condition causing lasting difficulty learning or using reading, writing, or mathematics skills despite appropriate teaching. The difficulty interferes with school, work, or daily activities and is not explained only by poor instruction, vision, hearing, or intellectual disability.

\medterm{Specific Phobia} A disorder involving marked, persistent fear or anxiety about a particular object or situation.
The fear is excessive relative to the actual danger, leads to avoidance or intense distress,
and interferes with daily functioning.

\medterm{Specific Reading Disability} Alternate index label; see \textit{Dyslexia}.

\medterm{Specificity} A test's ability to correctly identify people who do not have the condition being tested.
\medterm{Specimen Adequacy} The degree to which a collected clinical specimen is sufficient in quantity, correctly
identified, properly transported, and representative of the site being tested.
Inadequacy can cause invalid, misleading, or falsely negative results.

\medterm{Specimen Container} A clean, suitable container used to collect, identify, protect, and transport a patient sample for laboratory testing. The container depends on the sample and test; it should be labeled correctly, closed securely, and kept free from contamination.

\medterm{Specimen Transport Medium} A special liquid or gel that keeps germs or cells suitable for testing while a patient sample travels to the laboratory. The medium, temperature, container, and delivery time depend on the sample and suspected infection.

\medterm{SPECT Scan (Single Photon Emission CT)} A nuclear-medicine scan that uses a small radioactive tracer and a rotating camera to create three-dimensional pictures of blood flow or organ function. It can examine the heart, brain, bones, kidneys, or other organs.

\synonyms
Single Photon Emission CT

\medterm{Spectral Analysis} Examination of the different energy levels within radiation or another signal. In medical imaging, it can separate useful information from background scatter and help identify materials or improve image interpretation.

\medterm{Speculum} An instrument used to gently separate or hold open the walls of a body passage so that the area can be visualized, examined, sampled, or treated. Common types include vaginal, nasal, ear, and rectal specula; design and size depend on the body site and intended procedure.
\medterm{Speculum, Ear} A small cone-shaped tip fitted to an otoscope to direct light into the external ear canal and permit examination of the canal and tympanic membrane. Disposable and reusable types are available in different sizes.

\medterm{Speech} The production of spoken language using breathing, voice, and movements of the tongue, lips, palate, and other structures. Hearing, brain function, teeth, illness, and dental appliances can affect speech.

\medterm{Speech Audiometry} Hearing testing that measures detection or recognition of speech at different loudness levels. It commonly includes
speech reception threshold and word-recognition testing, complementing pure-tone audiometry.

\synonyms
Speech testing; speech audiometric testing.

\medterm{Speech Delay} Speech or language development that is slower than expected for a child's age. Causes may include hearing loss, developmental differences, autism, neurologic disease, or limited language access; loss of previously learned skills needs prompt assessment.

\medterm{Speech And Language Delay} Slower-than-expected development of speech sounds, vocabulary, understanding, or use of language in a child. Causes may include hearing loss, developmental disorders, autism, neurologic disease, or other medical and environmental factors.

\medterm{Speech Discrimination} A hearing test that measures how well a person understands spoken words at a comfortable loudness. The result helps assess hearing function beyond the ability to detect pure tones.

\medterm{Speech Disorders} Conditions that affect making speech sounds, speaking smoothly, using the voice, or controlling the movements needed for speech. They may follow hearing loss, developmental differences, brain injury, stroke, nerve disease, or other conditions and can often be assessed by a speech therapist.

\medterm{Speech Impairment In Adults} Difficulty producing or understanding speech caused by aphasia, dysarthria, apraxia, hearing loss, neurologic disease, structural problems, or cognitive disorder; sudden onset requires stroke assessment.

\medterm{Speech Therapy} Assessment and treatment to help with speaking, understanding language, voice, communication, or swallowing. A speech and language therapist may use exercises, practice, communication aids, and advice for the person and family.
\medterm{Speech-Language  Pathologist} A clinician who evaluates and treats disorders of speech, language, voice, communication, cognition, and swallowing across childhood, adulthood, and older age.

\medterm{Speech-Language Pathology Assessment} An evaluation of speech, language, voice, communication, swallowing, and sometimes related cognitive skills. It combines interview, clinical observation, and standardized or functional tests to guide therapy.


\medterm{Sperm} Mature male reproductive cells that can fertilize an egg. Semen contains sperm plus fluids from the reproductive glands; sperm are produced in the testes.
\medterm{Spermarche} The onset of sperm production and the first ejaculation, marking an important stage of male puberty.

\medterm{Spermatic} Relating to sperm or structures that produce, transport, or support sperm. Spermatic structures include the testes, ducts, blood vessels, nerves, and tissues involved in male reproductive function.
\medterm{Spermatic Cord} A cord-like structure carrying the vas deferens, testicular blood vessels, nerves, lymph vessels, and supporting tissue. It extends through the groin to the testis and may be affected by hernia, twisting, or infection.
\medterm{Spermatocele} A usually harmless, fluid-filled cyst beside or within the epididymis, the coiled tube behind the testicle. It often causes no symptoms but may feel like a lump and should be checked if newly noticed or painful.

\medterm{Sperm Retrieval} A procedure that obtains sperm from semen, the epididymis, or testicular tissue for fertility treatment or laboratory testing. The method depends on whether sperm are present in the semen and on the cause of infertility.

\synonyms
Cyst, Spermatic

\medterm{Spermatogenesis} The production and maturation of sperm in the seminiferous tubules of the testes. It depends on germ cells, Sertoli cells, hormones, and suitable testicular temperature.
\medterm{Spermatogonium} An early male reproductive cell in the testis that divides and develops into primary spermatocytes, beginning sperm production.

\medterm{Spermatorrhea} Involuntary leakage or discharge of semen, sometimes occurring during sleep or without orgasm. Occasional nocturnal emissions can be normal; persistent discharge, pain, blood, or urinary symptoms needs assessment.

\medterm{Spermatozoa} Mature male reproductive cells, also called sperm. Each normally carries half the usual chromosome number and has a head, middle section, and tail that helps it move toward an egg.
\medterm{Spermatozoon} A single mature male reproductive cell, or sperm cell. The singular form of spermatozoa: one mature male reproductive cell capable of carrying genetic material to an egg.
\medterm{Spermicide} A contraceptive substance that kills or immobilizes sperm, often used with condoms or other barrier methods.
\medterm{Sphenoid} Relating to the sphenoid bone, a central bone at the base of the skull. It helps form the eye sockets and skull base and contains the hollow spaces called the sphenoid sinuses.
\medterm{Sphenoid Bone} A central bone at the base of the skull that helps form the eye sockets and skull base. It contains the sphenoid sinuses and a hollow that holds the pituitary gland.

\medterm{Sphenoid Sinusitis} Inflammation or infection of the sphenoid sinus, which lies deep behind the nasal cavity. It may cause
deep headache, facial pain, or visual symptoms and can rarely spread to nearby nerves or the brain;
persistent or severe symptoms require medical assessment.

\medterm{Sphenoidal Sinus} A paired air-filled sinus within the sphenoid bone at the base of the skull. It lies near important nerves and blood
vessels, so inflammation or surgery in this area requires careful assessment.

\medterm{Spherocytes} Red blood cells that are unusually round and dense instead of having their normal flattened shape. They may occur when the cell membrane is damaged or removed, especially in hereditary spherocytosis or immune red-cell destruction.

\medterm{Spherocytosis (Hereditary, HS)} Alternate index label for Hereditary, HS.

\medterm{Sphincter} A ring-shaped muscle that opens and closes a body passage. Sphincters help control urine, stool, food, bile, and other fluids; examples include the muscles around the anus, bladder outlet, and stomach outlet.

\medterm{Sphincter Of Oddi} A muscular valve controlling the release of bile and pancreatic juice into the first part of the small intestine. It opens during digestion and abnormal blockage or tightening can cause pain or pancreatitis.

\medterm{Sphygmograph} An instrument or sensor that records the shape and timing of the pulse in an artery. The tracing can provide information about circulation and arterial stiffness.
\medterm{Sphygmomanometer} A device used to measure arterial blood pressure with an inflatable cuff and pressure-measuring system. Manual versions commonly use an aneroid or mercury manometer with a stethoscope, while electronic versions use sensors to analyze pressure oscillations and display the reading.

\synonyms
Blood Pressure Gauge
\medterm{Spider Angioma} A spider angioma is a central red blood vessel with radiating fine vessels that blanches with pressure and may occur normally or with liver or hormonal disease.

\medterm{Spider Bite} An injury caused by a spider, usually producing a small painful or itchy area. Some species can cause worsening tissue damage, severe muscle cramps, sweating, vomiting, breathing problems, or other serious effects; rapidly worsening symptoms need emergency care.

\medterm{Spider Bites (Black Widow And Brown Recluse)} Alternate index label for Black Widow and Brown Recluse.

\medterm{Spina Bifida} A birth defect in which the neural tube and the bones protecting the spinal cord do not close completely during early development. It can range from mild to severe and may cause weakness, loss of feeling, or problems controlling the bladder or bowel.

\medterm{Spina Bifida Meningocele} A form of spina bifida in which the protective coverings around the spinal cord bulge through an opening in the backbone and form a fluid-filled sac. The spinal cord is usually not inside the sac, so weakness or loss of bladder and bowel control may be less severe.

\medterm{Spina Bifida Myelomeningocele} The most severe common form of spina bifida, in which the spinal cord and its coverings bulge through an opening in the backbone. It can cause weakness or paralysis, loss of feeling, and difficulty controlling the bladder or bowel below the affected area.

\medterm{Spina Bifida Occulta} A usually mild form of spina bifida in which the vertebral arch does not fully close
without herniation of the spinal cord or meninges. It is often asymptomatic but may be associated with a skin dimple,
hair tuft, or other finding over the lower spine.

\medterm{Spinal Analgesia} Pain relief produced by injecting numbing medicine, sometimes with an opioid, into the fluid around the lower spinal nerves. It may cause temporary leg weakness, itching, difficulty urinating, low blood pressure, headache, or slowed breathing.

\medterm{Spinal And Bulbar Muscular Atrophy} An inherited disorder that gradually damages the nerve cells controlling movement. It usually affects adult men, causing weakness and wasting of the arms, legs, face, and throat, with muscle twitching, difficulty speaking or swallowing, and sometimes breast enlargement or reduced fertility. Treatment supports movement, swallowing, breathing, and daily activities.

\synonyms
Kennedy Disease; SBMA

\medterm{Spinal And Epidural Anesthesia} Spinal and epidural anesthesia deliver local anesthetic near spinal nerves to numb the lower body; spinal is a single injection, while epidural uses a catheter for adjustable dosing.

\medterm{Spinal Anesthesia} An injection of numbing medicine into the fluid around the lower spinal nerves to temporarily block pain and movement in the lower body. Possible effects include low blood pressure, difficulty urinating, headache, infection, bleeding, or rarely excessive numbness and breathing problems.

\synonyms
Spinal Anaesthesia

\medterm{Spinal Anesthesia For Cesarean Delivery} A regional anesthetic injected into the fluid around the lower spinal nerves to numb the lower body during a cesarean birth. The mother remains awake, while blood pressure, breathing, and the baby's condition are monitored.

\medterm{Spinal Angiography} X-ray imaging of blood vessels supplying the spinal cord after contrast is injected. It is used to locate vascular malformations, fistulas, aneurysms, or sources of bleeding and may guide treatment.

\medterm{Spinal Artery} An artery supplying the spinal cord and its coverings, including the anterior spinal artery and paired posterior spinal arteries; blockage can cause characteristic neurologic deficits.

\medterm{Spinal Block} A spinal block is injection of local anesthetic into cerebrospinal fluid in the lower back, producing temporary numbness and weakness for surgery or pain control.

\medterm{Spinal Canal} The spinal canal is the bony passage through the vertebral column containing the spinal cord, its coverings, nerve roots, and blood vessels.

\medterm{Spinal Cavity} The posterior body cavity formed by the vertebral column that contains and protects the spinal cord.

\medterm{Spinal Column} The vertebral column, consisting of vertebrae, discs, ligaments, and joints that support and protect the spinal cord.
\medterm{Spinal Cord} The cylindrical nervous tissue within the vertebral canal that carries signals between brain and body and mediates reflexes.
\medterm{Spinal Cord Abscess} A spinal cord or epidural abscess is a collection of infection near the spinal cord that can compress neural tissue and cause pain, fever, weakness, paralysis, or bladder problems.

\medterm{Spinal Cord Compression} Pressure on the spinal cord or its nerve roots, often from a tumor, infection, bleeding, fracture, or slipped disc. It can cause severe back pain, weakness, numbness, walking difficulty, or loss of bladder or bowel control and requires urgent medical assessment.

\medterm{Spinal Cord Disease} A disorder affecting the spinal cord, including compression, inflammation, vascular injury, infection, tumor, degeneration, or trauma, and potentially causing weakness, sensory loss, or autonomic dysfunction.

\medterm{Spinal Cord Injury} Damage to the spinal cord that can cause weakness, paralysis, altered sensation, or loss of bladder and bowel control. Severity is often graded from complete to incomplete using a standardized neurological examination.

\synonyms
SCI

\medterm{Spinal Cord Injury: Treatments And Rehabilitation} Spinal-cord injury can cause temporary or permanent weakness, sensory loss, autonomic dysfunction, and bowel or bladder problems. Acute care protects the spine and treats complications; rehabilitation addresses mobility, skin, breathing, function, equipment, sexuality, and community participation.

\medterm{Spinal Cord Neoplasms} Tumors arising in the spinal cord, its coverings, or nearby tissues. They may be noncancerous or cancerous and can cause pain, weakness, numbness, walking problems, or bladder and bowel changes by pressing on nerves.

\medterm{Spinal Cord Stimulation} An implanted treatment that delivers electrical impulses to the epidural space to alter pain signaling, used for selected chronic neuropathic or ischemic pain after careful assessment.

\medterm{Spinal Cord Trauma} Spinal cord trauma is injury to the cord from force such as fracture, dislocation, or penetrating damage, potentially causing weakness, sensory loss, autonomic dysfunction, or paralysis.

\medterm{Spinal Curvature} Alternate index label; see \textit{Scoliosis}.

\medterm{Spinal Disease} Any disorder affecting the vertebrae, discs, joints, ligaments, spinal cord, or nerve roots. It may cause back or neck pain, stiffness, weakness, numbness, walking difficulty, or changes in bladder or bowel control.
\medterm{Spinal Dysraphism} A group of birth defects caused by incomplete development or closure of the spine and spinal cord before birth. Some forms have an open sac on the back, while others are hidden and may later cause weakness, pain, or bladder and bowel problems.

\medterm{Spinal Epidural Abscess} A collection of pus between the covering of the spinal cord and the surrounding bone, usually caused by infection. Severe back pain, fever, weakness, numbness, difficulty walking, or loss of bladder or bowel control requires emergency assessment.

\medterm{Spinal Fluid Leak} Alternate index label; see \textit{CSF leak (Cerebrospinal fluid leak)}.

\medterm{Spinal Fracture} A fracture of one or more vertebrae caused by trauma, osteoporosis, tumor, or other bone-weakening disease.

\medterm{Spinal Fusion} Surgery that joins two or more spinal bones so they heal into one stable segment. Bone and often screws, rods, or cages are used; the operation may reduce movement and pain from selected conditions but can cause infection, nerve injury, or problems at nearby levels.

\medterm{Spinal Ganglion} A small group of nerve-cell bodies beside the spinal cord. It contains sensory nerve cells that carry information about touch, pain, temperature, and body position from the skin, muscles, and joints.

\medterm{Spinal Injury} Damage to the spinal cord, vertebrae, discs, ligaments, or nearby nerves caused by trauma or another disease. It may cause pain, weakness, numbness, paralysis, or bladder and bowel problems and requires urgent assessment when severe.


\medterm{Spinal Muscular Atrophy} An inherited disorder causing progressive loss of motor neurons and muscle weakness, often beginning in infancy
or childhood but sometimes appearing in adulthood. Severity varies; breathing, swallowing, and mobility
may be affected. Early diagnosis allows genetic counseling and disease-modifying care.

\medterm{Spinal Neoplasms} Tumors arising in the spinal cord, its coverings, vertebrae, or nearby tissues. They may be noncancerous or cancerous and can cause back pain, weakness, numbness, walking difficulty, or bladder and bowel changes when they press on nerves.
\medterm{Spinal Nerve} One of 31 pairs of nerves leaving the spinal cord through openings between the spinal bones. Each carries messages for sensation, movement, and automatic functions between the body and spinal cord.

\medterm{Spinal Precautions} Measures used after possible spinal injury to limit movement until the spine and spinal cord have been assessed. They may include careful positioning, support of the neck, and safe transfer; the exact approach depends on the injury and examination.

\medterm{Spinal Shock} A temporary loss of reflexes, muscle tone, and automatic functions below a spinal-cord injury. The affected limbs may initially be limp, and blood-pressure, bladder, bowel, or sexual problems can occur while the injury is assessed and treated.

\medterm{Spinal Stenosis} Narrowing of the passage containing the spinal cord or nerve roots. It may cause back pain, numbness, weakness, or leg pain with walking that improves when sitting or leaning forward.

\synonyms
Stenosis, Spinal

\medterm{Spinal Tap} Alternate index label; see \textit{Lumbar puncture}.

\medterm{Spinal Tumor} An abnormal growth in the spinal bones, spinal cord, its coverings, or nearby tissues. It may be noncancerous or cancerous; symptoms depend on its location and may include pain, weakness, numbness, or loss of bladder control.

\medterm{Spindle Pole} One end of the temporary structure formed during cell division. Fibers extending from each pole attach to duplicated chromosomes and pull the copies into the two new cells.

\medterm{Spindle Pole Body} A structure in fungal cells that organizes the fibers used to separate duplicated chromosomes during cell division. It performs a role similar to the centrosome in animal cells.

\medterm{Spindle-Cell Sarcoma} A group of cancers made of long, narrow cells under the microscope. They can arise from muscle, fibrous tissue, or nerve coverings; their behavior and treatment vary by the exact tumor type and location.

\medterm{Spine} The vertebral column, consisting of stacked vertebrae that support the body and protect the spinal cord.

\medterm{Spine And Spinal Cord, Diseases And Injuries Of} Diseases or injuries affecting the spinal bones, discs, ligaments, spinal cord, or nerve roots. They may cause pain, weakness, numbness, balance problems, deformity, or changes in bladder and bowel function.
\medterm{Spine Disorders} Conditions affecting the bones, disks, joints, ligaments, nerves, or spinal cord. They may cause back or neck pain, stiffness, weakness, numbness, balance problems, or changes in bladder and bowel control.

\medterm{Spine Injury} Damage to the bones, discs, ligaments, nerves, or spinal cord. It may cause back pain, weakness, numbness, loss of movement, or bladder and bowel problems and can require emergency care.
\medterm{Spinocerebellar Ataxia} Alternate index label; see \textit{Spinocerebellar Ataxias}.

\medterm{Spinocerebellar Ataxias} A group of rare inherited disorders that progressively damage the cerebellum and sometimes other parts of the nervous system. The cerebellum helps control balance, coordination, and precise movement. These disorders may cause unsteady walking, poor hand coordination, slurred or unclear speech, shaking, abnormal eye movements, muscle stiffness or weakness, and difficulty swallowing. The age of onset, symptoms, and rate of progression vary widely. Many types are inherited in an autosomal dominant pattern, but the inheritance pattern differs among types.

\medterm{Spinous Process} The backward-pointing bony projection of a vertebra that can be felt along the middle of the back. Muscles and ligaments attach to it and help support and move the spine.
\medterm{Spiradenocarcinoma} A rare malignant sweat-gland tumor that may arise in a long-standing spiradenoma, presenting as a rapidly enlarging, painful, or ulcerated nodule.

\medterm{Spiradenoma} A usually benign painful sweat-gland tumor presenting as a dermal nodule; multiple spiradenomas may occur in Brooke-Spiegler syndrome.


\medterm{Spiral Organ} The organ of Corti, a sensory structure in the cochlea containing hair cells that convert sound vibrations into nerve signals.

\medterm{Spiramycin} A macrolide antibiotic used for selected bacterial infections and sometimes during pregnancy for toxoplasmosis. Choice depends on the infection, local guidance, pregnancy stage, allergies, resistance, and other treatment needs.
\medterm{Spirillum Minus} Spirillum minus is a spiral Gram-negative bacterium that causes rat-bite fever, characterized by fever, rash, and migratory joint pain after rodent exposure; culture is difficult and diagnosis may be clinical or molecular.

\medterm{Spirit} A volatile liquid preparation or alcohol; the meaning depends on the medical, pharmaceutical, or general context.
\medterm{Spiritual Pain} Distress involving meaning, purpose, beliefs, identity, or connection, often addressed in palliative and supportive care.
\medterm{Spiritual Support} Support that helps a person address meaning, values, faith, hope, or existential distress during illness; it may be provided by chaplains, spiritual leaders, clinicians, or trusted people.

\medterm{Spirochaetaceae} A family of spiral-shaped bacteria that includes medically important spirochetes such as Borrelia, Treponema, and Leptospira.

\medterm{Spirochaetales Infection} Infection caused by spiral-shaped bacteria in the order Spirochaetales, including syphilis, Lyme disease, relapsing fever, and leptospirosis.

\synonyms
Spirochaetal Infections; Spirochetal Infections

\medterm{Spirochaete} Another spelling of spirochete, a thin, flexible, spiral-shaped bacterium. Important genera include Treponema, Borrelia, and Leptospira.
\medterm{Spirochete} A thin, spiral-shaped bacterium that moves with a corkscrew-like motion. Important examples cause syphilis, Lyme disease, relapsing fever, and leptospirosis; diagnosis depends on the organism and may use blood tests or other samples.

\medterm{Spirogram} A graph recording how much air a person breathes and how quickly it moves during a lung-function test. It helps identify patterns suggesting airway narrowing, restricted lung expansion, or poor effort.

\medterm{Spirometer} An instrument that measures airflow and the volume of air moved during specified breathing maneuvers, especially forced expiration. It does not by itself measure all lung volumes; residual volume and total lung capacity require additional pulmonary-function testing.
\medterm{Spirometry} A breathing test that measures how much air a person can blow out and how quickly. It helps identify narrowing of the airways, as in asthma or chronic obstructive pulmonary disease, and can monitor treatment. The person must follow instructions and blow forcefully several times; results are interpreted with age, height, and symptoms.

\medterm{Spironolactone} An aldosterone antagonist and potassium-sparing diuretic used for heart failure, hypertension, hyperaldosteronism, and other indications.
\medterm{Spit} Saliva or other fluid expelled from the mouth. Increased saliva or repeated spitting may occur with mouth irritation, nausea, reflux, difficulty swallowing, medication effects, or habit.

\medterm{Spitting} Forceful expulsion of saliva or other material from the mouth. New or excessive spitting may occur with oral disease, reflux, nausea, difficulty swallowing, medication effects, or habit.


\medterm{Spitz Nevus} A benign melanocytic nevus, often a pink or pigmented papule in children or young adults, that may resemble melanoma clinically or histologically.

\medterm{Splanchnic} Relating to internal organs, especially abdominal organs and their blood vessels or nerves. Splanchnic circulation supplies the digestive organs, while splanchnic nerves carry involuntary signals that regulate their function.
\medterm{Splanchnic Nerve} A nerve carrying automatic signals between the spinal cord and internal organs, especially organs in the chest and abdomen. It helps regulate digestion, blood-vessel tone, and other functions that do not require conscious control.

\medterm{Spleen} A lymphoid organ in the upper-left abdomen. Its red pulp filters blood and removes aged or abnormal blood cells; its white pulp supports immune responses. It also stores platelets and helps clear blood-borne microorganisms.
\medterm{Spleen Development} The formation and maturation of the spleen before and after birth. The spleen gradually develops its blood-filtering and immune functions, including removal of aged blood cells and responses to germs in the bloodstream.

\medterm{Spleen Disorder} A disease affecting the spleen, such as enlargement, infection, infarction, rupture, inflammation, cyst, or tumor, with possible effects on blood cells and immunity.

\medterm{Spleen Enlarged (Enlarged Spleen)} Alternate index label for Enlarged Spleen.

\medterm{Spleen Hemangioma} A spleen hemangioma is a usually benign blood-vessel tumor of the spleen, often incidental but occasionally associated with enlargement, anemia, or rupture risk.

\medterm{Spleen Hilum} The indented area on the inner surface of the spleen where its blood vessels, nerves, and lymph vessels enter or leave. The splenic vein carries blood from the spleen toward the liver through the portal circulation.

\medterm{Spleen Removal} Surgery to remove the spleen, which may be needed for injury, an enlarged or overactive spleen, blood disease, or a tumor. Without a spleen, the risk of severe infection increases, so vaccination and medical advice about fever are important.


\medterm{Spleen, Diseases Of} Disorders affecting splenic structure or function, including splenomegaly, infarction, infection, cysts, abscess, rupture, vascular disease, and hematologic or infiltrative disorders.
\medterm{Splenectomy} Surgical removal of the spleen, performed for selected traumatic, hematologic, malignant, vascular, or other disorders. It increases the risk of severe infection with encapsulated bacteria and requires preventive planning.
\medterm{Splenic Abscess} A localized collection of pus in the spleen, usually from bloodstream infection, trauma, infarction, or spread from nearby disease and potentially causing fever, left-upper-quadrant pain, and sepsis.

\medterm{Splenic Cyst} A fluid-filled lesion within the spleen, classified as true cysts with an epithelial
lining, including congenital, epidermoid, and dermoid cysts, or pseudocysts without a
lining, typically resulting from trauma, infection, or pancreatitis. Most are
asymptomatic and discovered incidentally on imaging.

\medterm{Splenic Flexure} The bend of the colon where the transverse colon becomes the descending colon, near the spleen.

\medterm{Splenic Flexure Syndrome} Pain, fullness, or bloating in the upper-left abdomen caused by gas collecting at the bend where the transverse colon meets the descending colon. Symptoms may improve after passing gas or stool.
\medterm{Splenic Infarction} Ischemic necrosis of splenic tissue caused by obstruction of splenic arterial or venous blood flow. Causes include embolism, thrombosis, hematologic disease, hypercoagulability, infection, trauma, or vasculitis; it often causes sudden left-upper-quadrant pain and is evaluated with imaging.

\medterm{Splenic Infection} Infection of the spleen, often an abscess caused by bacteria spreading through the bloodstream or from nearby tissue. It may cause fever, chills, left-upper-abdominal pain, and severe illness and usually needs urgent medical treatment.
\medterm{Splenic Neoplasm} An abnormal growth in the spleen. It may be noncancerous, cancerous, or cancer that has spread from another organ; symptoms and treatment depend on its type, size, and effect on nearby tissues or blood cells.

\medterm{Splenic Rupture} A tear of the spleen, usually after blunt abdominal trauma but sometimes spontaneous in an enlarged or diseased spleen. It can cause life-threatening internal bleeding; diagnosis and treatment depend on hemodynamic stability and imaging, with urgent surgery or embolization sometimes required.

\medterm{Splenomegaly} Abnormal enlargement of the spleen, usually secondary to infection, liver or portal-venous disease, hemolysis, hematologic malignancy, storage disease, or inflammation. It may cause left-upper-quadrant fullness or pain and increases rupture risk.
\medterm{Splenomegaly Gaucher (Gaucher Disease)} Alternate index label for Gaucher Disease.

\medterm{Splice-Site Variant} A variant affecting the sequences required for accurate removal of introns and joining
of exons during RNA processing. It may cause abnormal transcripts, exon skipping, intron
retention, or reduced functional protein.

\medterm{Spliceosome} The spliceosome is a nuclear RNA-protein complex that removes introns from pre-messenger RNA and joins exons to form mature RNA.

\medterm{Splint} A device that supports and limits movement of an injured or painful body part. Splints may be temporary or definitive and should not be so tight that they impair circulation or sensation.

\medterm{Splint (Dental)} Fixation of mobile teeth or fractured tooth. Wire and composite, orthodontic bracket splint. A dental device or wire used to stabilize loose teeth or a fractured tooth while healing occurs.
\synonyms
Dental splint; tooth splint

\medterm{Splinter Haemorrhages} Small, narrow lines of bleeding beneath a nail. They often follow minor injury but can also occur with heart-valve infection, blood-vessel inflammation, psoriasis, or other illnesses and need context for interpretation.
\medterm{Splinter Hemorrhages} Thin, lengthwise lines of bleeding beneath a nail. Minor injury is common, but they can also occur with heart-valve infection, blood-vessel inflammation, psoriasis, or other illnesses; their meaning depends on the overall clinical picture.

\medterm{Splinter Removal} Taking a piece of wood, metal, glass, or another material out of the skin. Small superficial splinters may be removed with clean tweezers, while deep, painful, infected, or eye- or nail-area splinters need clinical care.

\medterm{Splints} Devices that support, immobilize, align, or protect an injured body part, reducing pain and allowing healing.
\medterm{Split Brain} A state after surgical or pathological disruption of the corpus callosum in which information transfer between cerebral hemispheres is reduced, producing characteristic disconnection findings.

\medterm{Split-Thickness Skin Graft} A graft containing the epidermis and part of the dermis, used to cover wounds or burns and allowing the donor site to heal from remaining skin structures.

\medterm{Splitting} A psychological defense pattern in which a person or situation is viewed as entirely good or entirely bad, with difficulty integrating
positive and negative qualities. It can occur in several mental-health conditions and is
not diagnostic by itself.

\medterm{Spondee} A two-syllable word with equal emphasis on both parts, used during speech-hearing tests. Examples help measure the quietest loudness at which a person can recognize spoken words accurately.

\medterm{Spondylarthritis} A family of inflammatory rheumatic diseases affecting the spine, sacroiliac joints, entheses, peripheral joints, skin, eyes, or bowel, including ankylosing spondylitis and psoriatic arthritis.

\medterm{Spondylitis} Inflammation of one or more vertebrae or parts of the spine. Causes include infection, autoimmune disease, and inflammatory arthritis; pain, stiffness, fever, or neurologic symptoms may occur.
\medterm{Spondyloarthritis} A group of inflammatory joint diseases that often affect the spine and joints between the spine and pelvis. They may also cause tendon pain, swollen fingers or toes, eye inflammation, psoriasis, or bowel inflammation.

\medterm{Spondylolisthesis} Slipping of one spinal bone forward or backward compared with the bone below it. It may cause back pain, stiffness, leg pain, numbness, or weakness, although some people have no symptoms.
\medterm{Spondylolysis} A small stress crack or defect in part of a spinal bone, usually in the lower back. It may cause back pain in adolescents or athletes and can sometimes allow one spinal bone to slip on another.

\medterm{Spondylosis} Wear-related changes in the bones, discs, and small joints of the spine. It becomes more common with age and may cause stiffness or pain, but many people have no symptoms; nerve pressure can cause numbness or weakness.
\medterm{Spondylosis, Cervical} Alternate index label; see \textit{Cervical spondylosis}.

\medterm{Sponge} A sponge is a porous material used to absorb fluid, clean or dress wounds, or deliver a product; retained surgical sponges are a preventable complication.

\medterm{Sponge (Vascular)} A sterile absorbent material used during blood-vessel procedures to absorb blood, apply pressure, or protect delicate tissue. It must be counted and removed when required to prevent retained-material complications.
\synonyms
Vascular surgical sponge

\medterm{Spongiform Encephalopathy} A serious progressive brain disease in which brain tissue develops tiny sponge-like spaces. Prion diseases are important examples and can cause worsening memory, movement, coordination, behavior, and eventually loss of independence.
\medterm{Spongy Bone} Porous trabecular bone found inside many bones, especially near joints.

\synonyms
Cancellous bone; trabecular bone

\medterm{Spontaneous Abortion} Pregnancy loss before fetal viability that occurs without intentional intervention, commonly caused by chromosomal abnormalities but also by uterine, endocrine, immune, infectious, or other factors.

\medterm{Spontaneous Abortion (Miscarriage)} Unintentional loss of a pregnancy before fetal viability. Causes include chromosomal abnormalities, uterine or cervical problems, infection, and maternal medical conditions.

\synonyms
Miscarriage

\medterm{Spontaneous Bacterial Peritonitis} An infection of ascitic fluid without an identifiable abdominal source, usually occurring in people with cirrhosis and ascites. It may cause fever, abdominal pain, tenderness, or worsening illness. Diagnosis is made by testing ascitic fluid, usually showing at least 250 neutrophils per mm\textsuperscript{3}.

\medterm{Spontaneous Coronary Artery Dissection} A tear within the wall of an artery supplying the heart, which can narrow or block blood flow and cause a heart attack. It often affects women and may be associated with pregnancy, severe stress, or fibromuscular dysplasia. Treatment depends on blood flow and stability; medicines or a stent may be used, and emergency care is needed for chest pain.

\synonyms
SCAD


\medterm{Spontaneous Breathing Trial} A monitored period during which a person on a breathing machine does most of the breathing with little or no machine assistance. It helps clinicians judge whether the person may breathe safely after the tube is removed.

\medterm{Spontaneous Perforation} Unexpected rupture of a hollow organ without a clear external injury, allowing contents to escape and potentially causing peritonitis, mediastinitis, sepsis, or hemorrhage.

\medterm{Spontaneous Pneumothorax} Air leaking into the space around a lung without a direct injury. It may occur in an otherwise healthy person or because of lung disease and can cause sudden chest pain or breathlessness.

\medterm{Sporadic} Occurring irregularly or as an isolated case rather than in a predictable pattern or family cluster.

\medterm{Sporothrix} A genus of environmental fungi that can cause sporotrichosis, usually after inoculation through the skin and sometimes through inhalation or dissemination.

\medterm{Sporothrix Schenckii} A group of fungi that causes sporotrichosis, a typically nodular skin and lymphatic infection acquired from plants, soil, or animal scratches.

\medterm{Sporotrichosis} Sporotrichosis is a fungal infection caused by Sporothrix species, often introduced through plant matter and producing nodules along lymphatic channels.

\medterm{Sporozoa} An older group name for spore-forming protozoan parasites, including Plasmodium.
\medterm{Sporozoites} Moving, infective forms produced by certain parasites. For example, mosquitoes inject Plasmodium sporozoites into people, after which they travel to the liver and begin the malaria infection cycle.
\medterm{Sports Cream Overdose} Excessive use or swallowing of a muscle or joint pain cream, often containing menthol, methyl salicylate, or capsaicin. It may cause skin burns, salicylate poisoning, vomiting, ringing ears, confusion, or seizures.

\medterm{Sports Physical} A preparticipation health evaluation that reviews medical history, symptoms, medications, family risks, and examination findings to identify conditions needing evaluation before sports.

\medterm{Sports Rehabilitation} Treatment and training that help a person recover safe movement, strength, fitness, and confidence after a sports injury. It may include exercises, hands-on care, education, gradual return to activity, and steps to reduce the chance of another injury.

\medterm{Spots Before The Eyes} Visual floaters, flashes, or scotomata perceived in the visual field; sudden onset requires urgent assessment.
\medterm{Spots In Front Of The Eyes} Floaters or visual spots caused by shadows from vitreous opacities, inflammation, bleeding, retinal tears, migraine, or other eye or neurologic conditions; sudden flashes or a curtain-like loss are urgent.

\medterm{Spotted Fever} A group of infections usually spread by ticks and caused by Rickettsia bacteria. They can produce fever, headache, rash, muscle aches, and serious illness requiring prompt diagnosis and antibiotic treatment.
\medterm{Spousal Abuse} Spousal abuse is physical, sexual, emotional, financial, or coercive harm by an intimate partner; immediate safety planning and confidential support may be needed.

\medterm{Spouse} A spouse is a person’s legally or socially recognized partner; involving a spouse in healthcare requires the patient’s consent unless law provides otherwise.

\medterm{Sprain} An injury in which a ligament that supports a joint is stretched or torn. It can cause pain, swelling, bruising, and looseness; severity ranges from minor stretching to a complete tear and joint instability.

\medterm{Sprain Neck (Whiplash)} Alternate index label for Whiplash.

\medterm{Sprained Ankle} Stretching or tearing of one or more ankle ligaments, usually after rolling the ankle inward. It may cause pain, swelling, bruising, and instability.

\medterm{Sprains} Injuries in which ligaments around a joint are stretched or torn. They can cause pain, swelling, bruising, and joint looseness; severity ranges from minor stretching to complete tearing.
\medterm{Spray Drying} Spray drying converts a liquid mixture into a dry powder by atomizing it into a heated gas stream, often used in pharmaceutical and food manufacturing.

\medterm{Sprengel Deformity} A birth difference in which one shoulder blade sits higher than usual because it did not move to its normal position during development. It may cause uneven shoulders, limited arm movement, or neck and back discomfort.

\medterm{Spring Fever (Henoch-Schonlein Purpura)} Alternate index label for Henoch-Schonlein Purpura.

\medterm{Sprue} A condition causing poor absorption of nutrients from the small intestine. The term may describe celiac disease or tropical sprue, which have different causes and require appropriate testing and treatment.
\medterm{Sprue, Nontropical} Nontropical sprue is an older term for celiac disease, an immune reaction to gluten that damages the small-intestinal lining and impairs absorption.


\medterm{Spur Heel (Heel Spurs)} Alternate index label for Heel Spurs.

\medterm{Spur-Cell Hemolytic Anemia} A severe anemia in which abnormal, spiky red blood cells are rapidly removed from the blood, usually because of advanced liver disease. It can cause jaundice, weakness, and worsening anemia.

\medterm{Sputum} Mucus or other material brought up by coughing from the airways or lungs. Its amount, color, smell, and blood content may provide clues, but sputum alone cannot identify the cause; laboratory testing may be needed.
\medterm{Sputum Culture} A laboratory test that grows germs from mucus coughed up from the lungs or collected by suction. It can help identify bacterial, fungal, or mycobacterial infection and guide treatment, but a poor sample can give misleading results.

\medterm{Sputum Direct Fluorescent Antibody Test} A laboratory test that uses glowing antibody markers to look for selected germs in mucus from the lungs. It can provide a rapid result, but culture or molecular testing may also be needed.

\medterm{Sputum Fungal Smear} Examination of coughed-up lung mucus under a microscope to look for fungi. A culture, blood test, or molecular test may be needed to confirm the specific infection.

\medterm{Sputum Gram Stain} A quick microscope test that colors a sputum sample to show some bacteria and inflammatory cells and to judge whether the sample is suitable for culture.

\medterm{Sputum Stain For Mycobacteria} A microscope test that looks for tuberculosis-causing and related bacteria in mucus from the lungs. A positive result supports infection, but a negative result does not exclude it; culture or molecular testing may also be needed.



\medterm{Squamous Cell} A flat cell that forms part of the covering or lining of the skin and many body passages. These cells protect against friction and infection; changes in them can cause inflammation, precancer, or squamous cell cancer in organs such as the skin, mouth, cervix, lungs, or esophagus.

\medterm{Squamous Cell Carcinoma} A cancer of flat cells that line the skin or certain body passages. On the skin it may appear as a rough, scaly, crusted, wart-like, or nonhealing sore and can invade nearby tissue or spread if untreated.

\medterm{Squamous Cell Carcinoma Of The Lung} A type of non-small-cell lung cancer that usually begins in the lining of a larger airway. It is strongly linked to smoking but can occur in nonsmokers. Symptoms may include a persistent cough, coughing blood, chest pain, breathlessness, or repeated chest infections; treatment depends on stage and tumor features.


\medterm{Squamous Cell Carcinoma Of The Skin} A skin cancer arising from squamous cells, often appearing as a scaly, crusted, wart-like, or ulcerated growth.

\synonyms
Cancer, Squamous Cell

\medterm{Squamous Cell Skin Cancer} Squamous cell skin cancer is a malignant keratinocyte tumor often arising on sun-exposed skin and capable of local invasion or metastasis, especially when high risk.

\medterm{Squamous Cells} Flat epithelial cells with scale-like appearance that line surfaces such as skin, mouth, cervix, and esophagus; abnormal numbers or forms can reflect irritation, infection, dysplasia, or cancer.

\medterm{Squamous Intraepithelial Lesion Cervical (Cervical Dysplasia)} Alternate index label for Cervical Dysplasia.

\medterm{Square Foot} A square foot is a unit of area equal to a square measuring one foot on each side, used to describe surfaces, rooms, wounds, or treatment areas.

\medterm{Square Inch} A square inch is a unit of area equal to a square measuring one inch on each side, used for small surfaces or wound measurements.


\medterm{Squint} Misalignment of the eyes so they do not point in the same direction. It may be constant or intermittent and can cause double vision, poor depth perception, or reduced vision in one eye.

\medterm{SSRI} An SSRI is a selective serotonin-reuptake inhibitor antidepressant used for depression and several anxiety-related disorders; effects and withdrawal symptoms require monitoring.

\medterm{SSRIs} Selective serotonin reuptake inhibitors, a group of medicines that increase serotonin signaling in the brain. They are used for depression and several anxiety-related disorders; effects and side effects develop over time, not immediately.
\medterm{St John's Wort} A herbal product made from Hypericum perforatum and sometimes used for depressive symptoms. It can cause side effects and reduce the effectiveness of many medicines, including some contraceptives, anticoagulants, transplant medicines, and HIV treatments.
\medterm{ST Segment} The part of an electrocardiogram between the electrical activation and recovery signals of the heart's lower chambers. Abnormal elevation or depression can occur with a heart attack, inflammation, or other conditions and must be interpreted with symptoms and the full tracing.

\medterm{St Vitus's Dance} An older name for Sydenham chorea, involuntary movements following group A streptococcal infection.
\medterm{ST-Segment Elevation} Elevation of the ECG ST segment above the isoelectric baseline, reflecting abnormal
ventricular repolarization. It may indicate acute coronary occlusion but can also occur
in pericarditis, early repolarization, ventricular aneurysm, or other conditions;
interpretation requires clinical context and serial testing.


\medterm{St. John's Wort} An herbal product sometimes used for depressive symptoms. It can cause side effects and reduce the effectiveness of many medicines, including some contraceptives, anticoagulants, transplant medicines, and HIV treatments.
\medterm{Stab Wound} A deep puncture caused by a pointed or sharp object. The wound may damage blood vessels, nerves, organs, or bone even when the skin opening looks small; heavy bleeding, severe pain, or a wound to the chest or abdomen requires emergency care.

\medterm{Stable Angina} Repeated, predictable chest pressure or discomfort brought on by activity or stress because the heart muscle receives too little blood. It usually improves with rest or prescribed medicine; new, severe, prolonged, or rest pain may be a heart attack and needs emergency care.

\synonyms
Chronic stable angina

\medterm{Stable Coronary Artery Disease} Coronary artery narrowing causing a consistent, predictable pattern of angina.

\synonyms
Stable ischemic heart disease

\medterm{Stable Disease} An oncology response category in which a cancer has neither shrunk enough to meet criteria for partial response nor grown enough to meet criteria for progressive disease. It describes change over a specified assessment interval and does not mean that the cancer is cured or inactive.

\medterm{Stabs} Brief, sharp pains or penetrating injuries described as stabbing. Causes vary with location and may include nerve irritation, muscle strain, infection, inflammation, injury, or serious disease requiring urgent assessment.
\medterm{Stage 4 Prostate Cancer} Alternate index label; see \textit{Metastatic (stage 4) prostate cancer}.

\medterm{Stage I} A cancer stage that generally indicates limited disease, although the exact definition and prognosis depend on the cancer type and staging system.

\medterm{Stage II} A cancer stage that usually indicates a larger or more locally advanced tumor than stage I, without distant spread. The exact tumor size, lymph-node findings, prognosis, and treatment vary by cancer type and staging system.

\medterm{Stage III} A cancer stage that often indicates more extensive local or regional disease, with the precise criteria depending on the cancer type.

\medterm{Stage IV} A cancer stage that generally indicates spread to distant organs or sites. Prognosis and treatment vary by cancer type and individual factors.

\medterm{Stages Of Change} A model describing common steps people may pass through when changing a behavior. These include not yet considering change, thinking about it, preparing, taking action, and maintaining the change; people may move backward or forward between stages.

\synonyms
Transtheoretical model

\medterm{Stages Of Labor} The first stage is cervical opening, the second is birth of the baby, and the third is delivery of the placenta. The first stage has early and active phases; length and symptoms vary, and the uterus is monitored after delivery.

\medterm{Staghorn Calculi} Large branching kidney stones that fill part or all of the kidney's urine-collecting system. They are often linked with infection, can damage kidney tissue, and usually require specialist treatment to remove them.

\medterm{Staghorn Calculus} A branching kidney stone that fills part or all of the renal pelvis and calyces, often associated with infection.
\medterm{Staging} Determining how far a cancer has grown or spread by assessing the original tumor, nearby lymph nodes, and distant organs. The stage helps estimate outlook and choose treatment.

\medterm{Staining} A laboratory method that uses colored dyes to make cells, tissue structures, or microorganisms easier to see under a microscope. Different stains highlight different features and can help guide diagnosis, although results may need culture or other tests for confirmation.

\medterm{Stainless Steel} A corrosion-resistant alloy containing iron and chromium, used for instruments, implants, orthodontic appliances, and some dental crowns. Its suitability depends on the grade, strength, and tissue exposure.

\medterm{Stammering} Alternate name for stuttering, a speech-fluency difficulty involving repeated sounds or words, prolonged sounds, or blocks that interrupt speaking. It is common in childhood and may continue into adulthood; speech-language therapy can help.
\medterm{Standard Deviation} A number that shows how widely measurements differ from their average. A small standard deviation means values are close together, while a large one means they are more spread out; it does not by itself show whether results are good or bad.
\medterm{Standard Eye Exam} A comprehensive eye examination assessing visual acuity, refraction, pupils, eye movements, external structures, anterior eye, intraocular pressure, and ocular fundus as indicated.

\medterm{Standard Of Care} The conduct, skill, diligence, and judgment a reasonably prudent practitioner of the
same specialty would exercise under similar circumstances; the benchmark for medical
negligence. It is established through guidelines, peer practice, and expert testimony
rather than fixed statute.

\medterm{Standard Precautions} Basic infection-control measures used for every patient, because blood, body fluids, nonintact skin, and mucous membranes may carry germs. They include hand hygiene, appropriate protective equipment, safe injections, sharps safety, and respiratory etiquette.
\medterm{Standardized Extract} An herbal extract adjusted to contain a measured amount of selected ingredient, improving consistency but not guaranteeing safety or effectiveness.

\medterm{Standardized Uptake Value} A number used on a PET scan to estimate how much radioactive tracer a tissue has taken up, adjusted for the injected amount and body size. It can help compare areas or follow change, but it does not by itself prove cancer or another disease.

\medterm{Stannosis} A usually mild form of pneumoconiosis caused by long-term inhalation of tin oxide dust. It may produce visible lung changes on imaging, often without major breathing impairment.
\medterm{Stanozolol} A synthetic anabolic androgenic steroid with limited medical use and potential liver, lipid, and cardiovascular toxicity.
\medterm{Stapedectomy} Surgical removal or replacement of the stapes, usually to treat otosclerosis and conductive hearing loss.
\medterm{Stapedius} A small middle-ear muscle that helps limit movement of the stapes and may protect the inner ear from loud sounds.

\medterm{Stapes} The smallest middle-ear bone, passing sound vibrations from the incus to the inner ear for hearing.
\medterm{Staph Infection} Alternate informal name; see \textit{Staphylococcal Infection}.



\medterm{Staph Infections In The Hospital} Infections caused by Staphylococcus bacteria that begin during or after healthcare treatment, or require hospital care. They may involve wounds, blood, lungs, bones, or devices and can be difficult to treat when the bacteria resist antibiotics.

\medterm{Staphylococcal Food Poisoning} An acute foodborne intoxication caused by the ingestion of preformed heat-stable enterotoxins produced by Staphylococcus aureus in improperly refrigerated or handled foods (such as dairy products, pastries, processed meats, and salads). Characterized by an abrupt onset of severe nausea, spasmodic abdominal cramping, and violent vomiting within 1 to 6 hours of consumption, usually resolving spontaneously within 24 to 48 hours.

\synonyms
Staphylococcal Enterotoxicosis; Staph Food Poisoning

\medterm{Staphylococcal Infection} An infection caused by Staphylococcus aureus or another Staphylococcus species. It may affect
the skin, bone, joints, lungs, heart valves, blood, or other tissues. Some strains resist
multiple antibiotics. Diagnosis and antimicrobial treatment depend on the site, severity,
and susceptibility of the organism.

\synonyms
Staph Infection; Staphylococcus Infection; Staphylococcal Infections

\medterm{Staphylococcal Meningitis} Staphylococcal meningitis is infection and inflammation of the meninges caused by Staphylococcus bacteria, often after surgery, trauma, bloodstream infection, or an invasive device.

\medterm{Staphylococcal Pneumonia} Pneumonia caused by Staphylococcus bacteria, sometimes after influenza or another viral infection. It may cause fever, cough, chest pain, breathlessness, lung abscesses, or infected fluid around the lung and needs medical treatment.
\medterm{Staphylococcal Scalded Skin} Alternate index label; see \textit{Staphylococcal Scalded Skin Syndrome}.

\medterm{Staphylococcal Scalded Skin Syndrome} A toxin-mediated skin disorder caused by Staphylococcus aureus, usually affecting infants or young children. It produces widespread redness, tenderness, fragile superficial blisters, and peeling that can resemble a burn and requires prompt treatment.

\medterm{Staphylococcus} A genus of gram-positive cocci that includes S. aureus and common skin colonizers and pathogens.
\medterm{Staphylococcus Aureus} A bacterium that commonly lives on the skin or in the nose but can cause infections of the skin, blood, lungs, bones, joints, or heart valves. Some strains are resistant to several antibiotics.

\medterm{Staphylococcus Infection (Staph Infection)} Alternate index label; see \textit{Staphylococcal Infection}.

\medterm{Staple (Vascular)} A small metal or polymer device used to close tissue or join selected blood-vessel structures. It provides rapid closure but can cause narrowing, bleeding, thrombosis, or tissue injury if unsuitable or misplaced.
\synonyms
Vascular staple; vessel-closure staple

\medterm{Stapling} Closing or joining tissue with metal or absorbable surgical staples. Staples may be used for skin, blood vessels, or parts of the digestive tract and must be selected and removed or absorbed appropriately.

\medterm{Starch Poisoning} Illness after swallowing a starch product, most often causing choking, blockage, or stomach upset rather than systemic poisoning. Inhaled powder can irritate the airways; breathing difficulty or persistent vomiting requires medical care.

\medterm{Stargardt Disease} An inherited eye disorder that gradually damages the macula, the central part of the retina. It usually begins in childhood or young adulthood and causes loss of detailed central vision; side vision is often preserved.

\medterm{Starling Forces} The balance of pressures that pushes water out of small blood vessels or pulls it back in. Changes can cause swelling in tissues, dehydration of tissues, or fluid buildup around organs.

\medterm{Starvation} Severe lack of food or usable energy over time, causing weight and muscle loss, weakness, changes in body chemistry, reduced immunity, and eventually organ failure. Restarting nutrition after prolonged starvation must be done carefully because dangerous shifts can occur.
\medterm{Stasis} Slowing or stopping of normal movement of blood, lymph, intestinal contents, or other body fluids.
\medterm{Stasis Dermatitis} Inflammation of the lower-leg skin caused by chronic venous hypertension and edema. It may cause itching, redness,
scaling, brown discoloration, weeping, or ulcers around the ankles. Management focuses on treating venous disease,
compression when safe, skin care, and evaluation for infection or arterial insufficiency.

\medterm{Stasis Dermatitis And Ulcers} Stasis dermatitis and ulcers result from chronic venous hypertension in the lower legs, causing swelling, brown skin change, inflammation, itching, and slow-healing wounds.

\medterm{Stasis Ulcer} A stasis ulcer is a chronic lower-leg wound caused by venous hypertension and poor return of blood, commonly near the ankle and associated with edema and skin changes.

\medterm{STAT} A medical instruction meaning immediately or with urgent priority, such as a stat test, medicine, or consultation.
\medterm{Statin} A cholesterol-lowering medicine that reduces the liver’s production of cholesterol and lowers LDL, often called bad cholesterol. Statins reduce the risk of heart attack and stroke in appropriate people; muscle symptoms and medicine interactions should be reviewed.

\medterm{Statin-Induced Myopathy} Muscle pain, weakness, or tenderness caused by a statin cholesterol-lowering medicine. Symptoms may occur without muscle damage or, rarely, with severe muscle breakdown that can injure the kidneys. A clinician may check the medicine, symptoms, muscle enzymes, and other causes before changing treatment.

\synonyms
Statin-Associated Muscle Symptoms (SAMS); Statin Myotoxicity

\medterm{Statins} A class of medicines that lowers LDL cholesterol by reducing cholesterol production in the liver. They reduce the risk of heart attack and stroke in appropriate people; muscle symptoms, liver problems, and medicine interactions should be reviewed.

\synonyms
HMG-CoA reductase inhibitors

\medterm{Statins In Diabetes} Cholesterol-lowering treatment for people with diabetes who may benefit from reduced risk of heart attack and stroke. The medicine and dose depend on age, cardiovascular risk, kidney function, other medicines, and possible side effects.

\medterm{Station} A measure of how far the baby’s presenting part, usually the head, has descended through the pelvis during labor. It is compared with the mother’s pelvic bones and helps clinicians follow progress toward birth.

\medterm{Status Asthmaticus} A severe asthma attack that does not improve enough with the first treatments. Breathing may become dangerously difficult, with exhaustion, inability to speak normally, or low oxygen, and emergency hospital treatment is required.
\medterm{Status Epilepticus} A prolonged seizure or repeated seizures without recovery, requiring emergency treatment to prevent brain injury and complications.
\medterm{Status Migrainosus} A severe migraine attack that persists for more than 72 hours and may require urgent medical treatment, especially when dehydration, vomiting, or medication overuse is present.


\medterm{Stay Away From Asthma Triggers} Avoiding asthma triggers such as smoke, allergens, infections, fumes, and cold air can reduce symptoms and attacks, alongside prescribed controller and reliever treatment.

\medterm{STDs} Alternate index label; see \textit{Sexually transmitted diseases (STDs)}.

\medterm{Steady State} The point during regular medicine use when the amount entering the body over time is about equal to the amount leaving it. Blood levels then become relatively stable; reaching this point usually takes several doses and depends on the medicine’s half-life.

\medterm{Steal Syndrome} Reduced blood flow to a hand, foot, or other area because a bypass graft or dialysis fistula diverts too much blood. It may cause pain, coolness, pale or blue skin, numbness, or weakness and requires vascular assessment.

\medterm{Steam Iron Cleaner Poisoning} Poisoning from a steam-iron cleaning product, which may contain acids or other corrosive chemicals. It can burn the mouth, digestive tract, skin, or eyes and may irritate the lungs; do not induce vomiting.

\medterm{Steatoma} A fatty or sebaceous mass; the term may refer to a lipoma or, in the ear, a cholesteatoma depending on context.
\medterm{Steatorrhea} The passage of excess fat in the stool, causing greasy, pale, foul-smelling, floating
stools—a hallmark of fat malabsorption.

\medterm{Steatorrhoea} Excess fat in stool, causing pale, bulky, greasy, foul-smelling bowel movements. It suggests poor fat digestion or absorption from pancreatic, liver, bile-duct, intestinal, or other digestive disease.
\medterm{Steele-Richardson-Olszewski Syndrome} Alternate index label; see \textit{Progressive supranuclear palsy}.

\medterm{Stein-Leventhal Syndrome} Alternate index label; see \textit{Polyendocrine metabolic ovarian syndrome}.

\medterm{Stem Cell} A cell that can make more stem cells and develop into one or more specialized cell types. Some stem cells produce blood cells; others help form or repair different tissues.

\medterm{Stem Cell Donation} The collection of healthy blood-forming stem cells from a donor for transplantation into another person. Cells may be collected from the bloodstream by apheresis after medication increases their number in the blood, or directly from bone marrow under anesthesia. Donor and recipient compatibility, especially human leukocyte antigen matching, helps improve transplant success.

\medterm{Stem Cell Transplantation} Hematopoietic stem cell transplantation replaces diseased or ablated bone marrow
with healthy donor or autologous stem cells.

\medterm{STEMI} A type of heart attack caused by sudden blockage of a heart artery, usually producing characteristic changes on an electrocardiogram. It is an emergency because quickly restoring blood flow can limit permanent heart damage.

\synonyms
ST-Elevation Myocardial Infarction

\medterm{Stenosing Tenosynovitis} Alternate index label; see \textit{Trigger finger}.

\medterm{Stenosis} Abnormal narrowing of a passage, blood vessel, heart valve, spinal canal, or other opening. Narrowing may restrict flow or compress nearby structures.
\medterm{Stenosis Lumbar (Lumbar Stenosis)} Alternate index label for Lumbar Stenosis.

\medterm{Stenosis Of The Left Anterior Descending Artery} Narrowing of a major artery supplying the front of the heart. Reduced blood flow may cause chest pain during exertion or a heart attack if the narrowing becomes severe or the artery suddenly blocks.

\medterm{Stenosis Spinal (Lumbar Stenosis)} Alternate index label for Lumbar Stenosis.

\medterm{Stenosis, Aortic Valve} Alternate index label; see \textit{Aortic valve stenosis}.

\medterm{Stenosis, Mitral Valve} Alternate index label; see \textit{Mitral valve stenosis}.

\medterm{Stenosis, Pulmonary Valve} Alternate index label; see \textit{Pulmonary valve stenosis}.

\medterm{Stenosis, Pyloric} Alternate index label; see \textit{Pyloric stenosis}.

\medterm{Stenosis, Spinal} Alternate index label; see \textit{Spinal stenosis}.

\medterm{Stenotic} Stenotic means narrowed, as in a stenotic heart valve, blood vessel, airway, intestine, or other passage.

\medterm{Stenotrophomonas Maltophilia Resistance} The ability of Stenotrophomonas maltophilia bacteria to survive several antibiotics. Treatment depends on the infected body site, illness severity, and laboratory testing that identifies which medicines are likely to work.

\medterm{Stent} A tube or scaffold placed inside a vessel or other passage to keep it open.
\medterm{Stent (Vascular)} A small mesh or tube placed inside a narrowed or blocked vessel to help keep it open.

\synonyms
Vascular stent

\medterm{Stent Balloon (Vascular)} A balloon on a catheter used to expand a vascular stent against the vessel wall and restore an open channel. Inflation can cause vessel injury, bleeding, clotting, or rarely rupture.
\synonyms
Stent-expansion balloon; vascular angioplasty balloon

\medterm{Stent Graft} A fabric-covered stent placed inside a blood vessel to reinforce the vessel wall or exclude an aneurysm from blood flow.

\medterm{Stent Placement} Insertion of a mesh tube or other support to keep a narrowed or blocked passage open. Stents may be used in arteries, bile ducts, ureters, airways, or the esophagus; bleeding, blockage, infection, movement, or repeated narrowing can occur.

\medterm{Stent Restenosis} Re-narrowing of a blood vessel within or adjacent to a previously placed stent, usually caused by tissue growth, recurrent atherosclerosis, or thrombosis.

\medterm{Stent-Graft (Vascular)} A fabric-covered metal frame placed inside a blood vessel to reinforce a weak wall or seal an aneurysm, tear, or leak. Follow-up imaging checks for movement, blockage, or leakage around the graft.
\synonyms
Vascular stent-graft; endovascular graft

\medterm{Stenting} Placement of a small mesh tube to keep a narrowed or blocked passage open. Vascular stents restore blood flow, while drug-coated types release medicine to reduce repeated narrowing.

\medterm{Stepping Reflex} A primitive reflex present from birth to 2 months, elicited by holding the infant
upright with feet touching a flat surface. The infant makes stepping-like movements as
if walking. Disappears by 2 months as voluntary motor control develops. Reappears around
9-10 months as voluntary walking emerges.

\medterm{Steps To Take Before You Get Pregnant} Preconception care includes reviewing health conditions and medicines, folic acid, vaccinations, dental and genetic considerations, substance avoidance, and planning prenatal care.

\medterm{Stereognosis} The ability to recognize a familiar object by feeling it without looking. It requires intact skin sensation, nerve pathways, and areas of the brain that interpret touch; loss may occur after nerve or brain injury.
\medterm{Stereotactic} Stereotactic describes a technique using three-dimensional coordinates to target a precise location in the brain or another body structure.

\medterm{Stereotactic Biopsy} Image-guided biopsy using three-dimensional coordinates to target a lesion that may be difficult to localize by direct palpation. It is used in the breast and, with specialized systems, the brain, lung, and other organs.

\medterm{Stereotactic Body Radiation Therapy} A focused form of external radiation that delivers a high dose to a small body tumor in a few treatments. It may treat selected lung, liver, kidney, pancreatic, or spinal tumors while limiting exposure to nearby tissue.

\medterm{SRT} Stereotactic radiotherapy, a precise form of external radiation treatment that uses three-dimensional imaging and careful positioning to target a small tumor or other abnormal area. It is often given in two to five sessions to limit exposure to nearby tissue.

\medterm{Stereotactic Brain Biopsy} A precision, minimally invasive neurosurgical procedure used to safely obtain small tissue samples from deep, delicate, or high-risk areas of the brain for microscopic examination. Using high-resolution pre-operative MRI or CT scans and a specialized three-dimensional stereotactic head frame or frameless optical infrared navigation system, a neurosurgeon maps exact mathematical coordinates to the suspected tumor or infection. Through a tiny circular burr hole (approximately 3 to 5 mm) in the skull, the surgeon gently passes a specialized hollow biopsy needle straight to the target lesion without disrupting surrounding healthy brain tissue. The extracted tissue specimens are examined by a neuropathologist to determine the exact tumor type, grade, genetic markers, or underlying infection, allowing doctors to design targeted chemotherapy, radiation, or medical therapy.
\synonyms
Stereotactic Needle Biopsy; Brain Biopsy; Frameless Stereotactic Biopsy

\medterm{Stereotactic Radiosurgery} Precise radiation treatment delivering one or a few high doses to a small target, often in the brain, without open surgery.

\medterm{Stereotactic Radiosurgery - CyberKnife} CyberKnife is a robotic system that delivers many precisely aimed radiation beams to a tumor or other target without an open operation.


\medterm{Stereotactic Radiosurgery - Gamma Knife} Gamma Knife is a focused radiation system that delivers multiple precisely aimed beams, mainly to selected brain tumors and vascular or other brain lesions.

\medterm{Stereotaxis} A technique using a three-dimensional coordinate system to target deep structures, especially in brain procedures.
\medterm{Stereotypic Movement Disorder} Repeated, apparently purposeless movements such as hand flapping, rocking, or head banging that interfere with activities or cause injury. It is more common in childhood and developmental conditions and requires assessment for other causes.

\medterm{Sterile Field} A designated area maintained free of viable microorganisms during an invasive procedure.
Only sterile equipment and materials may contact the field or the sterile portions of
instruments.

\medterm{Sterile Site} A body area or sample that normally contains no living microorganisms, such as blood, spinal fluid, joint fluid, or tissue from a closed space. Finding a germ there is usually important, although accidental contamination must be considered.

\medterm{Sterile Technique} Methods used to prevent introduction of microorganisms into a sterile field, wound, tissue, device, or body site through hand hygiene, barriers, aseptic handling, and appropriate preparation.

\medterm{Sterilisation} Destruction or removal of all forms of microbial life, or permanent contraception depending on context.
\medterm{Sterility} Inability to reproduce, or complete absence of living microorganisms from an object or preparation, depending on context. In healthcare, sterile equipment, implants, and medicines must remain free from contamination to prevent infection.
\medterm{Sterilization} A process that destroys or removes all forms of microbial life, including spores, from equipment, medicines, or other materials. In another context, sterilization means a permanent method intended to prevent pregnancy.

\medterm{Sternal Exploration Or Closure} Surgical reopening or repair of the breastbone after median sternotomy, performed to treat infection, bleeding, instability, dehiscence, or other postoperative chest complications.

\medterm{Sternoclavicular Joint} A saddle-type synovial joint between the clavicle and manubrium of the sternum. It is
the only bony connection between the upper limb and axial skeleton and permits movement
of the shoulder girdle.

\medterm{Sternocleidomastoid Muscle} A prominent muscle on each side of the neck that turns and bends the head and helps stabilize the neck. It can also assist forceful breathing; strain or spasm may cause neck pain and limited movement.

\medterm{Sternotomy} Surgical division of the sternum, usually along the midline, to provide access to the heart or mediastinal structures.

\medterm{Sternum} The breastbone forming the front of the central chest wall.
\medterm{Sternutation} Sneezing, a reflex forceful expulsion of air through the nose and mouth triggered by irritation of nasal or upper-airway sensory receptors.

\medterm{Steroid Myopathy} Muscle weakness caused by prolonged exposure to corticosteroid medicines or excessive cortisol. It usually develops gradually and affects the hips, thighs, shoulders, and upper arms, making it hard to rise from a chair or climb stairs. Weakness may improve when the cause is treated, but steroid doses must be changed only with medical guidance.

\synonyms
Corticosteroid-Induced Myopathy; Glucocorticoid Myopathy

\medterm{Steroids} A broad family of compounds with a characteristic four-ring structure, including body hormones, corticosteroids, and sex hormones. Their effects depend on the type; prolonged corticosteroid use can suppress immunity and bone health.
\synonyms
Corticosteroids

\medterm{Stertor} A low-pitched snoring, gurgling, or rattling sound caused by partly blocked air passages in the nose or throat. It may occur during sleep or illness; noisy breathing with severe sleepiness, blue lips, or breathing difficulty is an emergency.
\medterm{Stethoscope} An instrument used for auscultation, the examination of sounds produced within the body. It usually has a chest piece with a diaphragm and sometimes a bell, flexible tubing, and earpieces; electronic models convert sounds into amplified or recorded signals. It is used to examine heart, lung, bowel, and vascular sounds.
\medterm{Stevens-Johnson Syndrome} A rare, life-threatening reaction usually caused by a medicine or infection. It causes painful widespread skin and mucous-membrane damage, blistering, and peeling and requires emergency hospital treatment.

\synonyms
SJS

\medterm{STI} Alternate index label for Sexually Transmitted Infections.

\medterm{Stickler Syndrome} An inherited connective-tissue disorder that may cause severe nearsightedness, retinal detachment, hearing loss, early arthritis, and differences in the jaw or palate. Regular eye and hearing checks are important.

\medterm{Stiff Baby Syndrome} A nonspecific description of unusual stiffness or increased muscle tone in an infant, not a single diagnosis. Causes include brain, nerve, muscle, genetic, or metabolic disorders; persistent stiffness, poor feeding, seizures, or breathing problems need urgent assessment.

\medterm{Stiff-Person Syndrome} A rare immune-related nervous-system disorder causing persistent muscle stiffness and sudden painful spasms. Noise, touch, movement, or emotional stress may trigger spasms, leading to falls and difficulty walking. Blood tests may find related antibodies. Treatment can reduce stiffness and immune activity, but symptoms often require long-term care.

\synonyms
SPS; Stiff-Man Syndrome; Moersch-Woltman Syndrome

\medterm{Stiffness} Difficulty or discomfort when moving a body part. It may result from joint disease, muscle problems, nerve disorders, inflammation, injury, prolonged rest, or aging and can limit everyday activities.
\medterm{Stigma} A mark, characteristic, or social discrediting label; in medicine it may refer to a sign or to harmful social attitudes.
\medterm{Stilboestrol} An older name for diethylstilbestrol, a synthetic estrogen no longer used routinely because of serious risks.
\medterm{Stilet} A thin, pointed instrument placed inside a cannula or tube to guide insertion, clear a blockage, or provide firmness. It must be used carefully because it can injure tissue or blood vessels.
\medterm{Still Disease} An inflammatory disease that can affect the whole body.
In children it is called systemic juvenile idiopathic arthritis; in adults it is called
adult-onset Still disease. It can cause fever, joint pain or swelling, and a rash.

\synonyms
Still's Disease; Arthritis Still (Still's Disease)

\medterm{Stillbirth} The death of a fetus before or during delivery after the gestational-age or birth-weight threshold used by a country or health system. In the United States, it usually refers to loss at or after 20 weeks of pregnancy; the World Health Organization commonly uses 28 weeks for international comparisons. Causes may include problems with the placenta, fetal growth, infection, genetic or developmental conditions, or maternal illness, but some stillbirths remain unexplained.
\medterm{Still's Murmur} A common harmless heart murmur in children, usually heard between ages 2 and 7. It has a soft, musical sound, becomes louder with fever or exercise, and usually disappears as the child grows; the heart is otherwise normal.

\synonyms
Innocent Vibratory Murmur; Normal Pediatric Murmur; Twanging-String Murmur

\medterm{Stimulant} A substance that increases activity in the brain or other body systems. Stimulants may improve alertness and attention, but excessive or inappropriate use can cause dependence, anxiety, fast heartbeat, high blood pressure, or seizures.

\medterm{Stimulant Use Disorder} A pattern of using stimulant drugs that causes significant problems or loss of control despite harm. Stimulants include cocaine and amphetamine-type drugs; stopping may cause tiredness, low mood, increased appetite, or excessive sleep.

\medterm{Stimulants} Substances that increase activity in the brain or the body’s fight-or-flight system. They may improve alertness and attention but can also cause insomnia, anxiety, dependence, fast heartbeat, high blood pressure, or seizures.

\medterm{Stimulus} A stimulus is a physical, chemical, sensory, or social change that evokes a response from a cell, tissue, organism, or person.

\medterm{Stingray} A sea-animal injury caused by a stingray spine, usually producing severe pain, swelling, and a puncture wound. Medical care may be needed to clean the wound, remove retained spines, and treat whole-body reactions.

\medterm{Stings} Injuries caused when an insect, animal, or plant injects venom or irritating material. They may cause pain, swelling, itching, allergic reactions, tissue damage, or serious poisoning depending on the source.
\medterm{Stings, Bee} Alternate index label; see \textit{Bee sting}.

\medterm{Stings, Jellyfish} Alternate index label; see \textit{Jellyfish stings}.

\medterm{Stings, Scorpion} Alternate index label; see \textit{Scorpion sting}.

\medterm{Stippling (Tattooing)} Patterned dermal abrasions and punctate hemorrhages surrounding a firearm entrance
wound, caused by unburned and partially burned gunpowder grains striking the skin at
close range (15--75 cm).

\synonyms
Tattooing

\medterm{Stitch} A suture used to close tissue, or a sharp transient pain, often in the side during exercise.
\medterm{Stochastic Effects} Possible effects of radiation, especially cancer or inherited genetic changes, whose chance increases with dose. The severity cannot be predicted from the dose, and even low exposure may carry some risk.
\medterm{Stoddard Solvent Poisoning} Poisoning from Stoddard solvent, a petroleum-based cleaning solvent. Inhalation or swallowing can cause headache, dizziness, drowsiness, skin irritation, or aspiration-related lung injury; do not induce vomiting and seek advice.

\medterm{Stokes-Adams Syndrome} Sudden fainting caused by a temporary dangerous slowing or stopping of the heartbeat, often from heart block or another abnormal rhythm. It requires urgent heart assessment because injury or cardiac arrest can occur.
\medterm{Stoma} A natural opening or a surgically created opening connecting an organ to the skin or another organ.
\medterm{Stoma Complications} Problems involving an opening made from the intestine or urinary tract to the skin. They include skin irritation, bleeding, retraction, prolapse, narrowing, blockage, high output, infection, and a parastomal hernia; sudden severe pain, color change, or no output needs urgent care.

\medterm{Stomach} A muscular digestive organ between the oesophagus and small intestine that stores and begins digestion of food.
\medterm{Stomach Cancer} Cancer that begins in the stomach, most often in gland-forming cells lining the stomach. Risk factors include Helicobacter pylori infection, smoking, chronic inflammation, certain foods, and inherited conditions; symptoms and treatment depend on type and stage.

\synonyms
Cancer, Gastric
Cancer, Stomach

\medterm{Stomach Cramps} Intermittent gripping or spasmodic upper-abdominal pain caused by gas, infection, indigestion, ulcer disease, gallbladder or pancreatic disease, obstruction, or other causes.

\medterm{Stomach Fat (Belly Fat, Abdominal Fat)} Adipose tissue stored beneath the skin or around abdominal organs; excess visceral fat is metabolically active and associated with insulin resistance, fatty liver, and cardiovascular risk.

\medterm{Stomach Flu} A common informal term for viral gastroenteritis, an infection of the stomach and intestines caused by viruses rather than influenza viruses. It may cause watery diarrhea, vomiting, abdominal cramps, fever, and dehydration; some infections cause few or no symptoms.

\medterm{Stomach Paralysis} Stomach paralysis, or gastroparesis, is delayed emptying of stomach contents without obstruction, causing early fullness, nausea, vomiting, and bloating.

\medterm{Stomach Tube} A tube passed through the mouth or nose into the stomach for drainage, feeding, sampling, or treatment.
\medterm{Stomach Ulcer} An open sore in the lining of the stomach, commonly associated with Helicobacter pylori infection or medicines such as NSAIDs.

\medterm{Stomach Ulcer (Peptic Ulcer)} Alternate index label for Peptic Ulcer.

\medterm{Stomach Washout} A procedure that rinses and removes stomach contents through a tube. It is rarely used for selected serious poisonings because it can cause choking, lung injury, bleeding, or other complications.
\medterm{Stomach, Diseases Of} Disorders affecting the stomach, including inflammation, ulcers, movement problems, blockage, bleeding, and tumors. Symptoms may include pain, nausea, vomiting, early fullness, weight loss, indigestion, or blood in vomit or stool.
\medterm{Stomatitis} Inflammation, soreness, redness, or ulcers of the lining of the mouth. Causes include infection, irritation, nutritional deficiency, medicines, and cancer treatment; severe pain, dehydration, fever, or inability to eat or drink needs medical care.
\medterm{Stomatodynia} Alternate index label; see \textit{Burning mouth syndrome}.

\medterm{Stone} A hard lump formed inside an organ or tube, commonly from mineral salts or other body substances. Stones may block flow, cause pain or infection, and occur in kidneys, gallbladder, or ducts.
\medterm{Stonefish Sting} A puncture injury from the venomous spines of a stonefish, causing severe pain, swelling, bleeding, and sometimes weakness, breathing problems, or shock. Seek emergency care; hot-water immersion may reduce pain while help is obtained.

\medterm{Stool} Fecal material expelled from the bowel, composed of water, undigested residue, bacteria, cells, mucus, and bile pigments; color, form, frequency, and blood can provide clinical clues.

\medterm{Stool C Difficile Toxin} A stool test for harmful substances made by Clostridioides difficile bacteria. A positive result supports active infection when a person has compatible diarrhea, but bacteria can be present without causing illness.

\medterm{Stool Examination} Laboratory analysis of a stool sample for blood, infection, parasites, inflammation, malabsorption, or other digestive abnormalities.

\synonyms
Stool test

\medterm{Stool Gram Stain} Microscopic examination of a stool specimen after Gram staining; it may provide limited information about bacterial flora or inflammatory cells but is not a general test for infectious diarrhea.

\medterm{Stool Guaiac Test} A chemical fecal occult blood test that detects the peroxidase activity of heme in stool. Dietary factors, medicines, and specimen collection can affect results. It is a screening or detection test rather than a diagnosis; a positive result requires clinical follow-up and usually evaluation for a gastrointestinal source.

\medterm{Stool Holding} Alternate index label; see \textit{Encopresis}.

\medterm{Stool Occult Blood Test} Alternate index label; see \textit{Fecal occult blood test}.

\medterm{Stool Ova And Parasites Exam} Microscopic examination of stool for parasite eggs, cysts, larvae, or other forms, sometimes requiring multiple specimens to improve detection.

\medterm{Stools} Feces, the material expelled from the digestive tract. Changes in color, consistency, frequency, or blood can provide clues about digestive health but are not diagnostic alone.
\medterm{Stools - Foul Smelling} Strongly odorous feces that may result from diet, infection, bleeding, malabsorption, altered gut bacteria, or rapid transit; persistent change with weight loss or diarrhea warrants review.

\medterm{Stools - Pale Or Clay-Colored} Light-colored stool caused by reduced bile pigment reaching the intestine, potentially indicating bile-duct obstruction, gallbladder disease, liver disease, or certain transient conditions.

\medterm{Stools – Floating} Feces that remain on the water surface, often from trapped gas or diet but sometimes from excess fat in malabsorption; persistent greasy, pale, foul-smelling stool needs evaluation.

\medterm{Stop Smoking Support Programs} Counseling, quit plans, telephone or digital coaching, and medicines such as nicotine replacement, varenicline, or bupropion used to help people stop tobacco use.

\medterm{STOPP/START} A checklist used to review medicines in older adults. STOPP identifies medicines that may be inappropriate or harmful, while START identifies useful treatments that may have been missed; clinical judgment is still required.

\medterm{Storage Diseases} A group of inherited metabolic disorders characterized by accumulation of substrates
within cells due to specific enzyme deficiencies. Lysosomal storage diseases result from
defects in lysosomal hydrolytic enzymes, leading to accumulation of undigested
macromolecules within lysosomes.


\medterm{Stork Bite} A stork bite is a common harmless pink birthmark, usually a nevus simplex on the forehead, eyelids, nose, or neck, that often fades during childhood.

\medterm{Strabismus} A condition in which the eyes are not aligned and do not point in the same direction. One eye may turn inward, outward, upward, or downward. It may be constant or come and go. Strabismus can cause double vision in adults and may lead to amblyopia, or reduced vision in one eye, when it begins during childhood. Common types include esotropia, exotropia, hypertropia, and hypotropia.
\medterm{Strain} An injury from overstretching or tearing a muscle or tendon, or a force applied to tissue.
\medterm{Strain Imaging} Ultrasound or magnetic-resonance imaging that measures how much tissue deforms during movement, especially how heart muscle shortens or thickens with each heartbeat. It can reveal early weakness even when the total amount of blood pumped appears normal.

\medterm{Strain Neck (Whiplash)} Alternate index label for Whiplash.

\medterm{Strains} A strain is overstretching or tearing of a muscle or tendon, causing pain, tenderness, swelling, weakness, or reduced function after force or overuse.

\medterm{Strangulated Hernia} A hernia in which the blood supply to the trapped tissue is cut off, causing ischaemia and possible tissue death; it is a surgical emergency.

\medterm{Strangulation} Compression of a body structure, especially the neck or a hernia, that compromises blood flow or airway.
\medterm{Strapping} Application of adhesive tape or bandage to support, protect, or restrict movement of a body part.
\medterm{Strategies For Getting Through Labor} Practical approaches for coping with labor, including breathing, movement, position changes, support people, relaxation, analgesia, and communication with the maternity team.

\medterm{Stratum Corneum} The outermost skin layer, made mainly of flattened dead cells filled with keratin. It forms a protective barrier that limits water loss and helps block chemicals, germs, and physical injury.

\medterm{Strawberry Hemangioma} Alternate index label; see \textit{Hemangioma}.

\medterm{Strength} The ability of a muscle or group of muscles to produce force. Healthcare professionals may test strength by comparing movements against resistance; reduced strength can result from pain, poor effort, muscle disease, nerve problems, or brain injury.

\medterm{Strength Testing} Assessment of how much force a muscle or muscle group can produce, using examination or equipment. It helps identify weakness and monitor recovery, progression, or response to treatment.

\medterm{Strength Training} Exercise against resistance to improve muscle force, endurance, power, bone health, and function. The program should be progressed safely according to ability and medical needs.

\synonyms
Resistance training; weight training

\medterm{Strep Throat} A bacterial infection of the throat and tonsils caused by group A Streptococcus, also called Streptococcus pyogenes. It commonly causes a sudden sore throat, pain when swallowing, fever, and swollen or red tonsils, sometimes with white patches.
\medterm{Streptococcal Cellulitis With Lymphangitis} A streptococcal skin infection with inflammation spreading along nearby lymphatic vessels. Redness, warmth, pain, swelling, fever, and red streaks may occur and require prompt antibiotic treatment.
\medterm{Streptococcal Infection} An infection caused by bacteria of the Streptococcus group. Depending on the species and site, it may cause strep throat, scarlet fever, impetigo, cellulitis, pneumonia, heart-valve infection, bloodstream infection, or severe tissue infection.

\medterm{Streptococcal Infections} Infections caused by Streptococcus bacteria, which may affect the throat, skin, lungs, heart valves, blood, or other tissues. Diagnosis and antibiotic choice depend on the site and species; rapidly worsening illness, severe pain, or breathing difficulty is urgent.

\medterm{Streptococcal Pharyngitis} Alternate index label for Strep Throat.

\medterm{Streptococcal Pharyngitis (Strep Throat)} Alternate index label for Strep Throat.

\medterm{Streptococcal Screen} A streptococcal screen looks for group A Streptococcus, usually from a throat swab, to help diagnose strep throat.

\medterm{Streptococcal Toxic Shock Syndrome} A fulminant, life-threatening invasive infection caused by Streptococcus pyogenes (group A Streptococcus) producing streptococcal pyrogenic exotoxins that act as superantigens. It is characterized by the sudden onset of hypotension, shock, and multiorgan failure, frequently in association with severe invasive soft-tissue infections such as necrotizing fasciitis; it carries an exceptionally high mortality rate.

\synonyms
STSS; Streptococcal TSS; Toxic Shock-Like Syndrome

\medterm{Streptococcus} A group of round bacteria that can cause throat infections, pneumonia, skin infections, heart-valve infection, or serious bloodstream illness. Different species spread and respond to antibiotics differently, so testing may guide treatment.
\medterm{Streptococcus Agalactiae} Group B Streptococcus, a bacterium that can live in the bowel or genital tract without causing symptoms. It can pass to a baby during birth and cause serious newborn infection, so pregnancy screening and preventive antibiotics are used when indicated.

\medterm{Streptococcus Pneumoniae (Pneumococcus)} A gram-positive bacterium that can cause pneumonia, meningitis, sinusitis, otitis media, and bloodstream infection; vaccines reduce disease risk.

\medterm{Streptococcus Pyogenes} Group A Streptococcus, a bacterium that causes strep throat, scarlet fever, impetigo, cellulitis, and sometimes severe invasive infection. Untreated infection can occasionally lead to rheumatic fever or kidney inflammation.

\medterm{Streptokinase} A thrombolytic protein that activates plasminogen and dissolves clots; use is limited by bleeding risk and alternatives.
\medterm{Streptomycin} An aminoglycoside antibiotic used in selected tuberculosis and other infections.
\medterm{Stress} A physical and emotional response to perceived demands or threats. It can change heart rate, breathing, muscle tension, attention, sleep, and digestion; persistent or overwhelming stress may affect health.


\medterm{Stress Cardiomyopathy} A temporary weakening of the heart muscle that often follows severe emotional or physical stress. It can cause chest pain, breathlessness, or fainting and may resemble a heart attack, even when no heart artery is blocked.

\medterm{Stress Echocardiography} An ultrasound scan of the heart performed while the heart works harder through exercise or medicine. It looks for changes in heart-muscle movement that may suggest reduced blood flow.
\medterm{Stress Fracture} A small crack in a bone caused by repeated activity or loading. Pain usually develops gradually in one spot and worsens with activity; continuing to stress the bone can make the crack larger or complete.

\medterm{Stress Fractures} Small cracks in bone caused by repeated loading that exceeds the bone’s ability to repair itself. They commonly affect the feet or legs and cause localized pain that worsens with activity.

\medterm{Stress Incontinence} Involuntary leakage of urine during activities that increase intra-abdominal pressure, such as coughing, sneezing, lifting, or exercise.

\medterm{Stress Injuries} Tissue injuries caused by repeated loading or overuse without sufficient recovery. Examples include stress fractures, tendinopathy, and muscle strain.

\medterm{Stress Response} Body changes that prepare a person to respond to a perceived threat or challenge. Hormones and nerve signals alter heart rate, breathing, blood flow, attention, digestion, immunity, and energy use.

\synonyms
Fight-or-flight response

\medterm{Stress Shielding} Loss or weakening of bone next to an implant because the implant carries more force than the bone normally would. Over time, this can reduce bone density and affect how securely the implant stays in place.

\medterm{Stress Tests} Diagnostic tests that assess organ function under controlled physical, pharmacologic, or physiologic stress, commonly evaluating the heart for inducible ischemia or abnormal rhythm.

\medterm{Stress Ulcer Prophylaxis} Measures used to reduce the risk of stomach or duodenal bleeding in selected critically ill people. The need for treatment depends on bleeding risk, other medicines, and the person's overall condition.

\medterm{Stress Urinary Incontinence} Leakage of urine when coughing, sneezing, laughing, lifting, exercising, or otherwise increasing pressure inside the abdomen. It commonly results from weakness of the pelvic-floor muscles or the bladder outlet and can often be treated.
\medterm{Stressors} Events, conditions, or circumstances that trigger a physical or emotional stress response. Examples include illness, pain, sleep loss, trauma, work demands, and relationship difficulties.

\medterm{Stretch Marks} Linear bands of atrophic skin, also called striae, caused by stretching and hormonal or corticosteroid effects; they commonly occur during pregnancy, growth, weight change, or steroid exposure.

\medterm{Stretch Reflex} An automatic tightening of a muscle when it is suddenly stretched, such as the lower-leg movement after tapping the knee tendon. It helps maintain posture and provides information about nerve and spinal-cord function.

\medterm{Stria} A line, streak, or band in tissue, such as a stretch mark. Striae may result from skin stretching, growth, hormone changes, medicines, or underlying disease and can fade over time.

\synonyms
Striae

\medterm{Striae} Linear bands or streaks in the skin, commonly called stretch marks. They develop when deeper skin layers stretch or change rapidly during growth, pregnancy, weight change, muscle gain, or hormone treatment.
\medterm{Striae Atrophicae} Linear bands of thin, wrinkled skin caused by stretching and disruption of dermal connective tissue; commonly called stretch marks.

\synonyms
stretch marks.

\medterm{Striae Gravidarum} Stretch marks that develop during pregnancy, commonly on the abdomen, breasts, thighs, or
buttocks. They may begin as red or purple lines and gradually fade to lighter bands.

\medterm{Stricture (Stenosis)} An abnormal narrowing of a passage in the body, such as the esophagus, intestine, bile duct, or urethra. Scarring, inflammation, injury, surgery, or cancer may cause blockage and symptoms related to the affected organ.

\medterm{Strictures Esophagus (Schatzki Ring)} Alternate index label for Schatzki Ring.

\medterm{Stridor} A harsh, high-pitched sound, usually heard during breathing in, caused by narrowing of the upper airway. New or severe stridor is an emergency because airflow may become blocked.
\medterm{String Test For Intestinal Parasites} A string test collects a sample from the upper digestive tract on a swallowed capsule and string to look for selected intestinal parasites.

\medterm{Stripping Of Membranes (Membrane Sweep)} A procedure performed during a cervical exam where the examiner inserts a finger through
the internal os and sweeps the membranes circumferentially, separating them from the
lower uterine segment. Releases endogenous prostaglandins, promoting cervical ripening
and potentially initiating labor.

\synonyms
Membrane Sweep

\medterm{Stroke} Brain injury caused by interrupted blood flow to part of the brain (ischemic stroke) or bleeding in or around the brain (hemorrhagic stroke). It may cause sudden facial drooping, arm or leg weakness or numbness, speech trouble, vision loss, severe headache, dizziness, or confusion. The effects vary according to the brain area involved and the extent of injury.
\medterm{Stroke Rehabilitation} Rehabilitation after a stroke helps improve movement, speech, swallowing, thinking, mood, and daily activities while reducing complications and preventing another stroke.
\medterm{Stroke Index} A measure of the amount of blood pumped by the heart with each beat, adjusted for body-surface area. It helps assess heart function and is interpreted with other examination, imaging, and hemodynamic findings.

\medterm{Stroke Risk Factors} Stroke risk factors include high blood pressure, smoking, diabetes, high cholesterol, atrial fibrillation, prior stroke, and certain blood or vascular disorders.


\medterm{Stroke Volume} The volume of blood ejected by the left ventricle with each heartbeat, normally 60-100
mL. It is determined by three factors: preload, afterload, and myocardial contractility.
Stroke volume equals end-diastolic volume minus end-systolic volume. Reduced stroke
volume occurs in heart failure, hypovolemia, and cardiogenic shock.

\medterm{Stroma} The supporting tissue of an organ that surrounds and holds its working cells. It may contain connective tissue, blood vessels, nerves, and immune cells; the exact structure and role vary between organs.
\medterm{Strongyloides Stercoralis} A soil-transmitted roundworm that can enter through skin and live in the human intestine. Infection may cause rash, abdominal symptoms, or no symptoms, but severe hyperinfection can occur when immunity is weakened.

\medterm{Strongyloidiasis} An intestinal worm infection that can spread throughout the body and become life-threatening when immunity is weakened.
\medterm{Structured Reporting} Reporting that uses a standard format, headings, and required fields to organize findings. It can improve completeness, consistency, and communication in imaging, pathology, and other clinical reports, but does not replace professional judgment.

\medterm{Struma Ovarii} A specialized mature ovarian teratoma composed predominantly of thyroid tissue; it may be hormonally active or develop malignant thyroid-type transformation.

\medterm{Struvite Stones} Kidney or urinary-tract stones formed during infection with certain bacteria. They can grow into large branching stones that fill the kidney and may damage it or block urine flow.

\medterm{Strychnine} A poison causing severe painful muscle spasms and convulsions while awareness may remain; emergency treatment is required.
\medterm{Student's Elbow} Olecranon bursitis, an inflammation of the bursa over the point of the elbow, often related to repeated pressure or trauma.


\medterm{Stuffy Nose} Blocked or difficult nasal breathing caused by swollen tissue, mucus, allergy, infection, irritants, or structural narrowing.

\medterm{Stuffy Or Runny Nose - Children} A stuffy or runny nose in a child is usually caused by a viral infection or allergy; breathing difficulty, dehydration, prolonged fever, or symptoms in an infant need prompt assessment.

\medterm{Stuffy Or Runny Nose – Adult} A stuffy or runny nose in an adult commonly results from viral infection, allergy, irritants, or sinus disease; persistent, severe, bloody, or one-sided symptoms need review.

\medterm{Stump Pain} Alternate index label; see \textit{Residual limb pain}.

\medterm{Stupor} A state of reduced consciousness in which the patient is unresponsive to most environmental stimuli but can be temporarily aroused by vigorous, repeated stimulation (painful stimuli, loud shouting).
\medterm{Sturge-Weber Syndrome} A condition present from birth involving abnormal blood vessels in the skin and sometimes on the surface of the brain. It may cause a port-wine birthmark, seizures, weakness, headaches, developmental problems, or glaucoma, depending on which areas are affected.

\medterm{Stuttering} A speech-flow problem involving repeated sounds or words, stretched sounds, pauses, or blocks when speaking. It is not caused by low intelligence; support may include speech therapy, communication strategies, and help with anxiety or teasing.
\medterm{Sty} Alternate index label for Stye.

\medterm{Stye} A tender, localized infection or inflammation of an eyelid gland, usually producing a painful red lump near the eyelid margin. It is also called a hordeolum.


\medterm{Styptics} Substances that help stop minor surface bleeding by tightening tissue, narrowing small blood vessels, or supporting clot formation. They are used for limited bleeding and do not replace emergency care for severe blood loss.
\medterm{Subacute} Developing or lasting longer than an acute condition but not as long as a chronic one. The exact time range varies with the disease.
\medterm{Subacute Bacterial Endocarditis} A slowly developing infection of the inner lining or valves of the heart, usually caused by bacteria entering the bloodstream. It may cause fever, tiredness, weight loss, a new heart murmur, or embolic complications and requires prolonged antibiotic treatment.

\medterm{Subacute Combined Degeneration} Damage to parts of the spinal cord, usually from vitamin B12 deficiency. It can cause numbness, weakness, poor balance, and difficulty walking.

\medterm{Subacute Combined Degeneration Of The Cord} Damage to the spinal cord, usually caused by vitamin B12 deficiency, affecting nerve pathways for vibration, position sense, and movement. It may cause tingling, weakness, poor balance, or difficulty walking and can become permanent without treatment.
\medterm{Subacute Rehabilitation} Rehabilitation provided after the early phase of illness or injury for people who still need nursing care and therapy but no longer need intensive hospital treatment. It may occur in a specialized facility or outpatient setting.

\medterm{Subacute Sclerosing Panencephalitis} A rare, progressive brain disease caused by measles virus remaining in the brain years after the original infection. It can cause behavior or learning changes, jerking movements, seizures, loss of coordination, and severe disability; prevention depends on measles vaccination.
\medterm{Subacute Thyroiditis} Temporary inflammation of the thyroid, often after a viral illness, causing painful neck swelling, fever, and tiredness. Thyroid hormone levels may first be high and later low before returning to normal; treatment relieves pain and other symptoms.

\medterm{Subaortic Stenosis} Narrowing just below the aortic valve that obstructs blood leaving the heart. It may cause a heart murmur, chest pain, fainting, breathlessness, or damage to the aortic valve and sometimes requires surgery.

\medterm{Subarachnoid} Situated beneath the arachnoid, one of the thin coverings around the brain and spinal cord. The subarachnoid area contains protective fluid and blood vessels; bleeding there can cause a sudden severe headache and is an emergency.

\medterm{Subarachnoid Haemorrhage} Bleeding into the space around the brain containing cerebrospinal fluid, often causing sudden severe headache.
\medterm{Subarachnoid Hemorrhage} Bleeding into the space around the brain, often from a ruptured blood-vessel bulge. It typically causes a sudden, extremely severe headache and may cause vomiting, neck stiffness, confusion, weakness, or loss of consciousness.

\medterm{Subarachnoid Space} The space between two thin coverings of the brain and spinal cord, called the arachnoid and pia mater. It contains cerebrospinal fluid and blood vessels and helps cushion the nervous system.

\medterm{Subareolar Abscess} A pus-filled infection beneath the areola, often associated with periductal inflammation and smoking, causing a painful lump, redness, drainage, or recurrent fistula.

\medterm{Subchorionic Hemorrhage} Bleeding between the chorion and the uterine wall caused by partial separation of the gestational sac membranes. It may produce vaginal bleeding in early or mid-pregnancy and is assessed according to bleeding, gestational age, and ultrasound findings.

\medterm{Subclavian} Relating to structures beneath the clavicle, especially the subclavian arteries and veins that supply or drain the upper limbs.
\medterm{Subclavian Artery} A major artery beneath the collarbone that supplies the arm and contributes blood flow to the neck, chest, and brain. Narrowing or blockage can cause arm pain, weakness, or symptoms from reduced blood flow to the brain.

\medterm{Subclavian Steal Syndrome} A condition in which narrowing of an artery supplying the arm redirects blood away from the back of the brain, especially during arm exercise. It may cause arm pain or weakness, dizziness, fainting, blurred vision, or imbalance. Treatment controls vascular risk factors and may restore blood flow when symptoms are significant.


\medterm{Subclavian Vein} The continuation of the axillary vein beyond the lateral border of the first rib. It
joins the internal jugular vein behind the sternoclavicular joint to form the
brachiocephalic vein.

\medterm{Subclavian Vein Thrombosis} A blood clot in the large vein beneath the collarbone. It may follow a catheter, repeated strenuous arm activity, or compression near the shoulder and can cause arm swelling, pain, bluish skin, or a dangerous clot in the lungs.

\medterm{Subclinical} Present without clear symptoms or signs that a person or clinician can notice. A subclinical problem may be found through a blood test, scan, screening examination, or other test and may or may not later cause symptoms.
\medterm{Subclinical Disease} Disease that produces no obvious symptoms or signs but can be detected by screening, examination, or laboratory testing. It may remain mild or later become clinically apparent, depending on the condition.

\medterm{Subclinical Hypothyroidism} A blood-test pattern showing that the thyroid is beginning to underperform: the pituitary sends extra stimulation, but thyroid-hormone levels remain normal. Many people have no symptoms; others have mild tiredness, cold intolerance, or constipation.

\medterm{Subclinical Infection} An infection that produces no noticeable symptoms or signs but may still be detectable by testing and may sometimes be transmissible.

\medterm{Subconjunctival Hemorrhage} Bleeding beneath the conjunctiva that appears as a sharply demarcated red patch on the white of the eye; it is usually harmless but recurrent cases or trauma require assessment.

\medterm{Subconjunctival Hemorrhage (Broken Blood Vessel In Eye)} A small blood vessel broken beneath the clear surface of the eye, causing a sharply demarcated red patch that usually resolves without treatment.

\synonyms
Broken Blood Vessel In Eye

\medterm{Subconscious} Mental processes outside immediate awareness; it is a psychological term rather than a specific diagnosis.
\medterm{Subcorneal Pustular Dermatosis} A rare recurring skin disorder causing shallow, noninfectious pus-filled blisters just beneath the outer skin layer, usually on the trunk or skin folds. It may resemble psoriasis and can be associated with blood or immune disorders.

\medterm{Subcu} Subcu is an abbreviation for subcutaneous, meaning beneath the skin; clear clinical instructions should spell out the route.

\medterm{Subculture} The transfer of microorganisms from one culture to fresh culture medium, used to
maintain pure cultures, isolate colonies, increase viability, prepare standardized
inocula, and perform identification/susceptibility testing. On solid media, subculture
aims to obtain isolated single colonies from a mixed or stock culture.

\medterm{Subcutaneous} Subcutaneous means beneath the skin; injections by this route deliver medicine into the fatty tissue between skin and muscle.

\medterm{Subcutaneous Injections} Administration of a medicine into the fatty tissue beneath the skin, commonly used for insulin, anticoagulants, and selected biologic medicines.

\medterm{Subcutaneous Emphysema} Air trapped beneath the skin, often after chest injury, surgery, a punctured lung, or a breathing procedure. It may cause swelling and a crackling feeling; treatment addresses the air leak, and breathing difficulty or rapid spread requires urgent care.

\medterm{Subcutaneous Implantable Cardioverter-Defibrillator} A device placed under the skin that detects and treats dangerous fast heart rhythms with an electrical shock. Its wires run outside the heart and blood vessels; infection, inappropriate shocks, movement, or device failure are possible.

\medterm{Subcutaneous Layer} The tissue layer beneath the dermis, made mainly of fat and connective tissue. It cushions deeper structures, insulates the body, stores energy, and contains larger blood vessels and nerves.

\medterm{Subcutaneous Tissue} The layer beneath the skin that connects it to deeper structures. It contains fat, connective tissue, larger blood vessels, and nerves and provides insulation, cushioning, and energy storage.

\medterm{Subdural} Located beneath the dura, one of the protective coverings around the brain and spinal cord. A subdural hematoma is bleeding in this space, often after head injury, and may cause headache, vomiting, confusion, weakness, seizures, or loss of consciousness. Significant bleeding needs urgent assessment.
\medterm{Subdural Effusion} Accumulation of usually clear fluid in the subdural space, often after trauma, surgery, infection, or cerebrospinal-fluid disturbance and sometimes causing headache or pressure effects.

\medterm{Subdural Haematoma} A collection of blood between protective coverings of the brain, usually after small veins tear during a head injury. It may develop suddenly or gradually and cause headache, confusion, drowsiness, weakness, seizures, or loss of consciousness.

\medterm{Subdural Hematoma} A collection of blood between two protective coverings of the brain, usually after small veins tear during a head injury. It can cause headache, confusion, weakness, seizures, or sleepiness and may require urgent treatment.

\medterm{Subdural Hemorrhage} Bleeding between two protective coverings of the brain, usually after tearing of small veins during a head injury. It may cause sudden or gradually worsening headache, confusion, drowsiness, weakness, seizures, or loss of consciousness and is diagnosed with brain imaging.
\medterm{Subendocardial Myocytes} Heart-muscle cells in the inner layer of the myocardium, which are especially vulnerable to reduced blood flow.

\medterm{Subependymal Giant Cell Astrocytoma} A usually slow-growing brain tumor that develops near a passage where fluid leaves the brain’s cavities. It is strongly associated with tuberous sclerosis and can block fluid flow, causing headache, vomiting, or enlarged brain cavities.

\medterm{Subependymoma} A rare, usually benign slow-growing ependymal tumor, commonly arising in the ventricular system of adults. Symptoms result from obstruction of cerebrospinal-fluid pathways or incidental detection.

\medterm{Subfertility} Alternate index label; see \textit{Infertility}.

\medterm{Subgaleal Hemorrhage} A potentially life-threatening accumulation of blood in the space beneath the scalp's
epicranial aponeurosis. It can expand rapidly in a newborn, particularly after vacuum-assisted delivery, and requires
urgent recognition and treatment.

\medterm{Subinvolution} Slower-than-expected return of an organ to its usual size, especially the uterus after childbirth. It may cause prolonged or heavy bleeding, pelvic pain, or an enlarged tender uterus and needs assessment for retained tissue, infection, or other causes.
\medterm{Subjacent} An anatomical term describing a structure, tissue, or lesion located beneath another structure. It indicates relative position and does not necessarily mean that the two structures are touching; the relationship may be described on physical examination, imaging, or during surgery.

\medterm{Subjective} Based on a person's own experience or report rather than a directly measured finding. Pain, nausea, fatigue, dizziness, and mood are subjective symptoms, while temperature and blood pressure are objective measurements.
\medterm{Sublimation} A psychological defense process in which unacceptable impulses are redirected into socially acceptable activity; in chemistry it is solid-to-gas transition.
\medterm{Sublingual} Sublingual means beneath the tongue; medicines given this way can be absorbed through oral mucosa and may bypass much first-pass liver metabolism.

\medterm{Sublingual Epidermoid Cyst} A benign cyst lined by stratified squamous epithelium beneath the tongue mucosa,
presenting as a slow-growing midline swelling that may require surgical excision if
symptomatic.

\medterm{Sublingual Gland} One of the paired major salivary glands, located beneath the tongue. It produces mainly mucous saliva that enters the
mouth through several small ducts and helps lubricate the oral cavity.

\medterm{Subluxation} Partial displacement of a joint in which some contact between the joint surfaces remains. It may cause pain, swelling, instability, limited movement, nerve injury, or repeated episodes after injury or disease.
\medterm{Submandibular Gland} One of the major salivary glands, located beneath the lower jaw. It makes much of the saliva present when a person is not eating and releases it into the floor of the mouth through a small duct.

\medterm{Submucosa} The supportive layer beneath the moist lining of organs such as the intestine, airway, and bladder. It contains blood vessels, nerves, glands, and connective tissue and helps nourish and anchor the lining.
\medterm{Subphrenic Abscess} A collection of pus beneath the diaphragm, usually caused by infection after abdominal surgery, a burst abdominal organ, or spread from nearby tissue. It may cause fever, upper-abdominal or shoulder pain, breathing discomfort, and sepsis and usually needs drainage and antibiotics.
\medterm{Subq} Subq is an abbreviation for subcutaneous, referring to tissue beneath the skin and the route used for some injections.

\medterm{Subscapular Artery} A branch of the axillary artery that supplies the shoulder blade region, chest wall, and nearby muscles. It contributes to the network of vessels supporting the shoulder and upper limb.
\medterm{Substance} A substance is a material with defined or variable chemical composition; medical effects depend on identity, dose, route, duration, and host factors.

\medterm{Substance Abuse} Legacy, nonspecific label; see \textit{Substance Use Disorder}.

\medterm{Substance Abuse Disorder} Legacy label; see \textit{Substance Use Disorder}.

\medterm{Substance Abuse Psychiatry} An older term for the clinical specialty concerned with substance use disorders and co-occurring mental-health conditions. Current terminology generally uses \textit{Addiction Psychiatry} or \textit{Addiction Medicine}.

\medterm{Substance Dependence} A state in which repeated substance use produces physiologic adaptation, so that reducing or stopping the substance may cause withdrawal. It is distinct from, but may occur with, substance-use disorder.

\medterm{Substance Intoxication} A temporary change in thinking, behavior, or body function caused by recent use of alcohol, a drug, or another substance. Effects depend on the substance, amount, timing, combinations, and the person’s health.

\medterm{Substance Use} The consumption of alcohol, tobacco or nicotine, cannabis, prescription or nonprescription drugs, inhalants, or other psychoactive substances. Use may be occasional, prescribed, nonmedical, or associated with a substance use disorder; the term alone does not imply harmful use.

\medterm{Substance Use - Amphetamines} Use of amphetamine-type stimulants, which can be associated with insomnia, anxiety, psychosis, high blood pressure, arrhythmias, overheating, dependence, and overdose.

\medterm{Substance Use - Cocaine} Use of cocaine, a stimulant that can cause dependence, agitation, seizures, heart attack, stroke, dangerous overheating, and severe psychiatric or cardiovascular complications.

\medterm{Substance Use - Inhalants} Inhalant misuse involves breathing fumes or gases from household or industrial products and can cause sudden cardiac arrest, brain injury, organ damage, dependence, or death.

\medterm{Substance Use - LSD} LSD is a potent hallucinogen that alters perception, mood, and thought and can cause panic, accidents, persistent perceptual symptoms, or psychosis in susceptible people.

\medterm{Substance Use - Marijuana} Marijuana or cannabis can impair attention, memory, coordination, and driving and may cause dependence, anxiety, psychosis, or respiratory harm when smoked.

\medterm{Substance Use - Phencyclidine} Phencyclidine is a dissociative drug that can cause agitation, hallucinations, loss of coordination, high blood pressure, seizures, dangerous behavior, or coma.

\medterm{Substance Use - Prescription Drugs} Use of prescription medicines in a way not directed by a clinician, including taking someone else’s medicine, taking extra doses, or using them for intoxication; dependence, overdose, and interactions may result.

\medterm{Substance Use Disorder} A clinically significant pattern of alcohol, tobacco or nicotine, cannabis, prescription or nonprescription drug, inhalant, or other psychoactive-substance use that causes impaired control, social or occupational impairment, risky use, or clinically significant distress. Features may include craving, unsuccessful efforts to reduce or stop use, substantial time spent obtaining or using the substance, failure to meet major responsibilities, continued use despite harm, tolerance, or withdrawal; tolerance and withdrawal are not required.

\medterm{Substance Use In Pregnancy} Use of alcohol, tobacco or nicotine, cannabis, opioids, cocaine, or other substances during pregnancy. Effects depend on the substance, dose, timing, route, and other factors and may include maternal complications, impaired fetal growth, preterm birth, or neonatal withdrawal.


\medterm{Substance Withdrawal} Physical or mental symptoms that occur after reducing or stopping heavy or prolonged substance use. Symptoms range from anxiety and nausea to seizures, severe confusion, dangerous blood-pressure changes, or breathing problems, depending on the substance.

\medterm{Substernal Goiter} An enlarged thyroid gland that extends from the neck into the upper chest. It may press on the windpipe or esophagus, causing shortness of breath, cough, or difficulty swallowing. Imaging helps assess its size and position.

\medterm{Substrate} A substance on which an enzyme acts or a material used by a biochemical pathway.
\medterm{Subungual Exostosis} A benign osteocartilaginous outgrowth arising beneath the nail, usually of a toe or finger. It causes a painful firm subungual mass and nail deformity and is treated by excision when symptomatic.


\medterm{Succenturial Lobe} An accessory lobe of the placenta, connected to the main placenta by blood vessels that
traverse the membranes. If undetected after delivery, the succenturial lobe may be
retained in the uterus, causing postpartum hemorrhage or infection. Diagnosis: careful
inspection of the placenta and membranes after delivery.

\medterm{Succenturiate Placenta} A succenturiate placenta has one or more accessory lobes separate from the main placenta, which can increase risk of retained tissue or fetal-vessel injury.

\medterm{Succinate Dehydrogenase} An enzyme inside mitochondria that helps cells make energy. It changes one small molecule into another while passing electrons through the energy-producing system; inherited or tumor-related changes can have clinical importance.

\medterm{Succinylcholine} A very fast, short-acting medicine that temporarily paralyzes skeletal muscles to help place a breathing tube or perform emergency procedures. It can dangerously raise potassium, stop breathing longer than expected, or trigger malignant hyperthermia in susceptible people.

\medterm{Succumb} To succumb means to die from or be overcome by a disease, injury, exposure, or other serious process.

\medterm{Sucking Blister} A blister or worn area of skin in a newborn, usually on a hand, finger, forearm, or lip, caused by
repeated sucking before or during feeding. It is usually harmless and heals on its own.
\medterm{Sucking Reflex} A primitive reflex present from birth, elicited by placing an object in the infant's
mouth. The infant rhythmically sucks on the object. Coordinated sucking and swallowing
typically develop by 34-36 weeks gestational age. Premature infants may lack coordinated
suck-swallow, requiring gavage feeding.

\medterm{Sucrose} Common table sugar, a disaccharide made of glucose and fructose. Digestive enzymes split it into these simpler sugars for absorption.
\medterm{Sucrose Lysis Test} An older blood test once used to look for paroxysmal nocturnal hemoglobinuria, a disorder in which red blood cells are unusually vulnerable to destruction. Flow-cytometry testing has largely replaced it because it is more specific and reliable.

\medterm{Suction} Removal of fluid, secretions, blood, or other material using negative pressure. It may be used in airway care, surgery, wound care, or drainage.
\medterm{Suction (Vascular)} A suction device used during a vascular procedure to remove blood or other fluid from the operative
field.

\synonyms
Vascular suction device

\medterm{Suction Device} Equipment that creates negative pressure to remove fluid, blood, mucus, or other material. It usually includes a pump, tubing, collection container, and suction catheter or tip and is used in airway care, surgery, or wound drainage.

\medterm{Suction Tip} The end attachment of a suction device used to remove blood, fluid, secretions, or debris.
Common types include the Yankauer tip for oropharyngeal suction and surgical suction tips such as the Frink tip.

\medterm{Sudden Cardiac Death} Unexpected death from a cardiac cause, usually occurring rapidly after symptom onset or without
witnessed symptoms. Causes include coronary disease, cardiomyopathy, inherited arrhythmia
syndromes, and bradyarrhythmia; prevention depends on the underlying risk.

\synonyms
Arrest Cardiac (Sudden Cardiac Death)
Death Sudden Cardiac (Sudden Cardiac Death)

\medterm{Sudden Infant Death Syndrome} The unexplained sudden death of an infant younger than one year after a complete investigation, including medical history, death-scene review, and autopsy. Safe sleep on the back on a firm, clear surface and avoiding smoke reduce risk, but do not prevent every case.

\medterm{Sudden Sensorineural Hearing Loss} A rapid loss of hearing caused by a problem in the inner ear or hearing nerve, often noticed on waking or over a few days. It may be accompanied by ringing or dizziness and needs urgent assessment because early treatment can improve recovery in some cases.

\medterm{Sudden Unexpected Death In Infancy} The sudden, unexpected death of an infant that requires investigation of the medical history, circumstances, death scene, and sometimes autopsy. The cause may be identified, or the death may meet the definition of sudden infant death syndrome.
\medterm{Sudden Unexpected Infant Death} The sudden and unexpected death of an infant younger than 1 year, whether or not the
cause is known before investigation. Sudden infant death syndrome is the unexplained
subtype after complete scene, history, and autopsy evaluation.

\medterm{Sudek's Atrophy} An older name for complex regional pain syndrome with pain, swelling, autonomic change, and regional bone loss.
\medterm{Sudorifics (Diaphoretics)} Medicines or other treatments that increase sweating. They have limited specific medical uses; excessive sweating caused by a treatment can lead to dehydration, overheating, or abnormal body salts.
\medterm{Suffocation} Inability to breathe enough because the airway is blocked, oxygen is unavailable, the chest is compressed, or another problem prevents normal breathing. It quickly causes unconsciousness and death without immediate rescue.
\medterm{Sugammadex} A medicine used during anesthesia to quickly reverse muscle paralysis caused by selected muscle-relaxing medicines. It binds those medicines so the body can remove them; allergic reactions and changes in heart rate are possible.

\medterm{Sugar-Water Hemolysis Test} An older laboratory test assessing red-cell resistance to hypotonic sugar solution, historically used in evaluating paroxysmal nocturnal hemoglobinuria and now largely replaced by flow-cytometric testing.

\medterm{Suicidal Ideation} Thoughts about dying or killing oneself, with or without a plan or intention to act. Any current plan, intent, access to a method, recent attempt, or inability to stay safe is an emergency requiring immediate support and professional help.

\medterm{Suicide} Death caused by an intentional action to end one's life. Suicidal thoughts, plans, or behavior are urgent medical concerns; immediate support, safety measures, and professional assessment can help protect the person.
\medterm{Suicide And Suicidal Behavior} Thoughts, plans, attempts, or acts intended to cause one’s own death; any current intent, plan, access to means, or inability to stay safe requires immediate emergency support.

\medterm{Suicide Risk Assessment} A clinical evaluation of suicidal thoughts, intent, plan, access to means, past
attempts, self-harm, psychiatric and medical illness, substance use, protective factors,
and immediate safety. It informs but does not replace ongoing clinical judgment and
safety planning.

\medterm{Sulcus} A groove or shallow furrow on an organ or body surface, especially the brain or heart. Brain sulci separate raised areas and increase the cerebral cortex's folded surface area.
\medterm{Sulfa Drugs} A group of medicines, especially sulfonamide antibiotics, that block a chemical pathway bacteria need to grow. They treat selected infections but can cause allergic reactions and other side effects; the term may also refer to nonantibiotic sulfonamides.

\medterm{Sulfamethoxazole} A sulfonamide antibiotic used commonly with trimethoprim in selected infections.
\medterm{Sulfasalazine} An anti-inflammatory medicine used for inflammatory bowel disease and rheumatoid or other inflammatory arthritis. It changes immune activity and can cause nausea, rash, low blood counts, liver problems, or kidney problems.
\medterm{Sulfatide} A fatty substance containing sulfate found mainly in myelin, the insulating covering of nerve fibers. Abnormal buildup or deficiency can disrupt nerve function and occurs in some inherited metabolic disorders.

\medterm{Sulfite Sensitivity} Sulfite sensitivity is an adverse reaction to sulfite preservatives, sometimes causing wheeze, flushing, headache, or other symptoms, especially in people with asthma.

\medterm{Sulfonamide} A medicine containing a sulfonamide chemical group. Some are antibiotics that stop bacteria making folate, which they need to grow; others reduce inflammation or increase urine flow. Allergic reactions and other side effects are possible.

\medterm{Sulfonamides} Sulfonamides are drugs containing a sulfonamide group, including antibacterial agents that inhibit bacterial folate synthesis; allergic reactions, severe skin reactions, cytopenias, renal effects, and crystalluria are possible.

\medterm{Sulfonylurea} An oral diabetes medicine that makes the pancreas release more insulin. It can lower blood sugar but may cause dangerously low blood sugar, hunger, or weight gain.

\medterm{Sulfonylureas} Oral medicines for type 2 diabetes that stimulate the pancreas to release more insulin. They can lower blood glucose but may cause low blood sugar, hunger, or weight gain.

\medterm{Sulfur Colloid Scan} A nuclear-medicine scan in which a small radioactive tracer is used to show activity in the liver, spleen, bone marrow, or lymph nodes. It can help assess how these organs work or identify abnormal tissue distribution.

\medterm{Sulfuric Acid Poisoning} Severe poisoning from sulfuric acid, a highly corrosive chemical. It can cause deep burns of the skin, eyes, mouth, throat, and digestive tract and dangerous fumes; emergency care is required and vomiting must not be induced.

\medterm{Sulindac Overdose} Taking too much sulindac, a nonsteroidal anti-inflammatory drug, may cause nausea, vomiting, stomach bleeding, drowsiness, kidney injury, liver injury, seizures, or metabolic acidosis. A suspected overdose needs prompt medical or poison-control advice.

\medterm{Sulphonamides} A group of medicines, including some antibiotics, that block a chemical pathway bacteria need to grow. They are used for selected infections and other conditions; allergic reactions, interactions, and blood or kidney problems can occur.

\medterm{Sulphonylureas} Oral medicines that stimulate pancreatic insulin release to lower blood glucose in type 2 diabetes.
\medterm{Sumatriptan} A migraine medicine that narrows certain blood vessels and changes nerve signaling to stop acute migraine or cluster headaches.
\medterm{Summer Diarrhoea} An older term for seasonal diarrheal illness, often referring to infectious diarrhea in children.
\medterm{Sun Poisoning} Alternate index label; see \textit{Polymorphous light eruption}.

\medterm{Sun Protection} Sun protection reduces ultraviolet exposure through shade, protective clothing, and broad-spectrum sunscreen and helps prevent sunburn, premature skin aging, and skin cancer.

\medterm{Sun Protection (Photoprotection)} Measures that reduce exposure to ultraviolet radiation, including shade, protective clothing, sunglasses, and broad-spectrum sunscreen. They help prevent sunburn, premature skin aging, eye damage, and ultraviolet-related skin cancers.

\synonyms
Photoprotection

\medterm{Sun-Sensitive Drugs (Photosensitivity To Drugs)} Alternate index label for Photosensitivity to Drugs.


\medterm{Sunburn} An acute inflammatory injury of the skin caused by excessive ultraviolet radiation, resulting in redness, pain, warmth, and sometimes blistering or peeling.
\medterm{Sundowning} Late-day or evening worsening of confusion, agitation, wandering, or behavioral symptoms in some people with dementia or delirium, influenced by fatigue, light, routines, and unmet needs.

\medterm{Sunflower Syndrome} A rare epilepsy in which a person repeatedly turns toward bright light and waves a hand near the eyes, often triggering absence seizures. It usually begins in childhood and requires neurologic assessment and seizure treatment.

\medterm{Sunken Chest} Alternate index label; see \textit{Pectus excavatum}.

\medterm{Sunstroke} A life-threatening form of heat illness in which body temperature rises dangerously high and the nervous system is impaired.

\medterm{Superantigen} A substance made by some bacteria that activates many immune cells at once. This can cause an excessive release of inflammatory chemicals and contribute to illnesses such as toxic shock syndrome or some forms of scarlet fever.




\medterm{Superficial} Near the surface of the body or an organ, rather than deep inside it. The term may describe tissues, veins, burns, wounds, or infections that mainly affect the skin and tissue just beneath it.
\medterm{Superficial Fascia} The soft tissue layer just beneath the skin. It contains fat, small blood vessels, nerves, and connective tissue and helps insulate the body and allow the skin to move.

\synonyms
subcutaneous tissue.

\medterm{Superficial Spreading Melanoma} The most common cutaneous melanoma subtype, characterized initially by radial growth across the epidermis and an irregularly pigmented patch or plaque before invasion develops.

\medterm{Superficial Thrombophlebitis} Inflammation and a blood clot in a vein just beneath the skin, causing a tender, firm, red cord. It often improves with treatment, but a spreading clot or symptoms of deep-vein thrombosis require medical assessment.

\medterm{Superimposed Preeclampsia} Preeclampsia developing in a person who already has chronic hypertension. It may
present with new or worsening proteinuria, thrombocytopenia, abnormal liver or kidney tests, neurologic symptoms, or
other features of maternal organ dysfunction.

\medterm{Superinfection} A new infection that develops during or after treatment of a previous infection, often
involving an antimicrobial-resistant organism or an organism not susceptible to the
original therapy, such as Clostridioides difficile after antibiotic exposure.

\medterm{Superior} Above or toward the head in the body, or higher than another structure. For example, the neck is superior to the chest. The opposite direction is inferior, meaning below or toward the feet.

\medterm{Superior Mesenteric Artery} A major artery branching from the abdominal aorta. It supplies blood to much of the small intestine and to parts of the large intestine, including the right colon and nearby transverse colon.

\medterm{Superior Mesenteric Artery Syndrome} Pressure on part of the small intestine where it passes between two major arteries, causing blockage. It may follow major weight loss and cause fullness, upper-abdominal pain, nausea, and vomiting; treatment restores nutrition and relieves the obstruction.

\medterm{Superior Semicircular Canal Dehiscence} A gap or very thin area in the bone covering one balance canal in the inner ear. Loud sounds or pressure may cause dizziness, imbalance, hearing changes, or hearing one’s own voice unusually loudly; diagnosis uses symptoms, hearing tests, and imaging.

\medterm{Superior Vena Cava} A large vein in the chest that carries oxygen-poor blood from the head, neck, arms, and upper chest to the right side of the heart.

\medterm{Superior Vena Cava Syndrome} Blockage or compression of the large vein carrying blood from the upper body to the heart. It can cause swelling of the face, neck, or arms, visible neck veins, cough, or shortness of breath and may require urgent treatment.
\synonyms
SVCS (Superior Vena Cava Syndrome)

\medterm{Superior Vena Caval Obstruction} Blockage or compression of the large vein returning blood from the head, neck, and arms to the heart. It is often caused by a chest tumor or a blood clot around a catheter and may cause facial or arm swelling, prominent neck veins, cough, or breathlessness.

\medterm{Supernumerary Nipples} Supernumerary nipples are extra nipples along the embryonic milk line, usually harmless but occasionally associated with underlying breast or urinary-tract developmental anomalies.

\medterm{Supernumerary Rib} An extra rib beyond the typical twelve pairs, most commonly a cervical rib above the
first thoracic rib, which may compress the brachial plexus or subclavian vessel causing
thoracic outlet syndrome.

\medterm{Supernumerary Teeth} Teeth that develop in addition to the usual number. They may remain hidden or appear in the mouth and can crowd, displace, or prevent nearby teeth from erupting; a common type develops between the upper front teeth.

\medterm{Supervised Exercise Therapy} A structured, monitored exercise program used to improve walking ability and
reduce claudication symptoms in people with peripheral artery disease. The exercise plan
is adjusted to the person’s symptoms and functional capacity.

\medterm{Supination} Rotation of the forearm so that the palm faces anteriorly in anatomical position, or upward when the elbow is flexed. The radius and ulna become parallel.
\medterm{Supine} Lying flat on the back with the face and front of the body facing upward. This position is commonly used for examinations, imaging, and surgery. Staying in this position for a long time can increase the risk of pressure injuries; late in pregnancy, it may also reduce blood return to the heart.
\medterm{Supine Hypotensive Syndrome} A drop in blood pressure that can occur when a pregnant person lies flat on the back, because the enlarged uterus reduces blood return to the heart. Dizziness, nausea, sweating, or faintness may improve by turning onto the side.

\medterm{Supine Position} Lying on the back with the face and front of the body upward. It is commonly used for examination, imaging, and surgery, but prolonged use can cause pressure injury and late pregnancy may reduce blood return to the heart.

\medterm{Supplemental Oxygen Therapy} Administration of oxygen at a concentration above that of room air to treat or prevent
clinically important hypoxemia. Delivery method and oxygen target depend on the illness,
baseline oxygenation, risk of carbon-dioxide retention, and monitoring results; excessive
oxygen can also cause harm.

\medterm{Supplementation} Use of a vitamin, mineral, protein, fiber, or other nutrient preparation to correct a
deficiency or meet a clinically identified increased requirement. The indication, dose,
interactions, product quality, and risk of toxicity should be assessed before use.


\medterm{Supportive Care} Treatment that maintains essential body functions, relieves symptoms, prevents complications, and supports recovery while the underlying cause is treated.

\medterm{Suppository} A solid preparation inserted into the rectum, vagina, or urethra where it dissolves or melts.
\medterm{Suppression} Reduction or inhibition of a function, response, symptom, or process.
\medterm{Suppuration} The formation or discharge of pus, a thick fluid made from immune cells, germs, and damaged tissue. It usually occurs with infection and may form an abscess that needs drainage and treatment.
\medterm{Supracondylar Humerus Fracture} A break in the upper part of the elbow end of the arm bone, common in children after a fall on an outstretched hand. Swelling, pain, deformity, or inability to move the elbow may occur. Doctors must check circulation and nerves because severe fractures can damage an artery or cause dangerous swelling; some need urgent surgery.

\synonyms
Gartland Supracondylar Fracture; Pediatric Elbow Fracture

\medterm{Supraglottic Airway Device} A breathing device placed above the vocal cords to help deliver oxygen and ventilation during anesthesia or emergencies. It can be inserted more quickly than a breathing tube but does not protect against aspiration as reliably.

\medterm{Supranuclear Ophthalmoplegia} Weakness or loss of control of eye movements caused by damage to brain pathways above the eye-movement nerves. It may make it difficult to look in certain directions.

\medterm{Supranuclear Palsy} Alternate index label; see \textit{Progressive supranuclear palsy}.

\medterm{Suprapubic} Located above the pubic bone in the lower abdomen. The term may describe pain or examination in this area, or a procedure such as inserting a catheter through the lower abdominal wall directly into the bladder when a tube cannot be passed through the urethra.
\medterm{Suprapubic Catheter Care} Care of a drainage tube that enters the bladder through the lower abdomen, including cleaning the site, keeping the tube unobstructed, securing it, and monitoring urine. Fever, increasing pain, leakage, bleeding, or no urine flow needs medical advice.

\medterm{Suprarenal Glands (Adrenal Glands)} Two small glands above the kidneys. Their outer layer makes hormones that regulate salt, blood pressure, stress, and metabolism; their inner layer makes adrenaline and related hormones.
\medterm{Suprascapular Nerve} A nerve from the brachial plexus supplying shoulder muscles that help lift the arm and rotate it outward.

\medterm{Supraventricular Tachycardia} A fast heart rhythm beginning above the ventricles, often starting and stopping suddenly, causing palpitations, dizziness, or weakness.

\synonyms
SVT

\medterm{Supraventricular Tachycardia Paroxysmal (Paroxysmal Supraventricular Tachycardia)} Alternate index label for Paroxysmal Supraventricular Tachycardia.

\medterm{Suramin} An antiparasitic medicine used mainly for selected early forms of African trypanosomiasis. It requires specialist use because it can cause serious kidney, nerve, allergic, and other adverse effects.
\medterm{Surfactant} A substance that lowers surface tension. In the lungs, it coats the tiny air sacs and helps keep them open when a person breathes out; babies born very early may have too little and need treatment.
\medterm{Surgeon} A physician trained to perform operations and invasive procedures, and to evaluate and care for patients before and after surgery.
\medterm{Surgery} Treatment in which a clinician operates on tissues or organs to diagnose, remove, repair, replace, or improve a body structure. Risks and recovery depend on the procedure and the patient's health.
\medterm{Surgery For Pancreatic Cancer} An operation to remove a pancreatic tumor and sometimes nearby organs or lymph nodes when the cancer can be safely removed. It is major surgery and may cause bleeding, infection, digestive problems, diabetes, or delayed stomach emptying.


\medterm{Surgery For Pilonidal Cyst} An operation to drain or remove a recurrent or complicated tunnel or cyst near the cleft of the buttocks. Wound care and keeping the area clean are important; infection, slow healing, and recurrence can occur.

\medterm{Surgical Drain} A tube or device placed during or after surgery to remove blood, pus, bile, urine, or other fluid. The amount and appearance of drainage are monitored; infection, bleeding, blockage, and tissue injury are possible.

\medterm{Surgical Drill} A powered instrument used to cut, shape, or make holes in bone during surgery. It requires careful control to avoid heat injury, damage to nearby tissues, infection, or unintended bone weakening.

\medterm{Surgical Excision} Cutting out a lesion, tumor, organ, or other unwanted tissue using an operation. The removed material may be examined under a microscope; bleeding, infection, scarring, damage to nearby structures, and incomplete removal are possible.

\medterm{Surgical Fixation} Stabilization of a bone, joint, or other structure with devices such as plates, screws, rods, wires, or external frames.

\medterm{Surgical Hemostasis} Control of bleeding during an operation using pressure, ligation, sutures, clips,
electrosurgery, topical agents, vessel sealing, or other techniques. Effective
hemostasis requires identification of the bleeding source and correction of coagulopathy
when present.

\medterm{Surgical Instruments} Tools used during operations, including scalpels, scissors, forceps, clamps, retractors, needle holders, and devices controlling bleeding. Each instrument is designed for a specific task and must be cleaned, sterilized, and handled safely.

\medterm{Surgical Margin} The outer edge of tissue removed during an operation. A laboratory examination checks whether disease reaches that edge: a clear margin suggests no disease was seen there, while an involved margin may mean more treatment is needed.

\medterm{Surgical Repair} An operation that restores or rebuilds damaged tissue, an organ, or a body function. It may close wounds, reconnect structures, stabilize bones, repair blood vessels, or correct defects present from birth.

\medterm{Surgical Risk Stratification} Estimation of the likelihood of perioperative complications using patient factors,
urgency, procedure risk, functional capacity, and disease-specific assessment. It
informs optimization, monitoring level, shared decisions, and postoperative planning but
cannot eliminate uncertainty.

\medterm{Surgical Safety Checklist} A structured checklist used at defined points before anesthesia, before incision, and
before leaving the operating room to verify patient identity, procedure, site, consent,
equipment, specimens, and team communication.

\medterm{Surgical Site Infection} Infection involving the incision, organ, or operative space within the surveillance
period after surgery. It may be superficial incisional, deep incisional, or organ-space
and requires appropriate culture, antimicrobial therapy, drainage, debridement, or
reoperation when indicated.

\medterm{Surgical Stapler} A device that places metal or absorbable staples to close skin, join tissue, or seal a hollow organ during an operation. Incorrect placement can cause bleeding, leakage, narrowing, or injury and must be checked before completion.

\medterm{Surgical Time-Out} A deliberate team pause immediately before incision to verify patient identity,
procedure, site, consent, allergies, antibiotics, imaging, implants, equipment, and
anticipated critical events. It is a core surgical-safety intervention.

\medterm{Surgical Wound Care - Closed} Care of an incision whose edges are closed with sutures, staples, adhesive, or strips, including keeping it clean and dry as instructed and watching for infection or separation.

\medterm{Surgical Wound Care - Open} Care of a surgical wound left partly or completely open to drain or heal gradually, usually with dressings and scheduled assessment for infection, bleeding, and healthy tissue growth.

\medterm{Surgical Wound Classification} Categorization of an operative wound as clean, clean-contaminated, contaminated, or
dirty/infected according to the degree of microbial exposure and existing infection. The
class helps estimate surgical-site infection risk and guide prophylaxis.

\medterm{Surgical Wound Infection - Treatment} Management of infection in or around a surgical wound, including assessment of depth and severity, drainage or debridement when needed, wound care, and targeted antimicrobial therapy when indicated.

\medterm{Surrogacy} Surrogacy is an arrangement in which a person carries a pregnancy for intended parent or parents; medical, legal, ethical, and psychosocial requirements vary by jurisdiction.

\medterm{Surrogate} A substitute or proxy. In medical research, a surrogate outcome is a measurable change, such as a blood-test result, used to estimate a patient-important outcome such as fewer symptoms or longer survival.
\medterm{Surrogate Decision-Maker} A person authorized to make healthcare decisions for someone who lacks decision-making capacity, guided by the person's
known values and preferences.

\medterm{Surveillance} The planned collection and review of health information to detect changes in disease, follow a known condition, or identify problems early. It may involve routine reports, active case finding, selected monitoring sites, or regular follow-up of an individual.

\medterm{Surveillance Imaging} Planned repeat imaging used to monitor disease, treatment response, or possible recurrence. The timing and type of scan depend on the condition, previous treatment, symptoms, risks, and expected usefulness of results.

\medterm{Survival Rate} The percentage of people alive after a specified time from diagnosis or treatment, commonly measured at five years.

\medterm{Survivorship} The period of living during and after cancer treatment, together with the care provided during that time. It may include checking for return of cancer, treating late effects, supporting emotional health, improving quality of life, and managing other illnesses.

\medterm{Survivorship Care Plan} A written summary of cancer treatment, recommended follow-up, possible late effects, screening, medicines, and contacts for ongoing care. It helps the person and healthcare team coordinate care after treatment.

\medterm{Susac Syndrome} A rare disease in which inflammation and blockage of tiny blood vessels can affect the brain, eyes, and inner ears. It may cause confusion or headache, vision loss, and hearing loss; all features may not appear at once.

\medterm{Susceptibility} The likelihood of being affected by a disease, exposure, medicine, or microorganism. Susceptibility can reflect genetics, age, immunity, health conditions, previous exposure, dose, or the organism’s drug resistance.
\medterm{Susceptibility Testing} Laboratory testing that shows which antibiotics are likely to stop a person’s bacterial infection. The result helps clinicians choose treatment, especially when resistance is possible.

\medterm{Suspected Cause} A condition, exposure, or factor believed to explain a symptom or disease but not yet proven. Further history, examination, testing, or observation may support or reject the suspected cause.
\medterm{Suspension} A liquid medicine in which fine drug particles are spread through the liquid rather than dissolved. The container may need shaking so each dose contains the intended amount; settling, inaccurate measuring, and contamination can cause problems.

\medterm{Suture} A stitch or material used to close or join tissue.
\medterm{Suture (Perineal)} A thread or strand of material used to approximate tissues (bring wound edges together)
and hold them in place until healing occurs. See Suture.

\synonyms
Perineal

\medterm{Suture Anchor} A small implant that secures a tendon, ligament, or other soft tissue to bone. It may be
made from metal, polymer, or bioabsorbable material and can fail through pullout,
breakage, migration, or inflammatory reaction.

\medterm{Sutures - Ridged} Ridged cranial sutures are raised skull seams in an infant and may reflect normal molding or premature suture fusion, which requires assessment for craniosynostosis.

\medterm{Sutures - Separated} Separated cranial sutures in an infant can reflect normal growth or increased intracranial pressure, hydrocephalus, prematurity, or other illness and require clinical assessment.

\medterm{SUV} Standardized uptake value, a number estimating how much radioactive tracer is concentrated in tissue during a PET scan. Clinicians compare it with surrounding tissue and clinical findings to assess tumors or inflammation.

\medterm{SVC Obstruction} Blockage or compression of the superior vena cava, the large vein returning blood from the head, neck, arms, and upper chest to the heart. It may cause facial or arm swelling, prominent neck veins, cough, headache, or breathing difficulty.

\medterm{SVC Syndrome} Blockage or compression of the superior vena cava, the large vein returning blood from the upper body to the heart. It may cause swelling of the face, neck, or arms, visible neck veins, cough, or breathing difficulty.

\medterm{SVT} Alternate index label; see \textit{Supraventricular tachycardia}.

\medterm{Swab} A small absorbent stick used to collect a sample from the skin, throat, nose, wound, or another body site. The sample can be tested for infection or other conditions, and correct collection affects the result.
\medterm{Swallow Study} An X-ray test that records a person swallowing food or liquid mixed with barium. It shows how the mouth and throat work, whether material enters the airway, and whether changes could make swallowing safer.

\medterm{Swallowing} The movement of food, liquid, or saliva from the mouth through the throat and esophagus into the stomach. Problems may cause coughing, choking, food sticking, or material entering the airway.

\synonyms
Difficulty Swallowing (Swallowing)

\medterm{Swallowing Assessment} An evaluation of how safely and effectively a person moves food, liquid, or saliva from the mouth to the stomach. It may include observation, X-ray or camera testing, and recommendations about exercises, posture, or food and drink textures.

\medterm{Swallowing Chalk} Ingestion of chalk, which is usually low in toxicity in small amounts but can cause choking, constipation, or gastrointestinal upset; repeated eating may indicate pica.

\medterm{Swallowing Difficulty} Dysphagia, or impaired movement of food or liquid from the mouth to the stomach, caused by oropharyngeal, esophageal, neuromuscular, structural, or functional disorders.

\medterm{Swallowing Problems} Difficulty moving food, liquid, or saliva safely from the mouth to the stomach; dysphagia can cause choking, aspiration, weight loss, or dehydration and may need evaluation.

\medterm{Swallowing Soap} Ingestion of soap, which commonly causes mouth or stomach irritation, nausea, or diarrhea; concentrated or large exposures can injure tissue or affect breathing.

\medterm{Swallowing Sunscreen} Ingestion of sunscreen, usually causing mild gastrointestinal upset but potentially causing choking or aspiration depending on amount and product; breathing difficulty or persistent vomiting requires urgent help.

\medterm{Swallowing Syncope} Fainting triggered by swallowing, usually through an abnormal reflex that slows the heart or lowers blood pressure; recurrent episodes require evaluation for rhythm and structural heart disease.

\medterm{Swallowing Therapy} Assessment and treatment to make swallowing safer and more effective. A speech-language therapist may teach exercises, positioning, food-texture changes, and strategies to reduce aspiration into the airway.
\medterm{Swan-Ganz - Right Heart Catheterization} A procedure in which a thin catheter is passed through the right side of the heart into a lung artery to measure heart pressures and blood flow. It is used selectively in critically ill patients because bleeding, infection, rhythm problems, or vessel injury can occur.

\medterm{Swan-Ganz Catheter} A thin tube passed through the right side of the heart into a lung artery to measure pressures and estimate how much blood the heart pumps. It is used selectively in seriously ill patients because complications are possible.
\medterm{Sweat} Fluid produced by sweat glands, mainly to cool the body as it evaporates. Sweating can also increase with exercise, fever, stress, hormones, or some medicines.
\medterm{Sweat Chloride Test} A test that measures salt levels in sweat to help diagnose cystic fibrosis. A high result supports the diagnosis, but it is interpreted with symptoms and may be repeated or followed by genetic testing because other factors can affect results.

\medterm{Sweat Electrolytes Test} A sweat electrolytes test measures salts, especially chloride, in sweat and is commonly used to help diagnose cystic fibrosis.

\medterm{Sweat Gland} A skin gland that produces sweat, which contributes mainly to temperature regulation and skin surface function.

\medterm{Sweat Gland Tumor} A growth arising from cells of a sweat gland. Most are benign skin lumps, but some are cancerous and may grow, ulcerate, bleed, or spread. Examination and biopsy determine the type; complete removal is often used, with wider treatment and follow-up for cancerous tumors.

\medterm{Sweat Glands} Eccrine and apocrine glands that produce sweat or related secretions.
\medterm{Sweating} Production of sweat by eccrine and apocrine glands for heat regulation and other responses; excessive, absent, or asymmetric sweating can indicate medication, autonomic, endocrine, or neurologic disease.

\medterm{Sweating And Body Odor} Excess perspiration or noticeable odor produced when skin bacteria break down sweat; it may be influenced by hygiene, diet, medicines, hormones, or disease.

\synonyms
Body Odor And Sweating

\medterm{Sweating, Abnormally Excessive} Alternate index label; see \textit{Hyperhidrosis}.

\medterm{Sweating, Gustatory} Sweating triggered by eating or smelling food, usually over the cheek or temple after parotid-nerve injury or surgery and known as Frey syndrome.

\medterm{Sweet Syndrome} An inflammatory skin disorder characterized by sudden fever and tender red or violaceous plaques or nodules containing
many neutrophils. It may be associated with infection, medicines, inflammatory disease,
or cancer.

\medterm{Sweet's Syndrome} An inflammatory skin disorder causing sudden fever and painful red or purple raised patches or nodules. It may follow infection, result from a medicine, or occur with inflammatory disease or cancer, especially blood cancers.

\medterm{Sweetened Beverages} Sweetened beverages contain added sugars and can increase calorie intake, dental caries, obesity, and diabetes risk when consumed frequently.

\medterm{Sweeteners - Sugar Substitutes} Sugar substitutes provide sweetness with little or no sugar; types differ in energy, digestion, aftertaste, and effects on glucose and gastrointestinal symptoms.

\medterm{Sweeteners - Sugars} Sugars are simple carbohydrates that provide energy but can raise blood glucose and contribute to dental decay or excess calories when consumed in excess.

\medterm{Swelling} Abnormal enlargement caused by fluid accumulation, inflammation, bleeding, infection, tumor, or tissue injury; location, onset, pain, and associated symptoms identify the likely cause.

\medterm{Swimmer's Ear} Swimmer’s ear, also called otitis externa, is inflammation or infection of the skin lining the outer ear canal. It often develops when moisture remains in the canal, allowing bacteria or fungi to grow. Scratching, cotton swabs, or other irritation can also damage the canal’s protective skin. Symptoms include ear pain, itching, redness, swelling, drainage, and pain when the outer ear is touched.

\synonyms
Otitis Externa

\medterm{Swimmer's Itch} An itchy rash caused by an allergic reaction to microscopic parasites released by infected snails in freshwater or saltwater.

\synonyms
Cercarial Dermatitis
Dermatitis, Cercarial

\medterm{Swimmers Ear (Otitis Externa)} Alternate index label for Otitis Externa.

\medterm{Swimming Pool Cleaner Poisoning} Poisoning from pool chemicals such as chlorine compounds, acids, or alkalis. They may burn tissues or release toxic gases, causing cough, chest pain, breathing difficulty, vomiting, or shock; leave fumes and seek emergency help.

\medterm{Swimming Pool Granuloma} A slowly developing skin infection caused by a water-dwelling bacterium after exposure to an aquarium, swimming pool, or natural water. It can cause persistent bumps or sores, often spreading along the arm or hand, and needs specific antibiotic treatment.

\medterm{Swine Flu} Alternate informal term; see \textit{Influenza A Virus}.


\medterm{Swollen Ankles And Swollen Feet} Enlargement of the ankles or feet caused by injury, venous insufficiency, lymphatic disease, medication,
pregnancy, or systemic fluid retention from heart, kidney, or liver disease. Sudden one-sided swelling with pain,
redness, warmth, or breathlessness may indicate thrombosis or infection and requires urgent assessment.

\synonyms
Ankle Swollen (Swollen Ankles And Swollen Feet)

\medterm{Swollen Lymph Nodes} Swollen lymph nodes are enlarged immune glands, commonly from infection but sometimes from inflammation, medication, or cancer; persistent, hard, or generalized nodes need evaluation.

\synonyms
Lymph Nodes, Swollen
Lymph Swollen Glands (Swollen Lymph Nodes)
Lymph Swollen Nodes (Swollen Lymph Nodes)

\medterm{Swollen Tongue} Enlargement of the tongue caused by injury, infection, allergy, inflammation, bleeding, or another condition; sudden breathing difficulty is an emergency.

\medterm{Sycosis Barbae} Long-lasting inflammation or infection of hair follicles in the beard area, usually involving bacteria and ingrown hairs. It can cause painful bumps, pus, crusting, scarring, and hair loss and needs appropriate treatment.
\medterm{Sydenham Chorea} A delayed neurologic complication of group A streptococcal infection causing involuntary jerky movements, hypotonia, emotional lability, and coordination problems, usually in children.

\medterm{Sydenham's Chorea} A delayed immune complication of group A streptococcal infection, usually in children. It causes involuntary, irregular movements and may cause clumsiness, muscle weakness, emotional changes, or difficulty with speech and daily tasks.
\medterm{Symblepharon} Scar tissue that joins the inside of the eyelid to the eye’s surface, usually after a severe burn, inflammation, or eye disease. It can restrict eye movement, cause dryness or pain, and may need specialist treatment.


\medterm{Symbolism} The use of an image, object, or action to represent an idea or feeling; in clinical settings, symbolic meaning may be relevant to communication, culture, or psychotherapy.

\medterm{Sympathectomy} Surgical interruption or destruction of sympathetic nerves to reduce abnormal vasoconstriction, sweating, or pain in selected conditions.

\medterm{Sympathetic} Relating to the part of the autonomic nervous system that prepares the body for activity or stress by changing heart rate, blood-vessel tone, sweating, breathing, and digestion. In everyday language, it can also mean showing concern for another person.
\medterm{Sympathetic Ganglion} A small group of nerve cells outside the brain and spinal cord that relays automatic signals to blood vessels, sweat glands, and internal organs. These signals help control heart rate, blood pressure, digestion, and body responses to stress.

\medterm{Sympathetic Nervous System} The autonomic division that prepares the body for activity or stress by increasing heart rate and altering vascular, airway, and digestive function.
\medterm{Sympatholytic} A medicine or treatment that reduces adrenaline-like nerve activity. It may lower blood pressure, slow the heart, or reduce excessive sweating, but the effects and risks depend on the specific medicine and require appropriate prescribing.

\medterm{Sympathomimetic Drugs} Medicines that imitate the body's adrenaline-like nerve activity, often by stimulating alpha or beta receptors. They may open airways, raise blood pressure, or reduce congestion but can cause fast heartbeat, anxiety, or dangerous effects.
\medterm{Sympathomimetic Toxidrome} A poisoning pattern caused by excessive adrenergic stimulation, producing agitation, sweating, dilated pupils, rapid heart rate, high blood pressure, elevated temperature, tremor, or seizures.

\medterm{Symphysis} A joint where two bones are joined by a pad of tough cartilage. The pubic symphysis connects the front pelvic bones and normally allows slight movement, especially during pregnancy and childbirth.

\medterm{Symphysis-Fundal Height} The distance from the top of the pubic bone to the top of the uterus, measured through the abdomen during pregnancy. From about 24 to 36 weeks, the measurement often roughly matches the number of weeks of pregnancy. A larger or smaller result may lead to an ultrasound to check fetal growth, fluid, or the pregnancy dates.

\synonyms
Fundal Height Measurement; SFH

\medterm{Symphysis Pubis} The fibrocartilaginous joint connecting the left and right pubic bones at the front of the pelvis.

\medterm{Symptom} A change or problem felt or reported by a patient that may suggest illness, such as pain, nausea, or breathlessness.
\medterm{Symptomatic} Having symptoms caused by a disease or condition; symptomatic treatment relieves those symptoms without necessarily treating the underlying cause.

\medterm{Symptomatic Treatment} Care aimed at relieving symptoms such as pain, fever, nausea, itching, or cough without directly removing the underlying cause. It may be used alone or alongside treatment for the disease itself.

\medterm{Synapse} A connection where one nerve cell communicates with another nerve cell, muscle, or gland. Signals cross the small gap using chemicals or electrical changes to control movement, sensation, thought, and body functions.
\medterm{Synapse Assembly} The developmental process by which neurons form organized connections with other neurons or target cells, using cell-adhesion molecules and signaling proteins.

\medterm{Synapsins} A family of neuronal proteins that help organize synaptic vesicles and regulate the release of neurotransmitters.

\medterm{Synaptic Cleft} The very small gap between a presynaptic neuron and its target cell across which neurotransmitters travel.

\medterm{Synaptic Membrane} The specialized membrane at a neuronal junction, including the neurotransmitter-releasing presynaptic membrane and receptor-containing postsynaptic membrane.

\medterm{Synaptic Plasticity} The ability of connections between nerve cells to become stronger, weaker, or structurally different in response to activity. It supports learning, memory, adaptation, and some changes in chronic pain.

\medterm{Synaptic Transmission} Communication between nerve cells or between a nerve cell and another target across a synapse. Electrical signals trigger release of chemical messengers, which bind receptors and change the receiving cell’s activity.
\medterm{Synaptic Vesicle} A small membrane-bound sac in a nerve ending that stores neurotransmitters and releases them into the synaptic cleft when the neuron is activated.

\medterm{Synaptonemal Complex} A temporary protein structure that pairs matching chromosomes during egg or sperm formation. It helps chromosomes exchange genetic material accurately before cells divide and is important for normal reproductive development.


\medterm{Synaptosome} A sealed nerve-terminal structure produced when brain tissue is mechanically disrupted; it is used to study neurotransmitter release and synaptic function.

\medterm{Synaptotagmin} A calcium-sensing protein on nerve-cell storage sacs that helps trigger rapid release of chemical nerve signals. It is essential for communication between nerve cells and between nerves and muscles.

\medterm{Synarthrosis} A joint with little or no movement, such as a skull suture or tooth anchored in its socket, in which bones are joined by fibrous tissue or cartilage.



\medterm{Syncope} A sudden, brief loss of consciousness and postural tone caused by temporary global cerebral hypoperfusion, followed by spontaneous, usually complete recovery. Common mechanisms are reflex syncope, orthostatic hypotension, and cardiac disease; prolonged confusion or seizure-like features may indicate another cause.
\medterm{Syncope (Fainting)} A brief loss of consciousness and muscle control caused by a temporary reduction in blood flow to the brain, followed by spontaneous recovery. Common causes include reflex fainting, a sudden blood-pressure drop, dehydration, and heart disease; prolonged confusion may suggest another cause.

\synonyms
Fainting

\medterm{Syncope, Coughing} A brief faint caused by coughing, which raises chest pressure, reduces venous return and cardiac output, and transiently lowers cerebral blood flow.

\medterm{Syndactyly} Congenital joining or webbing of two or more fingers or toes. It may occur alone or with a genetic syndrome and sometimes requires reconstructive surgery.
\medterm{Syndecan} A group of proteins on the cell surface that help cells attach to their surroundings and receive signals. They contribute to tissue development, healing, and responses to injury.


\medterm{Syndesmophyte} A thin, bony growth that forms along the edges of the vertebrae, often in ankylosing spondylitis, and may bridge adjacent vertebrae.

\medterm{Syndrome} A group of symptoms, signs, or test findings that commonly occur together. A syndrome may have several possible causes and does not always name one specific disease.
\medterm{Syndrome Asperger (Asperger Syndrome)} Alternate index label for Asperger Syndrome.

\medterm{Syndrome Cauda Equina (Cauda Equina Syndrome)} Alternate index label for Cauda Equina Syndrome.

\medterm{Syndrome Compartment (Compartment Syndrome)} Alternate index label for Compartment Syndrome.

\medterm{Syndrome Cyclic Vomiting (Cyclic Vomiting Syndrome)} Alternate index label for Cyclic Vomiting Syndrome.

\medterm{Syndrome Loeys-Dietz (Loeys-Dietz Syndrome)} Alternate index label for Loeys-Dietz Syndrome.

\medterm{Syndrome Of Inappropriate ADH Secretion} Excessive antidiuretic hormone activity causing water retention and low blood sodium despite no major swelling or dehydration. It may result from lung disease, medicines, cancer, or brain disorders and can cause confusion or seizures.
\medterm{Syndrome Of Inappropriate Antidiuretic Hormone} Excessive release or effect of antidiuretic hormone causing impaired free-water excretion, low serum sodium, and inappropriately concentrated urine in the absence of an appropriate physiologic stimulus.

\medterm{Syndrome Of Inappropriate Antidiuretic Hormone Secretion} Excessive activity of the water-retaining hormone vasopressin causes the body to retain water and lowers the blood sodium level. Causes include medicines, lung disease, cancer, and brain disorders; severe cases can cause confusion or seizures.

\synonyms
Syndrome of Inappropriate Antidiuretic Hormone Secretion

\medterm{Syndrome Pfeiffer (Pfeiffer Syndrome)} Alternate index label for Pfeiffer Syndrome.

\medterm{Syndrome Restless Legs (Restless Leg Syndrome)} Alternate index label for Restless Leg Syndrome.


\medterm{Syndrome X} Alternate index label; see \textit{Metabolic syndrome}.

\medterm{Syndrome, Nervous Colon} An older name for irritable bowel syndrome, a disorder causing recurrent abdominal pain with diarrhea, constipation, or both without a structural explanation. Symptoms may improve with individualized diet, medicines, and stress management.

\medterm{Syndrome, Post-Polio} New or worsening weakness, fatigue, muscle pain, breathing problems, or swallowing difficulty occurring years after recovery from polio. It reflects late effects on previously affected motor neurons and is managed with pacing, rehabilitation, and support.

\medterm{Syndrome, TAR} Thrombocytopenia-absent-radius syndrome, an inherited condition causing low platelet counts and absent or severely shortened radius bones in the forearms. It may also cause limb differences, heart defects, kidney problems, or milk intolerance.

\medterm{Synechia} An abnormal band of scar-like tissue joining surfaces that should remain separate. In the eye, it may attach the iris to the lens or cornea and interfere with fluid drainage or vision.

\medterm{Synechiae} Abnormal bands of scar-like tissue joining surfaces that should remain separate. They may occur in the eye, uterus, nose, or other sites and can restrict movement, block passages, or impair function.
\medterm{Synergist} A muscle, drug, or factor that works together with another to produce or enhance an effect.
\medterm{Synesthesia} Synesthesia is a stable perceptual phenomenon in which stimulation of one sensory pathway automatically produces an experience in another, such as seeing colors with sounds.


\medterm{Synostosis} Fusion of two bones by bone tissue, present from birth or developing later. It may be harmless, but fusion in the skull, forearm, or other important area can restrict growth or movement.
\medterm{Synovectomy} An operation to remove part or all of an inflamed lining inside a joint. It may reduce pain and swelling in selected cases of arthritis or repeated joint inflammation and can be done through small cuts or open surgery.
\medterm{Synovial Biopsy} A synovial biopsy removes a small sample of the joint lining to investigate infection, inflammatory arthritis, crystal disease, or an unusual joint lesion.

\medterm{Synovial Bursa} A small, fluid-filled sac near a joint that reduces rubbing between moving tissues such as tendons, muscles, skin, and bone. Irritation or inflammation of a bursa causes bursitis, pain, and swelling.

\medterm{Synovial Chondromatosis} A usually noncancerous condition in which the joint lining forms many small cartilage nodules that may break loose. It can cause pain, swelling, catching, and limited movement, often in the knee, hip, or elbow.

\medterm{Synovial Cyst} A synovial cyst is a fluid-filled sac arising from or near a joint or tendon sheath, sometimes compressing nearby nerves and causing pain or weakness.

\medterm{Synovial Fluid} A slippery fluid inside a freely movable joint that reduces friction and helps nourish cartilage. Changes in its amount or contents can occur with injury, infection, bleeding, crystals, or arthritis.
\medterm{Synovial Fluid Analysis} Testing fluid removed from a swollen or painful joint. The laboratory examines its appearance, cells, crystals, and germs to help distinguish injury, arthritis, gout, bleeding, and joint infection; a very hot painful joint needs urgent assessment.

\medterm{Synovial Joint} A freely movable joint in which the articulating bones are separated by a fluid-filled joint cavity enclosed by a capsule.

\medterm{Synovial Membrane} The thin lining inside a freely moving joint. It makes synovial fluid, which reduces rubbing and helps nourish the smooth joint surfaces; inflammation of this lining causes swelling and pain.
\medterm{Synovial Reaction To Implanted Prostheses} Inflammation or tissue change around an artificial joint caused by wear particles, loosening, corrosion, infection, or an immune reaction. Pain, swelling, or reduced function after replacement needs evaluation to identify the cause.

\medterm{Synovial Sarcoma} A cancer of soft tissue that often develops near a joint in an arm or leg, especially in adolescents and young adults. Diagnosis requires tissue testing; treatment may include surgery, radiation, or chemotherapy.

\medterm{Synoviocytes} Synoviocytes are cells lining the joint synovium that produce lubricating fluid and inflammatory mediators and help maintain or damage joint tissue.

\medterm{Synovitis} Inflammation of the lining inside a joint, causing swelling, warmth, pain, or reduced movement. It may follow injury or occur with arthritis, gout, autoimmune disease, or infection; a hot, very painful joint needs urgent assessment.
\medterm{Synovium} Tissue lining a joint or tendon sheath that makes lubricating fluid and helps support smooth movement.
\medterm{Synovography} An imaging test in which contrast is placed inside a joint to outline its capsule, lining, or possible tear. It is used selectively; MRI, ultrasound, or other scans are often preferred for joint problems.


\medterm{Synthetic Androgen} A synthetic androgen is a manufactured compound that activates androgen receptors and may mimic testosterone; therapeutic and anabolic use can suppress gonadal function and cause acne, mood effects, infertility, liver injury, or cardiovascular harm.

\medterm{Syphilis} A sexually transmitted infection caused by Treponema pallidum that progresses through primary, secondary, latent, and tertiary stages and can affect the nervous system, eyes, heart, or fetus.
\medterm{Syphilis Serology (RPR/VDRL)} Blood tests used to screen for syphilis and monitor response to treatment. A reactive result usually needs confirmation with a different test and interpretation alongside symptoms, examination, and treatment history.

\synonyms
Syphilis Serology
\medterm{Syphilitic Alopecia} Non-scarring hair loss occurring as a manifestation of secondary syphilis, presenting
with diffuse thinning or patchy, moth-eaten alopecia of the scalp, eyebrows, or beard.
It is thought to result from direct spirochaetal invasion of hair follicles during the
spirochaetaemic phase.

\medterm{Syphilitic Aseptic Meningitis} Syphilitic aseptic meningitis is inflammation of the meninges caused by Treponema pallidum, producing headache, neck stiffness, cranial-nerve symptoms, or other neurologic findings.

\medterm{Syphilitic Gumma} A destructive lump or area of inflammation caused by long-standing untreated syphilis. It can affect the skin, bones, liver, brain, or other organs and may form a firm mass or ulcer; appropriate antibiotic treatment can prevent further damage.

\medterm{Syringe} A medical device with a hollow barrel, movable plunger, and tip that may connect to a needle, catheter, or other fitting. It is used to draw up, measure, inject, or remove fluids. Common types include disposable, oral, insulin, tuberculin, and safety syringes.

\synonyms
Medical Syringe
\medterm{Syringe Drivers} Small pumps that push medicine from a syringe continuously or at programmed times. They provide controlled delivery when precise dosing is needed, including pain relief, symptom control, or medicines given under the skin.
\medterm{Syringocystadenoma Papilliferum} A usually harmless sweat-gland skin growth, often appearing as a small bump or plaque on the scalp or face. It may develop within a birthmark and can resemble other lesions, so diagnosis may require examination or biopsy.

\medterm{Syringoma} A syringoma is a benign tumor of sweat-duct cells, usually appearing as small skin-colored or yellowish papules around the eyelids or elsewhere.

\medterm{Syringomyelia} A fluid-filled cavity that develops inside the spinal cord. It can slowly damage the cord and cause pain, weakness, loss of pain or temperature sensation, stiffness, or problems with bladder, bowel, or sexual function.
\medterm{Syrinx} A fluid-filled cavity within the spinal cord, usually related to disturbed movement of the fluid around the brain and spinal cord. It can damage nerve tissue and cause pain, weakness, stiffness, or loss of sensation; a syrinx in the cord is called syringomyelia.

\medterm{Syrup} A syrup is a viscous liquid preparation, often sweetened, used as a medicine vehicle or oral dosage form; concentration and ingredients determine safe dosing.


\medterm{Systemic} Affecting the whole body or relating to an entire system rather than a localized site.
\medterm{Systemic Antifungals} Antifungal medicines taken orally or given intravenously for infections that are extensive,
invasive, refractory, or unsuitable for topical treatment. Choice and duration depend on the
organism, site, severity, drug interactions, organ function, and local resistance patterns.

\medterm{Systemic Chemotherapy} Systemic chemotherapy uses medicines that travel through the bloodstream to kill or slow cancer cells throughout the body.

\medterm{Systemic Circuit} The circulation that carries oxygenated blood from the left ventricle through the body and returns deoxygenated blood to the right atrium.

\medterm{Systemic Disease} A disease that affects the whole body or several organs rather than one local area. Examples include autoimmune disease, widespread infection, hormonal disorders, and some cancers.
\medterm{Systemic Effects} Effects involving the whole body rather than one local area. A medicine, infection, toxin, or disease may produce systemic cardiac, neurologic, kidney, liver, blood, or other effects.

\medterm{Systemic Exertional Intolerance Disease} Alternate index label; see \textit{Myalgic encephalomyelitis/chronic fatigue syndrome (ME/CFS)}.

\medterm{Systemic Inflammatory Response Syndrome} A widespread inflammatory reaction that may follow infection, injury, burns, pancreatitis, or surgery. It can cause abnormal temperature, heart rate, breathing rate, or white blood-cell count and may progress to organ failure.

\medterm{Systemic Juvenile Idiopathic Arthritis} A childhood inflammatory disease causing repeated high fevers, a temporary pink rash, and swollen or painful joints. It can also enlarge lymph nodes, liver, or spleen and may affect other organs, requiring specialist monitoring and treatment.

\medterm{Systemic Lupus} Systemic lupus erythematosus is an autoimmune disease that can affect skin, joints, kidneys, blood, nervous system, lungs, and heart. Symptoms fluctuate, and treatment is tailored to organ involvement with sun protection, monitoring, and immunomodulatory medicines.

\synonyms
Arthritis SLE (Systemic Lupus)
Lupus (Systemic Lupus)
SLE (Systemic Lupus)

\medterm{Systemic Lupus Erythematosus} A long-term autoimmune disease in which the immune system mistakenly attacks the body’s own tissues and causes inflammation. It can affect the skin, joints, kidneys, heart, lungs, brain, and blood cells. Symptoms vary widely and may include fatigue, fever, joint pain or swelling, rashes, hair loss, mouth sores, or swelling. The disease often has periods of increased activity called flares and periods of fewer symptoms called remission; inflammation can cause organ damage in some people.
\medterm{Systemic Mastocytosis} A disorder in which too many mast cells build up in the skin, bone marrow, digestive tract, or other organs. The cells can release chemicals causing flushing, itching, faintness, stomach symptoms, bone problems, or severe allergic reactions.

\medterm{Systemic Scleroderma} An autoimmune disease that causes tight, thickened skin and may scar blood vessels and internal organs. The esophagus, lungs, kidneys, heart, fingers, and other areas can be affected, so care depends on the organs involved.

\medterm{Systemic Sclerosis} An autoimmune disease in which the body makes too much scar-like tissue. The skin may become tight and hard, and blood vessels, the esophagus, lungs, kidneys, or heart may also be affected.

\medterm{Systemic Therapy} Treatment that travels through the bloodstream or otherwise affects the whole body. Examples include chemotherapy, hormone treatment, targeted treatment, and immunotherapy; it differs from local treatment such as surgery or focused radiation.

\medterm{Systemic Vascular Resistance} The opposition to blood flow in the body’s arteries, controlled mainly by how narrow or wide the small arteries are. It helps determine blood pressure and the work required of the heart.

\medterm{Supracervical Hysterectomy} Hysterectomy in which the uterine body is removed while the cervix is left in place; it is also called subtotal or partial hysterectomy. The fallopian tubes and ovaries may be managed separately.

\medterm{Suprapubic Catheter} A tube placed through the lower abdominal wall directly into the bladder to drain urine. It may be used when urine cannot drain safely through the urethra or when longer-term bladder drainage is needed.

\medterm{Systemic Venous} Systemic venous refers to veins returning deoxygenated blood from the body to the right side of the heart, excluding the pulmonary circulation.

\medterm{Sustained Virologic Response} No detectable hepatitis C virus in the blood 12 or more weeks after antiviral treatment. This is considered a cure of that infection, but it does not protect against catching hepatitis C again.

\medterm{Systole} The part of each heartbeat when a heart chamber contracts and pushes blood forward. Ventricular systole sends blood to the lungs and body, while arterial pressure rises and the heart's valves coordinate flow.
\medterm{Systolic} Systolic means relating to contraction of the heart, especially the pressure in arteries when the ventricles eject blood.

\medterm{Systolic Anterior Motion} Movement of the mitral valve toward the path where blood leaves the heart during each heartbeat. It can partly block blood flow and cause the valve to leak, especially in hypertrophic cardiomyopathy. Symptoms may include breathlessness, chest pain, dizziness, or fainting; treatment depends on the cause and severity.

\synonyms
SAM; SAM of the Mitral Valve

\medterm{Systolic Murmurs} Extra heart sounds heard while the heart is contracting. They may result from fast blood flow, a narrowed heart valve, or another structural problem; their significance depends on the sound and the person’s symptoms.

\medterm{Systolic Pressure} The maximum arterial pressure during ventricular contraction, recorded as the upper number of a blood-pressure reading.


\medterm{Saffron} Saffron is a spice made from the dried stigmas of Crocus sativus. It is used in food and studied as a possible supplement, but it is not a substitute for proven medical treatment.

\medterm{Science Communication} Science communication is the clear sharing of scientific information with patients, professionals, policymakers, or the public.

\medterm{Sclerotomy} A sclerotomy is a surgical opening made through the sclera, the white outer layer of the eye, often during vitreoretinal surgery.

\medterm{Seasonal Affective Disorder (SAD)} Alternate index label for Seasonal Affective Disorder.

\medterm{Secondary Cancer} A secondary cancer is a tumor that has spread from its original site to another part of the body. It is also called a metastasis or metastatic cancer.

\medterm{Sedative} A sedative is a medicine or substance that reduces alertness and may cause relaxation or sleepiness. It can impair coordination and breathing, especially when combined with alcohol or opioids.

\medterm{Selective Serotonin Reuptake Inhibitor (SSRI)} A selective serotonin reuptake inhibitor is an antidepressant that increases serotonin signaling in the brain. Examples include fluoxetine, sertraline, and escitalopram.

\medterm{Self-Diagnosis} Self-diagnosis is identifying a health problem without a professional assessment. It can be inaccurate, so concerning or persistent symptoms should be evaluated by a clinician.

\medterm{Sentinel Node} A sentinel node is the first lymph node or group of nodes to which cancer cells are most likely to spread from a primary tumor. A sentinel lymph-node biopsy helps assess spread.

\medterm{Severe Combined Immunodeficiency (SCID)} A group of inherited disorders in which the immune system cannot make effective defenses against germs. Babies may develop repeated, severe infections, persistent diarrhea, or poor growth; early diagnosis and treatment can be lifesaving.

\medterm{Sexual Health} Sexual health includes physical, emotional, mental, and social well-being related to sexuality. It includes contraception, fertility, consent, relationships, and prevention of sexually transmitted infections.

\medterm{Sexually Transmitted Disease (STD)} A sexually transmitted disease is an infection spread mainly through sexual contact. The term sexually transmitted infection is also used because some infections cause no symptoms.

\medterm{Sheehan's Syndrome} Loss of pituitary hormones caused by damage to the pituitary gland after severe bleeding or dangerously low blood pressure during or after childbirth. It may cause inability to breastfeed, extreme tiredness, low blood pressure, menstrual changes, or other hormone deficiencies.

\medterm{Shilajit} Shilajit is a mineral-rich, resin-like substance used in some traditional medicines. Product purity and safety vary, and contamination with heavy metals is possible.

\medterm{Shisha} Shisha is tobacco or another smoking mixture used in a water pipe. The water does not remove the main toxic substances, and smoking shisha can expose users to nicotine, carbon monoxide, and cancer-causing chemicals.

\medterm{Short-Chain Fatty Acids} Small fatty substances made when bacteria in the large intestine break down dietary fiber. They help nourish the bowel lining and may influence inflammation and metabolism; their effects vary with diet and the person’s gut bacteria.

\medterm{Singing} Singing is producing musical sounds with the voice. It can support communication and social well-being, while excessive or improper voice use may cause vocal strain.


\medterm{Sirtuins} A family of enzymes that change other proteins and help regulate how cells use energy, respond to stress, and switch genes on or off. Their possible roles in aging and disease are still being studied.

\medterm{Sjogren's Syndrome} Older name for Sjögren disease, an autoimmune disease that commonly causes dry eyes and dry mouth.

\medterm{Skin Cells} Skin cells form the layers of the skin and include keratinocytes, melanocytes, immune cells, and sensory cells. They provide protection, pigmentation, sensation, and repair.

\medterm{Skin Cyst} A skin cyst is a closed sac under or within the skin, often filled with fluid, keratin, or other material. Many are harmless, but a painful, infected, or changing cyst should be assessed.

\medterm{Sleep Disorder} Alternate index label for Sleep Disorders.

\medterm{Small Incision Lenticule Extraction Surgery} Small incision lenticule extraction, or SMILE, is a laser eye procedure that removes a small piece of corneal tissue through a tiny incision to correct some cases of nearsightedness and astigmatism.

\medterm{Smoking and Pregnancy} Smoking during pregnancy exposes the fetus to nicotine and other harmful chemicals and increases risks such as poor growth, preterm birth, placental problems, and stillbirth. Stopping at any time is beneficial.

\medterm{Social Media} Social media consists of online platforms where people create, share, and discuss content. Its effects on health vary and may involve social support, sleep disruption, stress, or exposure to misinformation.

\medterm{Social Media Addiction} Social media addiction is a pattern of difficult-to-control use that causes distress or interferes with sleep, work, school, relationships, or daily activities. It is not a universally defined formal diagnosis.

\medterm{Solar Elastosis} Solar elastosis is degeneration and thickening of elastic fibers in skin after long-term ultraviolet exposure. It contributes to leathery, wrinkled, yellowish sun-damaged skin.

\medterm{Solitary Kidney} A solitary kidney means a person has one functioning kidney, because the other kidney is absent, removed, or not working. Many people remain healthy but may need blood-pressure and kidney-function monitoring.

\medterm{Sonohysterography} Sonohysterography is an ultrasound examination in which sterile fluid is placed in the uterus to outline the uterine cavity and assess problems such as polyps or fibroids.

\medterm{Soybeans} Soybeans are legumes that provide protein, fiber, and compounds called isoflavones. They are used in foods such as tofu and soy milk; allergy is possible.


\medterm{Spirulina} Spirulina is a type of cyanobacterium sold as a dietary supplement. Products may provide nutrients, but contamination and medication interactions are possible, and it is not a reliable substitute for a balanced diet.

\medterm{Splinter} A splinter is a small piece of wood, glass, metal, or another material embedded in the skin. It should be removed carefully; medical care may be needed if it is deep, infected, or near an eye or joint.

\medterm{Splinter Hemorrhage} A splinter hemorrhage is a thin line of bleeding under a fingernail or toenail. It may follow injury, but multiple or unexplained lesions can occur with medical conditions and deserve evaluation.

\medterm{Sports Medicine} Medical care focused on preventing, diagnosing, and treating exercise- and sports-related injuries and conditions. It also supports safe training, rehabilitation, performance, and return to activity.
\medterm{Squamous Cell Carcinoma in Situ} Squamous cell carcinoma in situ, also called Bowen disease, is an early skin cancer limited to the outer skin layer. It can become invasive if untreated.

\medterm{Stafne Bone Defect} A Stafne bone defect is a usually harmless, developmental indentation in the lower jaw caused by nearby salivary-gland tissue. It is often found incidentally on imaging.

\medterm{Stem Cells} Stem cells are cells that can self-renew and develop into specialized cell types. They are important in normal tissue repair and in some established treatments, such as blood-forming stem-cell transplantation.

\medterm{Still's Disease} An inflammatory illness that can affect children or adults. It commonly causes repeated high fevers, a fleeting rash, joint pain or swelling, and marked inflammation; it can affect other organs and usually requires specialist treatment.

\medterm{Stomach Pain} Stomach pain is discomfort in the upper abdomen. Causes range from indigestion and infection to ulcers, gallbladder disease, or pancreatitis; severe or persistent pain needs medical assessment.

\medterm{Streptococcus pneumoniae} Streptococcus pneumoniae is a bacterium that can cause pneumonia, meningitis, ear infection, sinusitis, and bloodstream infection. Vaccination helps prevent serious disease.

\medterm{Stress Related} Describes symptoms or health problems that are triggered or worsened by physical or emotional stress. Stress can affect sleep, pain, digestion, mood, heart rate, and blood pressure, but symptoms still require assessment for other causes.

\medterm{Stretch Mark} A stretch mark is a linear scar caused by stretching and changes in the deeper skin, often during pregnancy, growth, or rapid weight change. It is harmless and usually fades over time.

\medterm{Stunted Growth} Stunted growth is height substantially below expected levels for age, often due to long-term undernutrition, chronic illness, infection, or other factors. Children with poor growth should be evaluated.

\medterm{Stutter} Stuttering is a speech-fluency disorder involving repetitions, prolongations, or blocks in speech. Speech-language therapy can improve communication and confidence.

\medterm{Sudden Infant Death Syndrome (SIDS)} Sudden infant death syndrome is the sudden, unexplained death of an infant younger than one year, usually during sleep. Safe-sleep practices reduce risk: place babies on their backs on a firm, flat surface without loose bedding.

\medterm{Sudden Unexpected Death in Epilepsy} An unexpected death in a person with epilepsy that is not caused by drowning, injury, prolonged seizure, or another known condition. It is uncommon; controlling seizures, taking prescribed treatment, and discussing nighttime safety may reduce risk.

\medterm{Superfood} Superfood is a marketing term for a food promoted as especially nutritious. No single food provides all needed nutrients, and health claims should be judged against scientific evidence.

\medterm{Supplements} Dietary supplements contain substances such as vitamins, minerals, herbs, or amino acids. They can have side effects and interactions, so patients should discuss their use with a healthcare professional.

\medterm{Synaesthesia} Synaesthesia is a neurological trait in which stimulation of one sense or mental process automatically produces an additional sensory experience, such as seeing colors when hearing sounds.

\medterm{Synthetic Biology} Synthetic biology combines biology and engineering to design or modify biological systems. It has potential medical uses but also requires careful safety and ethical oversight.
